% Options for packages loaded elsewhere
\PassOptionsToPackage{unicode}{hyperref}
\PassOptionsToPackage{hyphens}{url}
\PassOptionsToPackage{dvipsnames,svgnames,x11names}{xcolor}
%
\documentclass[
  12pt,
  a4paper,
  extrafontsizes,
  onecolumn,
  openright]{book}
\usepackage{amsmath,amssymb}
\usepackage{lmodern}
\usepackage{iftex}
\ifPDFTeX
  \usepackage[T1]{fontenc}
  \usepackage[utf8]{inputenc}
  \usepackage{textcomp} % provide euro and other symbols
\else % if luatex or xetex
  \usepackage{unicode-math}
  \defaultfontfeatures{Scale=MatchLowercase}
  \defaultfontfeatures[\rmfamily]{Ligatures=TeX,Scale=1}
  \setmainfont[Extension=.otf,UprightFont=*-regular,BoldFont=*-bold,BoldItalicFont=*-bolditalic,ItalicFont=*-italic]{texgyretermes}
  \setmathfont[]{texgyretermes-math.otf}
\fi
% Use upquote if available, for straight quotes in verbatim environments
\IfFileExists{upquote.sty}{\usepackage{upquote}}{}
\IfFileExists{microtype.sty}{% use microtype if available
  \usepackage[]{microtype}
  \UseMicrotypeSet[protrusion]{basicmath} % disable protrusion for tt fonts
}{}
\makeatletter
\@ifundefined{KOMAClassName}{% if non-KOMA class
  \IfFileExists{parskip.sty}{%
    \usepackage{parskip}
  }{% else
    \setlength{\parindent}{0pt}
    \setlength{\parskip}{6pt plus 2pt minus 1pt}}
}{% if KOMA class
  \KOMAoptions{parskip=half}}
\makeatother
\usepackage{xcolor}
\usepackage{color}
\usepackage{fancyvrb}
\newcommand{\VerbBar}{|}
\newcommand{\VERB}{\Verb[commandchars=\\\{\}]}
\DefineVerbatimEnvironment{Highlighting}{Verbatim}{commandchars=\\\{\}}
% Add ',fontsize=\small' for more characters per line
\usepackage{framed}
\definecolor{shadecolor}{RGB}{248,248,248}
\newenvironment{Shaded}{\begin{snugshade}}{\end{snugshade}}
\newcommand{\AlertTok}[1]{\textcolor[rgb]{0.94,0.16,0.16}{#1}}
\newcommand{\AnnotationTok}[1]{\textcolor[rgb]{0.56,0.35,0.01}{\textbf{\textit{#1}}}}
\newcommand{\AttributeTok}[1]{\textcolor[rgb]{0.77,0.63,0.00}{#1}}
\newcommand{\BaseNTok}[1]{\textcolor[rgb]{0.00,0.00,0.81}{#1}}
\newcommand{\BuiltInTok}[1]{#1}
\newcommand{\CharTok}[1]{\textcolor[rgb]{0.31,0.60,0.02}{#1}}
\newcommand{\CommentTok}[1]{\textcolor[rgb]{0.56,0.35,0.01}{\textit{#1}}}
\newcommand{\CommentVarTok}[1]{\textcolor[rgb]{0.56,0.35,0.01}{\textbf{\textit{#1}}}}
\newcommand{\ConstantTok}[1]{\textcolor[rgb]{0.00,0.00,0.00}{#1}}
\newcommand{\ControlFlowTok}[1]{\textcolor[rgb]{0.13,0.29,0.53}{\textbf{#1}}}
\newcommand{\DataTypeTok}[1]{\textcolor[rgb]{0.13,0.29,0.53}{#1}}
\newcommand{\DecValTok}[1]{\textcolor[rgb]{0.00,0.00,0.81}{#1}}
\newcommand{\DocumentationTok}[1]{\textcolor[rgb]{0.56,0.35,0.01}{\textbf{\textit{#1}}}}
\newcommand{\ErrorTok}[1]{\textcolor[rgb]{0.64,0.00,0.00}{\textbf{#1}}}
\newcommand{\ExtensionTok}[1]{#1}
\newcommand{\FloatTok}[1]{\textcolor[rgb]{0.00,0.00,0.81}{#1}}
\newcommand{\FunctionTok}[1]{\textcolor[rgb]{0.00,0.00,0.00}{#1}}
\newcommand{\ImportTok}[1]{#1}
\newcommand{\InformationTok}[1]{\textcolor[rgb]{0.56,0.35,0.01}{\textbf{\textit{#1}}}}
\newcommand{\KeywordTok}[1]{\textcolor[rgb]{0.13,0.29,0.53}{\textbf{#1}}}
\newcommand{\NormalTok}[1]{#1}
\newcommand{\OperatorTok}[1]{\textcolor[rgb]{0.81,0.36,0.00}{\textbf{#1}}}
\newcommand{\OtherTok}[1]{\textcolor[rgb]{0.56,0.35,0.01}{#1}}
\newcommand{\PreprocessorTok}[1]{\textcolor[rgb]{0.56,0.35,0.01}{\textit{#1}}}
\newcommand{\RegionMarkerTok}[1]{#1}
\newcommand{\SpecialCharTok}[1]{\textcolor[rgb]{0.00,0.00,0.00}{#1}}
\newcommand{\SpecialStringTok}[1]{\textcolor[rgb]{0.31,0.60,0.02}{#1}}
\newcommand{\StringTok}[1]{\textcolor[rgb]{0.31,0.60,0.02}{#1}}
\newcommand{\VariableTok}[1]{\textcolor[rgb]{0.00,0.00,0.00}{#1}}
\newcommand{\VerbatimStringTok}[1]{\textcolor[rgb]{0.31,0.60,0.02}{#1}}
\newcommand{\WarningTok}[1]{\textcolor[rgb]{0.56,0.35,0.01}{\textbf{\textit{#1}}}}
\usepackage{longtable,booktabs,array}
\usepackage{calc} % for calculating minipage widths
% Correct order of tables after \paragraph or \subparagraph
\usepackage{etoolbox}
\makeatletter
\patchcmd\longtable{\par}{\if@noskipsec\mbox{}\fi\par}{}{}
\makeatother
% Allow footnotes in longtable head/foot
\IfFileExists{footnotehyper.sty}{\usepackage{footnotehyper}}{\usepackage{footnote}}
\makesavenoteenv{longtable}
\usepackage{graphicx}
\makeatletter
\def\maxwidth{\ifdim\Gin@nat@width>\linewidth\linewidth\else\Gin@nat@width\fi}
\def\maxheight{\ifdim\Gin@nat@height>\textheight\textheight\else\Gin@nat@height\fi}
\makeatother
% Scale images if necessary, so that they will not overflow the page
% margins by default, and it is still possible to overwrite the defaults
% using explicit options in \includegraphics[width, height, ...]{}
\setkeys{Gin}{width=\maxwidth,height=\maxheight,keepaspectratio}
% Set default figure placement to htbp
\makeatletter
\def\fps@figure{htbp}
\makeatother
\setlength{\emergencystretch}{3em} % prevent overfull lines
\providecommand{\tightlist}{%
  \setlength{\itemsep}{0pt}\setlength{\parskip}{0pt}}
\setcounter{secnumdepth}{5}
\ifLuaTeX
\usepackage[bidi=basic]{babel}
\else
\usepackage[bidi=default]{babel}
\fi
\babelprovide[main,import]{french}
% get rid of language-specific shorthands (see #6817):
\let\LanguageShortHands\languageshorthands
\def\languageshorthands#1{}
\usepackage{booktabs}
\usepackage{graphicx}
\usepackage{caption}
\captionsetup[figure]{labelfont=bf, font=small, justification=centering} % Met en gras le label et centre la légende
\usepackage{float} % Pour contrôler la position des figures
\floatplacement{table}{H} % Forcer les tableaux à apparaître exactement là où ils sont insérés
\captionsetup[table]{labelfont=bf, font=small, justification=centering}

\usepackage{fvextra}  % Meilleure gestion de verbatim
\RecustomVerbatimEnvironment{Verbatim}{BVerbatim}{breaklines=true,fontsize=\small}
\DefineVerbatimEnvironment{Highlighting}{Verbatim}{breaklines,commandchars=\\\{\}}
\usepackage{mdframed}
\usepackage{listings}
\usepackage{xcolor}

\definecolor{lightgray}{gray}{0.95}

\surroundwithmdframed[
  linewidth=1pt,
  backgroundcolor=lightgray,
  innerleftmargin=5pt,
  innerrightmargin=5pt
]{lstlisting}

\lstset{
  basicstyle=\ttfamily\footnotesize,
  breaklines=true,
  frame=single
}
\usepackage{tcolorbox}

% Définition d'un callout block personnalisé
\newtcolorbox{callout}[1][]{
  colback=blue!5!white, colframe=blue!75!black, 
  fonttitle=\bfseries, title=#1
}
\usepackage[authoryear]{natbib}


\usepackage{booktabs}
\usepackage{longtable}
\usepackage{array}
\usepackage{multirow}
\usepackage{wrapfig}
\usepackage{float}
\usepackage{colortbl}
\usepackage{pdflscape}
\usepackage{tabu}
\usepackage{threeparttable}
\usepackage{threeparttablex}
\usepackage[normalem]{ulem}
\usepackage{makecell}
\usepackage{xcolor}
\ifLuaTeX
  \usepackage{selnolig}  % disable illegal ligatures
\fi
\usepackage[authoryear]{natbib}
\bibliographystyle{apalike}
\usepackage{csquotes}
\IfFileExists{bookmark.sty}{\usepackage{bookmark}}{\usepackage{hyperref}}
\IfFileExists{xurl.sty}{\usepackage{xurl}}{} % add URL line breaks if available
\urlstyle{same} % disable monospaced font for URLs
\hypersetup{
  pdftitle={DISTRIBUTION DE LA DIVERSITE LIGNEUSE ET DES STOCKS DE CARBONE DANS LES ZONES DE TRANSITION FORETS-SAVANE: CAS DES FORETS COMMUNALES DE NANGA-EBOKO ET NDIKINIMEKI},
  pdfauthor={Nkouankam Aurelien},
  pdflang={fr-FR},
  colorlinks=true,
  linkcolor={Maroon},
  filecolor={Maroon},
  citecolor={Blue},
  urlcolor={Blue},
  pdfcreator={LaTeX via pandoc}}

\title{DISTRIBUTION DE LA DIVERSITE LIGNEUSE ET DES STOCKS DE CARBONE DANS LES ZONES DE TRANSITION FORETS-SAVANE: CAS DES FORETS COMMUNALES DE NANGA-EBOKO ET NDIKINIMEKI}
\author{Nkouankam Aurelien}
\date{30/03/2025}

\usepackage{amsthm}
\newtheorem{theorem}{Theorem}[chapter]
\newtheorem{lemma}{Lemma}[chapter]
\newtheorem{corollary}{Corollary}[chapter]
\newtheorem{proposition}{Proposition}[chapter]
\newtheorem{conjecture}{Conjecture}[chapter]
\theoremstyle{definition}
\newtheorem{definition}{Definition}[chapter]
\theoremstyle{definition}
\newtheorem{example}{Example}[chapter]
\theoremstyle{definition}
\newtheorem{exercise}{Exercise}[chapter]
\theoremstyle{definition}
\newtheorem{hypothesis}{Hypothesis}[chapter]
\theoremstyle{remark}
\newtheorem*{remark}{Remark}
\newtheorem*{solution}{Solution}
\begin{document}
\maketitle

{
\hypersetup{linkcolor=}
\setcounter{tocdepth}{2}
\tableofcontents
}
\listoffigures
\listoftables
\hypertarget{introduction}{%
\chapter*{Introduction}\label{introduction}}
\addcontentsline{toc}{chapter}{Introduction}

La concentration des Gaz à Effet de Serre (GES) dans l'atmosphère et particulièrement le gaz carbonique (\(CO_2\)) a connu une nette augmentation à partir du début de l'ère industriel avec un taux de 277 parties par million (ppm) en 1750 \citep{Joos_Spahni2008}, pour 405 ± 0,1 ppm en 2017 \citep{LeQuere2018}. Cette augmentation exceptionnelle des émissions de GES pendant la période post-industrielle est principalement d'origine anthropique et a conduit à un brusque rechauffement global qui perturbe l'équilibre général de la planète. Des changements aussi soudain codéterminent l'impact du rechauffement global sur les systèmes naturels et socioéconomiques. L'une des conséquences les plus palpable et irréversible est la perte de la biodiversité. Les espèces endémiques des hauts lieux de la biodiversité ont un risque d'extinction très élevé et qui peut même doublé lorsque la température moyenne mondiale augmente de 1,5 °C à 2 °C \citep{Anonyme2022a}.En effet, une lente augmentation de température facilite l'adaptation aux changements contrairement aux augmentations brusques qui sont généralement à l'origine des cataclismes \citep{Joos_Spahni2008}. Les études menées par \citep{Griscom2017, Seymour_Busch2016} montrent qu'il est difficile d'atteindre les objectifs de reduction des émissions de GES sans la prise en compte des forêts, la reduction de la déforestation et la dégradation des forêts, la regenération des forêts et l'amélioration des méthodes de gestion forestière. De nos jours seules les forêts présentent une capacité d'absorption et de stockage de \(CO_2\) atmosphérique à une échelle significative pour contribuer aux efforts de reduction du rechauffement global. Les forêts absorbent environ 25\% des émissions de \(CO_2\) d'origine anthropique \citep{LeQuere2018}. A côté des prises de mesures drastiques telles que la reduction des émissions liées aux énergies fossiles, seules les forêts présentent de grands potentiels de reduction avec des bénéfices importants pour les sols, l'air, l'eau, la biodiversité et le développement \citep{VanGoor_Martijn2019}). Les forêts ont autrefois suscitées beaucoup d'interêts à cause des nombreuses ressources dont elles disposent notamment : les grumes, le bois énergie et les produits forestiers non ligneux. De même les inventaires forestiers ont généralement porté sur la densité des arbres par unité de surface, leur âge, leur taux de croissance, leur biomasse et la densité du bois des espèces ligneuses \citep{Watson2009}. Cependant le rôle des forêts dans la conservation de la biodiversité, le cycle de l'eau, la survie des populations riveraines, les cycles biogéochimiques, particulièrement le cycle du carbone et son lien avec les émissions de gaz à effet de serre a été moins bien mis en exsergue avant les premières manifestations des changements climatiques. Les forêts tropicales renferment environ 25\% du carbone de la biosphère représentant ainsi 33\% de la production primaire net \citep{Bonan2008}. Avec une superficie estimée à 178 millions d'hectares \citep{Vancutsem2021}, les forêts du bassin du Congo constituent le deuxième plus grand massif de forêts denses tropicales humides au monde après l'amazonie et représentent un reservoir exceptionnel de carbone et de biodiversité \citep{Ebaa2022}. Elles s'étendent des côtes du golfe de Guinée à l'ouest, aux montagnes du rift Albertin à l'est, et couvrent près de 7° de latitude de part et d'autre de l'équateur. En majeur partie, elles appartiennent à l'ensemble des forêts Guinéo-congolaises dont elles constituent plus de 80\% de la superficie \citep{Devers_Vande2006}. Elles offrent des moyens de subsistance à 100 millions de personnes qui y vivent ou résident à proximité \citep{Nasi2011} et contribuent plus indirectement à alimenter 40 millions d'autres personnes vivant dans les centres urbains avoisinant \citep{Ebaa2022}. Avec environ 60\% des réservoirs de carbone terrestre \citep{VanGoor_Martijn2019}, la biomasse forestière représente la quantité potentielle de carbone pouvant être ajouté au \(CO_2\) atmosphérique si elle venait à disparaitre \citep{Brown1997}. Les forêts d'Afrique centrale sequestrent environ 40 Gt de carbone \citep{Saatchi2011}. Elles ont une capacité de fixation de \(CO_2\) atmosphérique supérieure à celle des forêts amazonniènes malgré leur superficie comparativement plus reduite \citep{Hubau2020}. Cependant ces écosystèmes sont exposés à des processus de dégradation qui contribuent fortement à la diminution de leur étendue. En effet, depuis l'an 2000 on a enregistré une disparition de près 9\% de leur superficie soit 18 millions ha \citep{Ebaa2022}et d'après \citet{Seymour_Busch2016} , si aucune mesure n'est prise nous assisterons à une disparition de près de 2,9 millions de km² de forêts tropicales à l'horison 2050. Ces reductions de surfaces forestières sont liées à la progression des surfaces cultivées, la croissance démographique, le développement des infrastructures urbaines à proximité des forêts, l'exploitation forestière et l'extration du bois énergie \citep{Mayaux2013}.
En Afrique centrale la déforestation est beaucoup plus liée à des activités à petites échelles telles que: l'agriculture de subsistance, la production du charbon et le ramassage du bois de chauffe contrairement à l'amérique latine où des grands espaces sont directement concernés \citep{Curtis2018}.
Les menaces auxquelles sont exposées les forêts ont retenu l'attention de la communauté internationale en vue de developper des stratégies permettant la réduction des émissions liées à la déforestation et la dégradation des forêts.
Cette mobilisation s'est manifestée par l'institution de la Convention Cadre des Nations Unies pour les Changements Climatiques (CCNUCC), qui encourage les pays signataires à participer à l'aténuation des effets des changements climatiques, en mettant en oeuvre de manière volontaire, des processus visant à reduire les émissions de Gaz à Effet de Serre (GES) et à renforcer la fixation et la sequestration durable de ces gaz.
De même, elle recommande à tous les Etats parties de s'engager à développer et à publier périodiquement et selon des méthodes comparables, leurs inventaires nationaux des émissions et stockage de tout leur GES \citep{Brown2002}.
Pour faciliter la réalisation des engagements souscrits par les pays développés, trois mécanismes de flexibilités ont été mis sur pieds:

\begin{itemize}
\item
  le mécanisme des permis négociables qui provient directement du principe de préservation de la Convention Cadre des Nations Unies pour les Changements Climatiques (CCNUCC) et de la création du crédit carbone.
  Il correspond à la mise en place du marché carbone à différentes échelles où s'échangent les crédits carbones;
\item
  la Mise en Oeuvre Conjointe (MOC) qui est un mécanisme de financement de projets ayant pour objectif principal le stockage du carbone ou la réduction des GES. Il permet entre pays développés de procéder à des investissements visant à réduire les émissions de GES en dehors de leur territoire national et à bénéficier des crédits d'émission générés par les réductions ainsi suscitées;
\item
  le Mécanisme de Développement Propre (MDP) qui permet aux pays industrialisés et les entités qui en sont issues, d'aider au financement et à la réalisation de projets de réduction d'émissions dans les pays en développement.

  L'accord de Paris de 2015 a été décisif dans les négociations internationales sur le climat par l'initiation d'une approche systémique conduite par les pays pour atteindre des objectifss climatiques collectifs dans un cadre de transparence, de suivi et de communication.
  Les Contributions Déterminées au niveau National (CDN) issuent de cet accord représentent l'un des principaux instruments par lequel les contributions des forêts sont suivies par tous les pays du monde.
  De même l'agenda 2030 de l'Organisation des Nations Unies (ONU) avec ses 17 Objectifs de Développement Durable (ODD) et spécifiquement l'ODD 13 sur la lutte contre les changements climatiques représente également un instrument international par lequel l'importance des forêts peut être mise en évidence \citep{Mbairamadji2022}.
  Dans le secteur forestier les efforts de reduction des émissions des GES sont beaucoup plus manifestes à travers le processus REDD+ qui est un mécanisme de reduction des émissions de gaz liées à la déforestation et la dégradation des forêts dans les pays en développement.
  Ce processus a entre autre pour objectif de: conserver les stocks de carbone; gérer durablemeent les forêts et renforcer les stocks de carbone forestiers.
  L'idée soutenue par cette approche est d'établir un lien direct entre les incitations financières à la conservation des forêts et les stocks de carbone.
  Dans ce sens la REDD+ est souvent consideré comme un mécanisme à gain quadruple en ceci que les bénéfices climatiques obtenus pourraient être associés à une reduction de la pauvreté, une meilleure gouvernance forestière et à une conservation de la biodiversité \citep{Duchelle2018}.
  De même une protection exclusive des forêts tropicales sans l'amélioration des conditions de vies des couches défavorisées et les plus vulnérables ne saurait être éfficasse pour l'atteinte des objectifs de reduction des effets des changements climatiques \citep{Seymour_Busch2016}.
  En effet 28\% du revenu total des ménages des communautés riveraines proviennent des forêts \citep{Angelsen2014}. Des études recentes ont montré que stopper la déforestation et ses corollaires pourraient contribuer à une reduction des émissions de l'ordre de 37\% par rapport aux objectifs de 2030 stipulés dans l'accord de Paris sur le maintien du rechauffement global en dessous de 2°C au dessus du niveau de la période préindustrielle, la diminution des émissions de 43\% d'ici 2030 et l'atteinte de la neutralité carbone à l'horizon 2050 \citep{Griscom2017}.

  Le Cameroun a adhéré au processus REDD+ et a publié son document de proposition de mesures pour l'état de préparation (R-PP) de sa stratégie nationale en juin 2014. Dans ce document, il affirme son attachement à la stabilisation du climat et l'amélioration des conditions de vie des populations locales.Pour lui, le mécanisme REDD+ doit adopter une approche multisectorielle, intégrée et partenariale mettant en exergue tous les secteurs de développement et fondée sur une démarche participative et inclusive \citep{R-PP2013}. Cette transformation fait de la gestion des ressources naturelles un pôle de développement. Depuis 1994 le Cameroun a entamé une vaste réforme de sa politique forestière et la gestion durable des ressources forestières. Cette réforme s'est traduite par l'institution des outils de valorisation de l'implication des différents acteurs dans la gestion des forêts. Il s'agit entre autre, de la valorisation d'une foresterie sociale à travers la création des forêts communales \citep{Milol2008}. Ce concept renvoie aux formations végétales ayant fait l'objet d'un acte de classement pour le compte d'une commune ou qui ont été plantées par la collectivité concernée. Avec une superficie totale cumulée de 2327891 ha les forêts communales représentent 25\% de la superficie des forêts de production du Cameroun \citep{Anonyme2021i} et font partie intégrante du Domaine Forestier Permanent (DFP) en ce sens qu'elles sont définitivement affectées à la forêt. Cependant, la gestion de la plupart de ces forêts communales créées reste très inadéquate, en raison du fait que les acteurs des communes locales ne disposent ni outils, ni capacités suffisantes pour mettre en œuvre une gestion forestière durable, prenant en compte les opportunités qu'offrent les services écosystémiques et les problématiques actuels liées à la biodiversité et au changement climatique \citep{Anonyme2022b}. Bien que le Cameroun soit impliqué dans des programmes internationaux relatifs à la préparation de la REDD+, il n'a pas encore mis en place un système opérationnel de surveillance et de comptabilisation du carbone. En effet un système de comptabilisation et de surveillance du carbone est un outil d'estimation et de prédiction des stocks de carbone forestier, s'appuyant sur des scénarios des taux de perturbation futurs et des actions de gestion ainsi que les changements dans les stocks de carbone pour les prochaines décennies \citep{Watson2009}. En d'autre terme il s'agit d'un système de Mesure, de Notification et de Vérification (MRV) afin de générer des informations vérifiables et de qualité à des fins diverses sur l'ensemble des pools de carbone. Les forêts communales du Cameroun peuvent constituer des puits de carbone permettant de réduire l'impact de la déforestation sur le stock de carbone. De ce fait elles sont éligibles au MDP et la REDD+ dont la finalité est la gestion durable, la conservation des forêts, mais également la possibilité de bénéficier des Réductions d'Emission Certifiés (REC), encore appelées crédits carbone \citep{Cuny2011}. La mise en œuvre des programmes environnementaux tels que le mécanisme REDD+ implique les riverains et les propriétaires forestiers pour s'assurer que ces derniers participent activement à la gestion forestière. L'évaluation de la performance des systèmes d'aménagement et de gestion durable mis en place dans les forêts communales et suivant le contexte des changements climatiques, nécessite une tenue régulière des inventaires floristiques et la caractérisation des stocks de carbone dans ces espaces. La plupart des études antérieures ont généralement porté sur l'étude des impacts de l'exploitation forestère sur les stocks de carbone ( \citep{Kankeu2013, Nanda2018}. D'autres travaux se sont intéressés à l'étude de l'évolutiondu couvert végétal par des techniques de télédetection \citep{Momo2018, Djiongo2019, Zekeng2020} et peu d'études ont été faites sur la biodiversité et la croissance des ligneux en rapport avec les stocks de carbone. La pertinence de cette étude découle du fait qu'en ciblant la dynamique de la biodiversité et des stocks de carbone dans les forêts communales situées dans la zone de contact forêt-savane, elle permettra de mieux comprendre leur fonction potentielle dans l'atténuation du rechauffement climatique ainsi que les mécanismes qui gouvernent l'évolution de la biodiversité et des stocks de carbone dans un contexte de gestion forestière décentralisée. La question générale qui se dégage est la suivante : quels sont les mécanismes qui soutendent la dynamique de la biodiversité et les stocks de carbone dans les zones de contact forêt-savane et dans des systèmes de gestion forestière associés aux populations locales? spécifiquement il s'agit de répondre aux questions suivantes:
\item
  quelle est la dynamique spatio-temporelle de la végetation dans les forêts communales des zones de contact forêt-savane?
\item
  quelle est la nature de la diversité ligneuse dans ces espaces?
\item
  quelles sont les quantités de carbone stockées dans ces forêts?
\item
  où sont localisés les stocks de carbone les plus importants dans ces forêts?
\end{itemize}

La situation géographique de la zone de contact forêt-savane nous amène à émettre l'hypothèse selon laquelle l'évolution de ces écosystèmes pourrait être un indicateur de la nature des influences aux quelles elle est soumise notamment; les changements climatiques et les actions anthropiques.
D'autres hypothèses que nous pouvons dégagées sont les suivantes:

\begin{itemize}
\item
  les espaces occupés par les écosystèmes forestiers subissent une augmentation au détriment des écosystèmes savanicoles;
\item
  la dynamique spatio-temporelle va entrainer une augmentation de la composition et de la diversité ligneuse;
\item
  les stocks de carbone seront moins élevés comparativement aux espaces entièrement couverts de forêts;
\item
  les stocks de carbone les plus importants seront localisés dans les zones forestières les moins accessibles à l'homme.
\end{itemize}

L'objectif général de cette étude est d'analyser la biodiversité et la distribution des stocks de carbone dans les zones de contact forêt-savane.
Spécifiquement il s'agit de:

\begin{itemize}
\item
  comprendre la dynamique spatio-temporelle de la couverture végétale et prédire ses changements;
\item
  caractériser et évaluer la diversité ligneuse;
\item
  évaluer les stocks de carbone;
\item
  et cartographier la distribution spatiale des stocks de carbone.
\end{itemize}

\mainmatter

\hypertarget{guxe9nuxe9ralituxe9s}{%
\chapter{Généralités}\label{guxe9nuxe9ralituxe9s}}

\hypertarget{pruxe9sentation-de-la-foruxeat-communale-de-nanga-eboko}{%
\section{Présentation de la Forêt Communale de Nanga-Eboko}\label{pruxe9sentation-de-la-foruxeat-communale-de-nanga-eboko}}

\hypertarget{localisation}{%
\subsection{Localisation}\label{localisation}}

La forêt communale de Nanga-Eboko est constituée de deux blocs distincts (Figure: \ref{fig:Localisation-FCNanga1}) couvrant une superficie totale de 20000 ha (11000 ha et 9000 ha). C'est le décret n°2009/2046/PM du 16 Novembre 2009 qui marque sa création et confi sa tutelle à la mairie de Nanga-Eboko. Elle est située dans la région du Centre, Département de la Haute Sanaga, Arrondissement de Nanga-Eboko \citep{Yemafack_Alemagi2016}.

\begin{figure}

{\centering \includegraphics[width=0.8\linewidth]{images/localisation_nanga-01} 

}

\caption{Localisation de la Forêt communale de Nanga-Eboko (Yemefack et Alemagi, 2016)}\label{fig:Localisation-FCNanga1}
\end{figure}

\hypertarget{milieu-biophysique}{%
\subsection{Milieu biophysique}\label{milieu-biophysique}}

\hypertarget{relief-et-hydrographie}{%
\subsubsection{Relief et hydrographie}\label{relief-et-hydrographie}}

C'est le relief inégal du grand plateau méridional qui est une vaste surface d'érosion en gradins ou en \enquote{mosaïque irrégulière de blocs} couvrant tout le Centre-Sud et le Sud du Cameroun \citep{Martin1966}. Ce sont des formations précambriennes plissées puis arasées par l'érosion \citep{Servant2000}. Elles ont été mises en place sur une surface subhorizontale où les altitudes varient entre 600 et 750 m. Avec son aspect vallonnée, le réseau hydrographique y est très ramifié et les eaux s'écoulent vers la Sanaga axe principal de drainage du plateau sud-Camerounais. Le principal cours d'eau est la Njékè et ses affluents: Menam, Abomvé, Membé et Tigbé qui drainent toute la FC. la Njékè est un affluent de la Sanaga, son bassin s'étend surtout en savane jusqu'à la plaine Yoko.

\hypertarget{climat}{%
\subsubsection{Climat}\label{climat}}

Par son appartenance au domaine subéquatorial, la Forêt Communale de Nanga-Eboko connait quatre saisons: deux saisons pluvieuses séparées par deux saisons sèches d'inégales importances. C'est le climat équatorial de type guinéen. La topographie de cuvette et l'abaissement imporatant de l'altitude par rapport aux hautes terres voisines favorisent l'échauffement et l'évolution dynamique diurne. Par conséquent c'est une zone relativement chaude: la température moyenne annuelle se situe autour de 24°C avec une amplitude de 3,4°C. Les fortes chaleurs sont enrégitrés pendant les mois de janvier, février et mars (25°C-26°C) tandis que les basses températures sont relevées en juillet et août (22,6°C - 22,8°C). Les valeurs de pluviosité sont faibles avec des précipitations oscillant entre 1113 et 1827 mm. Avec des pluies étalées d'avril à juin et de mi-août à novembre, le nombre de jours pluvieux sur l'ensemble de la zone d'étude est de 133 jours.

\hypertarget{vuxe9guxe9tation}{%
\subsubsection{Végétation}\label{vuxe9guxe9tation}}

La végétation de la Forêt Communale de Nanga-Eboko se caractérise par une forêt de type tropicale semi décidue guinéo-congololaise \citep{Letouzey1985} constituée de strates arborescentes, arbustives et d'espèces lianescentes. Elle appartient au domaine des forêts secondaires à savanes péri-forestières plus ou moins dégradées sur terres forestières \citep{Onana2019}. Le massif forestier renferme de nombreuses espèces de la famille des Sterculiacées telles que: \emph{Triplochiton scleroxylon}, \emph{Sterculia oblonga} et autres \emph{Sterculia}, \emph{Pterygota kamerunensis} et divers \emph{Cola}. Dans les autres familles on relève divers \emph{Celtis}, \emph{Chlorophora excelsa}, de nombreux \emph{Ficus} et \emph{Terminalia superba}. En dehors des forêts galerie,les autres surfaces forestières apparaissent dégradées soit par l'exploitation forestière ou les défrichements avec jachères forestières à divers stades de régénération, soit de forêts et brousses forestières reconstituées à partir de la savane. Elles sont caractérisées par les espèces de jachères suivantes: \emph{Musanga cecropiodes}, \emph{Albizzia zigia}, \emph{Anthocleista nobilis}, \emph{Xylopia aethiopica}, \emph{Sarcocephalus esculentus}, \emph{Ceiba pentadra} et \emph{Vernonia amygdalina}. A certains endroit au Nord-ouest de la Forêt Communale (village Bisaga), l'on peut observé des champs cultivés à partir des défrichements de savanes et des brousses forestières fermées caractérisées par les espèces telles que: \emph{Anthocleista nobilis}, \emph{Sarcocephalus esculentus}, \emph{Albizzia zigia}, \emph{Vernonia amygdalina}, \emph{Harungana madascariensis}. Quelques rares pieds de \emph{Terminalia glaucescens}, \emph{Annona arenaria} et \emph{Piliostigma reticulata} témoignent d'une présence ancienne de la savane. Les savanes proprement dites sont arborées et caractérisées par les espèces telles que: \emph{Bridelia ferruginea}, \emph{Crossopteryx febrifuga},\emph{Himinocardia acida}, \emph{Annona arenaria}, \emph{Piliostigma reticulata}. La graminéé dominante des savanes est \emph{Hyparrhenia rufa}. Les espèces comme \emph{Pennisetum purpureum} sont répandus dans les jachères et sur les termitières. Dans les bas-fonds marécageux on peut observer des formations forestières dominées par les \emph{Raphias}.

\hypertarget{sols}{%
\subsubsection{Sols}\label{sols}}

La grande majorité des sols de la Forêt Communale de Nanga-Eboko appartiennent à la sous-classe des sols ferralitiques et aux catégories des sols ferralitiques typiques et sols ferralitiques indurés. Ce sont éssentiellement des sols argileux présentant une richesse organique moyenne et richesse minérale faible leur conférant ainsi une fertilité intrinsèque moyenne \citep{Martin1966}. Dans les bas-fonds marécageux et les vallées des grandes rivières, on peut noter la présence des sols hydromorphes: ils ne représentent que des surfaces reduites et leur valeur agricole est généralement très médiocre.

\hypertarget{environnement-socio-uxe9conomique}{%
\subsubsection{Environnement socio-économique}\label{environnement-socio-uxe9conomique}}

La population autour de la Forêt Communale est hétérogène avec trois tribus (Vouté, Baveuk et Yangavok) occupant 23 villages appartenant à deux groupements: Emtsé et Bissaga. Le groupement Emtsé est le plus important en terme de population avec plus de 6000 habitants. Dans cet ensemble les femmes sont les plus nombreuses et représentent près de 52\% de la population. Chaque village est soumis à l'autorité d'un chef. Ce dernier est élu par les populations et les canditats à l'élection proviennent des descendants de la famille reignante. Le chef est garant du pouvoir traditionnel et assure le lien entre l'administration et les populations. Il dernier règle les problèmes de tenure foncière et tout autre type de conflits qui sont de deux types : les conflits en vue de l'appropriation des arbres, et les divergences entre les membres d'une famille ou tribu.

L'agriculture est la principale activité des riverains de la FC de Nanga-Eboko. Les cultures vivrières occupent la première place dans les activités agricoles. Bien que les techniques agricoles soient encore rudimentaires, le maïs occupe une place prépondérante dans ces cultures étant donné que l'alimentation de base de la zone est le couscous de maïs. Le cacao et le café sont les principales sources de revenus des populations. L'élevage est de faible importance et concerne surtout les caprins, les ovins et la volaille. La chasse et la pêche sont des activités de moyenne importance et fournissent des revenus non négligeables à ceux qui les pratiquent.

\hypertarget{pruxe9sentation-de-la-foruxeat-communale-de-ndikinimuxe9ki}{%
\section{Présentation de la Forêt Communale de Ndikiniméki}\label{pruxe9sentation-de-la-foruxeat-communale-de-ndikinimuxe9ki}}

\hypertarget{localisation-1}{%
\subsection{Localisation}\label{localisation-1}}

La Forêt communale de Ndikiniméki est située dans la région du Centre, Département du Mbam et Inoubou, Arrondissement de Ndikiniméki. Elle couvre une superficie de 20000 ha repartie en deux blocs (17000 ha et 3000 ha) entièrement gérés par la commune de Ndikiniméki (Figure: \ref{fig:Localisation-FCNdikinimeki1}). Cette commune est limitée au Nord par l'arrondissement de Makénéné, à l'Ouest par l'arrondissement de Yingui, au Sud par Ngambé et Yingui et à l'Est par les arrondissements de Ndom, Bokito, Nitoukou et Bafia.

\begin{figure}

{\centering \includegraphics[width=0.8\linewidth]{images/localisation_Ndiki-01} 

}

\caption{Localisation de la Forêt Communale de Ndikiniméki (Centre Technique de la Forêt Communale Cameroun, 2008)}\label{fig:Localisation-FCNdikinimeki1}
\end{figure}

\hypertarget{milieu-biophysique-1}{%
\subsection{Milieu biophysique}\label{milieu-biophysique-1}}

\hypertarget{relief-et-hydrographie-1}{%
\subsubsection{Relief et hydrographie}\label{relief-et-hydrographie-1}}

Le relief de la Forêt Communale de Ndikiniméki correspond au relief de la frange occidentale du plateau sud-Camerounais. C'est un relief escarpé et coupé par de nombreuses vallées très profondes. Il est rattaché dans sa partie Est au plateau Sud-Camerounais par la dépression de Bafia et dans sa partie Ouest aux plateaux Bamilékés et Bamoun où il prend un aspect réellement montagneux \citep{Olivry1986}. Le réseau hydrographique est dense et constitué de cours d'eaux permanents et aux débits variables (Nting-nting, Molo, Bassaben, Biadjang, Makombé\ldots). Le Makombé est la rivière la plus importante de la zone et constitue l'un des affluents les plus importants du Wouri dans sa partie continentale.

\hypertarget{climat-1}{%
\subsubsection{Climat}\label{climat-1}}

Le climat est équatorial de type guinéen, caractérisé par deux saisons sèches et deux saisons de pluies. La moyenne thermique se situe entre 20 et 25°C avec un maxima de 30°C au mois de février. Les précipitations moyennes annuelles se situent autour de 1440 mm.

\hypertarget{vuxe9guxe9tation-1}{%
\subsubsection{Végétation}\label{vuxe9guxe9tation-1}}

La végétation rencontrée d'une part est celle de la forêt ombrophile de basse et moyenne altitude à Lophira alata \citep{Letouzey1958}. Elle est constituée d'espèces de bois d'œuvre, de bois de chauffe et d'espèces à usage multiples. Quelques espèces rencontrées sont : le Movingui (\emph{Distemonanthus benthamanus}), le Padouk (\emph{Ptérocarpus soyauxii}), le Tali (\emph{Erythrophleum suaveolens}), le Sapelli (\emph{Entandrophragma cylindicum}), l'Okan (\emph{Cilicodiscus gabonensis}), le Fraké (\emph{Terminalia superba}), l'Iroko (\emph{Chlorophora excelsa}) et l'Ilumba (\emph{Pycnanthus angolensis}). Les familles typiques sont entre autre : les Violaceae, Euphorbiaceae, Anonaceae, Fabaceae et Ebénaceae. D'autre part on rencontre des jachères et des cultures. Les cultures sont vivrières (arachides, manioc, maïs, etc.) et pérennes (cacao). La faune terrestre est abondante et diversifiée et comprend : les mammifères, les oiseaux et les reptiles. La faune aquatique est riche et comprend : les clarias, les carpes, les capitaines, etc.

\hypertarget{sols-1}{%
\subsubsection{Sols}\label{sols-1}}

La grande majorité des sols de la FC de Ndikiniméki appartient à la classe des sols à sesquioxydes faiblement ferralitiques modaux sur roches acides. Dans les zones de savane on peut distinguer les sols ferraliques indurés caractérisés par des affleurements de dalles ou de gravillons. Les sols hydromorphes sont localisés dans les bas-fonds marécageux et les vallées \citep{Anonyme2012c}.

\hypertarget{environnement-socio-uxe9conomique-1}{%
\subsubsection{Environnement socio-économique}\label{environnement-socio-uxe9conomique-1}}

La population est constituée essentiellement de l'ethnie Banen, les autres ethnies étant halogènes. C'est une population cosmopolite qui vit pour sa grande majorité le long de la route nationale numéro 4, loin de la forêt à raison des troubles qu'a connue la zone à la veille des indépendances. Les habitants des villages qui autrefois jouxtaient la forêt communale (Ndokbou I et II) reviennent de temps en temps pratiquer l'agriculture à la lisière de la Forêt Communale. La commune est peu peuplée avec une densité de 19 habitants au km² pour une superficie de 2650 km². Chaque village est sous l'autorité d'un chef assisté par des notables. Les élites, les élus locaux, les patriarches et les responsables politiques, jouent également un rôle important dans les prises de décisions au sein des communautés. La vie des populations locales dépend énormément de la Forêt Communale. Elles y prélèvent des ressources pour leur alimentation, la pharmacopée, l'artisanat, l'énergie et le commerce. La pratique de l'agriculture est rudimentaire avec l'usage de très peu de fertilisants et de semences améliorées. L'essentiel de la production vivrière est destinée à la consommation tandis que la production de rente (cacao) est destinée à la vente. Les produits de la chasse et de la pêche sont généralement boucanés avant d'être vendus sur les marchés locaux (Makénéné et Ndikiniméki) ou Yaoundé.

\hypertarget{duxe9finition-des-concepts}{%
\section{Définition des concepts}\label{duxe9finition-des-concepts}}

\hypertarget{biodiversituxe9}{%
\subsection{Biodiversité}\label{biodiversituxe9}}

D'après la Convention des Nations Unies sur la Diversité Biologique , la biodiversité peut se définir comme étant la diversité des écosystèmes, des espèces et des gènes dans l'espace et dans le temps, ainsi que les interactions au sein de ces niveaux d'organisaion. Loin d'être un catalogue statistique des espèces, elle est le resultat de la dynamique du monde vivant (Figure: \ref{fig:Concept-Biodiversite1}).

\begin{figure}

{\centering \includegraphics[width=0.8\linewidth]{images/Biodiversite} 

}

\caption{Concept de biodiversité d'après di Castri et Younes, 1996.}\label{fig:Concept-Biodiversite1}
\end{figure}

Pour les scientifiques il existe trois niveaux d'organisation de la biodiversité: la diversité génétique en tant que diversité de l'information génétique portée par les organismes vivants;la diversité des espèces qui correspond à l'ensemble des être vivants identifiables dans le monde vivant;la diversité des écosystèmes qui correspond aux communautés animales et végétales vivants sous différentes conditions climatiques et géographiques \citep{Leveque2008}. Le terme \enquote{biodiversité} qui est la forme réduite de \enquote{diversité biologique}a été employé pour la première fois par des naturalistes inquiets de la destruction rapide des milieux naturels et des forêts tropicales humides en particulier. La grande majorité des espèces présentes sur Terre sont encore inconnues, de même leur rôle dans le fonctionnement des écosystèmes reste encore à déterminer. Au coeur d'écosystèmes évoluant en permanence sous l'effet des contraintes climatiques, géologiques et anthropiques, la diversité se crée soit par des mutations génétiques soit par des processus de selections naturelles.

\hypertarget{estimation-des-stocks-de-carbone}{%
\subsection{Estimation des stocks de carbone}\label{estimation-des-stocks-de-carbone}}

La capacité à pouvoir mesurer avec exactitude et précision, les stocks de carbone séquestrés dans les forêts et avec des méthodes comparables fait l'objet de nombreuses recherches. Beaucoup de pays tropicaux disposent de très peu de données d'inventaires de tout ou partie de leur surface forestière. La plupart des inventaires sont anciens et n'ont pas souvent fait l'objet d'une actualisation \citep{Brown2002, Brown1997}. Les méthodes d'estimation de stock de carbone les plus rependues concernent surtout la biomasse aérienne constituée des branches (houppier) et le tronc (Figure: \ref{fig:Stockage-Biomassearbre1}). Elles reposent pour la plupart sur des méthodes biométriques pour ce qui concerne les forêts tropicales, et consistent généralement à convertir les données d'inventaires en stocks de carbone \citep{Loubota2018}.

\begin{figure}

{\centering \includegraphics[width=0.8\linewidth]{images/stockage biomasse arbre} 

}

\caption{Stockage de biomasse au niveau de l'arbre (adapté de De Madron, 2013).}\label{fig:Stockage-Biomassearbre1}
\end{figure}

Pour l'évaluation des stocks de carbone, \citet{Penman2003} et \citet{MacDicken1997} définissent 5 réservoirs repartis en 3 groupes pour la quantification du carbone:

\begin{itemize}
\tightlist
\item
  Groupe1: biomasse vivante. Il englobe la biomasse aérienne constituée des arbres et des autres végétaux et la biomasse souterraine, ou les racines;
\item
  Groupe 2: matière organique morte. Il comprend le bois mort;
\item
  Groupe 3: carbone organique du sol. Il inclue le carbone du sol.
\end{itemize}

La première technique employée par \citet{Brown_Lugo1992}, consiste à estimer la biomasse à partir des volumes de bois dits « commerciaux » (volume de bois fort sur écorce) corrigés par un coefficient d'expansion (Figure: \ref{fig:Methode-BrownLugo1}). Elle s'applique pour les inventaires ayant un diamètre de précomptage égale à 10 cm.

\begin{figure}

{\centering \includegraphics[width=0.8\linewidth]{images/Methode Brown et Lugo} 

}

\caption{Méthode de Brown et Lugo (1992).}\label{fig:Methode-BrownLugo1}
\end{figure}

Cependant, son application pour des inventaires ayant un diamètre de précomptage allant jusqu'à 30 cm nécessite l'ajout d'un facteur correctif (VEF : \enquote*{Volume Expansion Factor}) permettant d'extrapoler le volume des classes de diamètre manquantes \eqref{eq:Brown-Lugo1}.

\begin{equation}
AGB=VOB×WB×BEF(VEF)
\label{eq:Brown-Lugo1}
\end{equation}

\(AGB\) : Biomasse épigée; \(VOB\) : Volume bois fort; \(WD\) : Densité du bois; \(VEF\):Facteur d'expansion de volume.

La seconde technique (Figure: \ref{fig:Methode-Chave2014}) est celle consistant à estimer la biomasse à partir des tables d'inventaires et d'un modèle allométrique \citep{Chave2005}. Elle repose sur l'utilisation du résultat d'un inventaire en nombre de tiges par essence et par catégorie de diamètre auquel on associe la biomasse épigée totale d'un arbre (AGB) estimée à partir d'une ou plusieurs données d'entrées : le Diamètre (D), la Hauteur (H), la Densité du bois (ρ). Equation: \eqref{eq:Chave-al2005}.

\begin{equation}
AGB=0,0673×(ρD^2 H)0,976≡0,559×(ρD^2 H)
\label{eq:Chave-al2005}
\end{equation}

\(D\) : diamètre à hauteur de poitrine ou au-dessus des contreforts (m); \(H\) : hauteur (m); \(ρ\) : densité spécifique (g/cm\textsuperscript{3}).

{[}l'intervalle de validité de l'équation: 5 à 156 cm de diamètre{]}.

Lorsque les données de hauteur ne sont pas disponibles, \citet{Chave2014} a récemment développé une équation permettant d'estimer la biomasse épigée à partir du diamètre de l'arbre et un facteur de stress environnemental. Equation: \eqref{eq:Chave-al2014}.

\begin{equation}
AGB=℮^[-1,803-0,976×E+0,976×ln⁡(ρ)+2,673×ln⁡(D)-0,0299×(ln(D))²]
\label{eq:Chave-al2014}
\end{equation}

\(AGB\), \(ρ\), \(D\), \(H\) et \(E\) sont respectivement la biomasse aérienne (kg), la densité du bois spécifique (g/cm\textsuperscript{3}), le diamètre de l'arbre (cm), la hauteur totale de l'arbre (m) et une variable qui quantifie le stress environnemental (elle varie entre -0,2 et 1). {[}L'intervalle de validité de cette équation est de 5 à 180 cm de diamètre{]}.

Comparativement à la technique de \citet{Brown_Lugo1992}, les équations de Chave permettent d'estimer la biomasse totale de l'arbre en fonction du diamètre et/ou la hauteur et un facteur de stress environnemental sans appliquer de facteur de conversion global (Figure: \ref{fig:Methode-Chave2014}). C'est une méthode qui s'applique aisément aux données brutes d'inventaires forestiers et est très adaptée aux dispositifs scientifiques du type « placettes permanentes ». Cependant sa mise en œuvre exige des corrections d'imprécision liées à l'affectation de chaque tige à une classe de diamètre plus ou moins large (5 ou 10 cm).

\begin{figure}

{\centering \includegraphics[width=0.8\linewidth]{images/Chave 2014-04} 

}

\caption{Méthode de Chave et al.(2014).}\label{fig:Methode-Chave2014}
\end{figure}

\hypertarget{zone-de-transition}{%
\subsection{Zone de transition}\label{zone-de-transition}}

La zone de transition au Cameroun correspond à la zone de contact forêt-savane. C'est un vaste domaine situé entre le sixième et le quatrième degré de latitude nord, et entre le onzième et le seizième degré de longitude est. Avec une altitude comprise entre 650 et 900 m, son relief est formé d'une succession de collines convexes et d'interfluves émoussés , parfois cuirassés \citep{Morin1979}. La végétation dominante est une forêt semi-décidue en équilibre instable avec les savanes périforestières. Le caractère semi-caducifolié de cette forêt constitue un trait dominant avec les espèces arborescentes caractéristiques telles que les Sterculiacées et les Celtis \citep{Olivry1986}.

\hypertarget{evolution-du-processus-redd}{%
\subsection{Evolution du processus REDD+}\label{evolution-du-processus-redd}}

La stratégie REDD+ est un ensemble de politiques, d'actions locales, infranationales, nationales et mondiales dont l'objectif premier est de réduire les émissions liées à la déforestation et la dégradation des forêts et d'augmenter les stocks de carbone forestier dans les pays en développement. Cette définition relève que le processus REDD+ est liée aux émissions et aux piégeages des Gaz à Effet de Serre (GES), et que les actions entreprises devraient englober des paiements conditionnels ou basés sur les résultats \citep{Angelsen2013}. En théorie, l'architecture globale de la stratégie REDD+ peut être comparée à un système qui consiste à payer les propriétaires forestiers et les usagers de la forêt sur la base de leurs performances en réduction des émissions \citep{Angelsen2010}. Ces paiements sont effectués conformément au respect d'un régime d'occupation des terres ou d'actions menés convenus et prévoient des restrictions d'usage. La stratégie REDD+ est un mécanisme qui a été initié à Montréal en 2005 lors de la 11e Conférence des parties (CdP-11) à la Convention Cadre des Nations Unies pour les Changements Climatiques (CCNUCC) par la Papouasie-Nouvelle-Guinée et le Costa Rica. Ces pays s'étaient appuyés sur une publication produite par un consortium d'organisations non gouvernementales, de centres de recherche brésiliens et l'université de Yale aux USA. Ce rapport rappelait le rôle clé joué par la déforestation dans les émissions de GES, mettait en évidence l'opportunité d'une approche sectorielle de ce problème et proposait le concept de « réduction d'émissions compensées » \citep{Santilli2005}. En effet, sur la totalité des émissions annuelles de GES générées par les activités humaines depuis 1990, 20\% proviennent des changements d'usage des sols \citep{GIEC2007} et en particulier de la déforestation et de la dégradation des forêts en milieu tropical. L'engagement et la participation du Cameroun à ce processus s'est concrétisé par la création du comité de pilotage REDD+ en juin 2012.

\hypertarget{matuxe9riel-et-muxe9thodes}{%
\chapter{Matériel et méthodes}\label{matuxe9riel-et-muxe9thodes}}

\hypertarget{matuxe9riel}{%
\section{Matériel}\label{matuxe9riel}}

\hypertarget{sites-de-luxe9tude}{%
\subsection{Sites de l'étude}\label{sites-de-luxe9tude}}

Les Forêts Communales de Nanga et Boko et Ndikiniméki sont toutes deux situées dans la région du centre Cameroun. C'est une zone correspondant à un climat équatorial de type guinéen. La végétation rencontrée est respectivement une forêt de type tropicale semi décidue guinéo-congololaise \citep{Letouzey1985} et une forêt ombrophile de basse et moyenne altitude à Lophira alata \citep{Letouzey1958}. Le choix de ces Forêts Communales s'est fait sur la base de: leur statut juridique (existence d'un acte de classement); la validation de leur plan d'aménagement; leur accessibilité (l'existence des pistes forestières qui ont facilité les choix et l'installation des placettes). D'autres part, le fait pour ces forêts d'être soit en arrêt d'exploitation (FC Nanga-Eboko) ou pas encore en exploitation (FC Ndikiniméki), constitue un atout indéniable pour l'élaboration des critères et indicateurs pour la désignation des sites de carbone, l'évaluation et le suivi du carbone dans les forêts communales. La formulation des Critères et Indicateurs (C\&I) de désignation des sites de carbone dans cette étude a été basée à la fois sur la gestion de l'activité forestière courante dans la forêt, la structure du peuplement, et la gestion des écosystèmes \citep{Nivert2012}.

\hypertarget{matuxe9riel-technique}{%
\subsection{Matériel technique}\label{matuxe9riel-technique}}

Pour atteindre les objectifs fixés, le matériel inscrits dans le Tableau : \ref{tab:Liste-desoutils2} a été utilisé

\begin{table}
\centering
\caption{\label{tab:Liste-desoutils2}Liste des outils utilisés et leurs fonctions}
\centering
\resizebox{\ifdim\width>\linewidth\linewidth\else\width\fi}{!}{
\begin{tabular}[t]{ll}
\toprule
Outils & Fonctions\\
\midrule
Global Positionning System (GPS) & enregistrement des coordonnées géographiques des points et des parcelles d'échantillonnage\\
Boussole & Orientation des layons d'ouverture\\
Appareil photo & Prise de vues\\
Machette & ouverture des layons et confection des jalons\\
Penta-décamètre de 50 m & Délimitation des aires d'échantillonnage\\
\addlinespace
Compas forestier & Mesure des diamètres ou circonférences des arbres\\
Ficelle & Délimitation des sous-parcelles\\
Clinomètre & Mesure les pentes des parcelles\\
Sécateur & Collecte des échantillons sur le terrain\\
Etuve & Séchage des échantillons\\
\addlinespace
Marqueurs & Etiquetage des échantillons collectés\\
Plastique polyéthylique & Collecte des échantillons de litière et sol\\
Planche à herbier, papiers-presse et alcool & Pressage, séchage et conservation des spécimens\\
Balance & Pesage des échantillons\\
Carnets, crayons et gommes & prises de notes\\
\bottomrule
\end{tabular}}
\end{table}

\hypertarget{muxe9thodes}{%
\section{Méthodes}\label{muxe9thodes}}

La méthodologie utilisée dans cette étude a consisté à une combinaison de trois opérations. La première, à partir des observations et l'analyse des données prélevées dans un ensemble de parcelles préétablies, a permis de décrire la composition floristique et la structure des peuplements forestiers (caractérisation et évaluation de la biodiversité). La seconde a consisté à mesurer les volumes de biomasse dans les pools préalablement définis pour en déduire les stocks de carbone existants (estimation des stocks de carbone). La dernière s'est focalisée sur la classification des images multibandes correspondantes aux zones d'étude pour générer des cartes de distribution de carbone.

\hypertarget{caractuxe9risation-et-uxe9valuation-de-la-biodiversituxe9}{%
\subsection{Caractérisation et évaluation de la biodiversité}\label{caractuxe9risation-et-uxe9valuation-de-la-biodiversituxe9}}

La caractérisation et l'évaluation de la biodiversité a nécessité l'établissement d'un dispositif de sondage d'inventaires forestiers. De manière pratique c'est une opération qui exige la mise sur pied des parcelles d'inventaires suivant un plan de sondage. Etant représentatif de grandes surfaces d'inventaires, l'échantillonnage a été le meilleur moyen du point de vue statistique, d'évaluer la structure et la composition des peuplements forestiers. Il a consisté à constituer à partir d'une population, un échantillon de « n » unités et chaque combinaison d'unités ayant la même probabilité d'être sélectionnée. Dans notre contexte, l'échantillonnage stratifié randomisé s'est avéré plus précis et moins couteux pour un inventaire de la biodiversité associé à l'estimation des stocks de carbone en comparaison aux autres modèles d'échantillonnage \citep{Anonyme2009d, MacDicken1997}. L'application de ce modèle a nécessité une stratification ou découpage spatiale des Forêts Communales en zones aussi homogènes que possible \citep{Picard_Gourlet-Fleury2008}. Chaque strate ou sous-population pouvant être définie par son type d'occupation de sol. Pour se faire, nous nous sommes appuyés sur des supports cartographiques de la végétation issus des images satellites vérifiées sur terrain, des photos aériennes et des cartes de sol ou de la topographie. Cette stratification a également pris en compte des critères qui sont liés aux variables mesurés et surveillés, à l'exemple des réservoirs de carbone. Les normes de stratification forestière élaborée en 2016 en République Démocratique du Congo \citep{DIAF2016}, les normes d'inventaire d'aménagement et de préinvestissement du Cameroun \citep{ONADEF1991} et le manuel pour le relevé intégré de données sur le terrain \citep{FAO2009} ont servi de base pour identifier les différentes unités d'occupation dans le présent travail.

\hypertarget{dispositif-duxe9chantillonnage}{%
\subsubsection{Dispositif d'échantillonnage}\label{dispositif-duxe9chantillonnage}}

Le dispositif d'échantillonnage utilisé dans le cadre de ce travail est un ensemble de placettes proposées par \citet{Hairiah2011} dans le cadre du partenariat ASB (Alternative to Slash and Burn) Figure: \ref{fig:Placettes-echantillonnages2}. En effet, cette méthode préconise l'installation des placettes principales rectangulaires de 2 000 m² (50 m × 40 m) pour l'échantillonnage des ligneux à grands diamètres (\(DHP\ge 30 cm\)) et des placettes secondaires rectangulaires de 200 m² (5 m × 40 m) pour l'échantillonnage des ligneux à petits diamètres (\(5 cm\le DHP < 30 cm\)).

\begin{figure}

{\centering \includegraphics[width=0.8\linewidth]{images/Placettes_echantillonnages} 

}

\caption{Placettes d'échantillonnages (Adaptée de Hairiah et al.(2010)}\label{fig:Placettes-echantillonnages2}
\end{figure}

Dans chaque parcelle, un inventaire a été réalisé dans le but d'identifier toutes les espèces ligneuses présentes. Les mesures dendrométriques faites étaient la hauteur de l'arbre et le Diamètre à Hauteur de poitrine (DHP). Dans chaque parcelle et sous-parcelle, tous les arbres de diamètre du tronc à hauteur de poitrine (1,30 m du sol et 0,30 m au-dessus des défauts) pour les individus de diamètre supérieurs à 10 cm et à 50 cm du sol pour les individus de diamètre inférieurs à 10 cm \citep{Zapfack2005, Hairiah2011} ont été mesurés en utilisant un dhp-mètre.

\hypertarget{installation-des-parcelles}{%
\subsubsection{Installation des parcelles}\label{installation-des-parcelles}}

Les parcelles ont été installées en respectant les conditions moyennes de densité des espèces proposées par \citet{Picard_Gourlet-Fleury2008}. En effet ce sont des parcelles de 0,2 à 0,5 ha ayant fait l'objet d'une stratification selon le milieu et la composition spécifique Figure: \ref{fig:Delimitation-parcelles2} . La mise en place des placettes s'est faite à partir de coordonnées GPS des points issus de la classification non supervisée des images satellites de la zone d'étude. La «méthode 3-4-5» a été utilisée pour la mise en place des parcelles. Une fois le point de reférence repéré, un jalon est fixé (Figure: \ref{fig:Protocole-placettes2}a). Ensuite un second jalon est placé à 4 m du premier suivant la direction de base principale (Figure: \ref{fig:Protocole-placettes2}b). Par la suite, on mesure 3 m à partir du premier jalon suivant une direction perpendiculaire à l'axe principal, puis on fixe le troisième jalon lorsque la distance qui sépare ce dernier du second est égale à 5 m (Figure: \ref{fig:Protocole-placettes2}c). Ainsi, l'extrémité de la placette formera un angle droit et on pourra matérialiser le reste de la surface aux dimensions requises (40m x 50m) par une ficelle colorée (Figure: \ref{fig:Protocole-placettes2}).

\begin{figure}

{\centering \includegraphics[width=0.8\linewidth]{images/Mise_en_place_placettes} 

}

\caption{Protocole de mise en place des placettes}\label{fig:Protocole-placettes2}
\end{figure}

\begin{figure}

{\centering \includegraphics[width=0.8\linewidth]{images/installation_parcelle} 

}

\caption{Délimitation d'une parcelle d'échantillonnage}\label{fig:Delimitation-parcelles2}
\end{figure}

\hypertarget{mesure-des-paramuxe8tres-dendromuxe9triques-des-arbres}{%
\subsection{Mesure des paramètres dendrométriques des arbres}\label{mesure-des-paramuxe8tres-dendromuxe9triques-des-arbres}}

Sur chaque placette installée, les informations générales (coordonnées géographiques, description du paysage, etc.) ont été enregistrées, puis les parcelles principales et secondaires ont été délimitées. Pour chaque arbre identifié, les paramètres tels que : le diamètre, les noms scientifiques ou vernaculaires et/ou les noms commerciaux ont été notés (Figure: \ref{fig:Mesure-diametre2}). Tous les arbres ont été identifiés sur la base de leurs caractères discriminants en utilisant les flores et les clés d'identification. Pour limiter le risque de double comptage des ligneux dans les parcelles en forêt, tous les arbres ont été marqués à la peinture. Des échantillons d'herbier des espèces non identifiées sur le terrain ont été collectés, pressés, puis conservés dans de l'alcool 70°. Ces derniers ont été utilisés pour confirmer plus tard leur identification à l'Herbier National de Yaoundé. La nomenclature taxonomique adoptée a été la classification botanique phylogénétique des angiospermes \citep{APGIII2009}.

\begin{figure}

{\centering \includegraphics[width=0.8\linewidth]{images/mesures_dendrometrriques} 

}

\caption{Mesure du diamètre des arbres}\label{fig:Mesure-diametre2}
\end{figure}

\hypertarget{muxe9thode-destimation-de-la-biomasse-auxe9rienne}{%
\subsection{Méthode d'estimation de la biomasse aérienne}\label{muxe9thode-destimation-de-la-biomasse-auxe9rienne}}

Pour l'évaluation des stocks de carbone, \citet{CIFOR2009}, \citet{Penman2003} et \citet{MacDicken1997} définissent 5 réservoirs repartis en 3 groupes pour la quantification du carbone :

\begin{itemize}
\item
  groupe 1 : biomasse vivante. Il englobe la biomasse aérienne constituée des arbres et des autres végétaux et la biomasse souterraine, ou les racines ;
\item
  groupe 2 : matière organique morte. Il comprend le bois mort ;
\item
  groupe 3 : carbone organique du sol. Il inclue le carbone du sol.
\end{itemize}

Pour estimer les stocks de carbone dans la biomasse vivante, nous avons utilisé la méthode biométrique employée par \citet{Fayolle2018} dans le bassin du Congo qui est une approche basée sur les équations allométriques permettant d'estimer directement la biomasse totale ou partielle d'un arbre à partir de prédicteurs tels que le diamètre du tronc, la densité du bois, la hauteur totale et une constante environnementale. Ensuite nous avons multiplié la biomasse obtenue par un facteur d'expansion et de conversion de la biomasse 0,47 \citep{Zapfack2013}, pour obtenir la masse de carbone stockée. A l'échelle du paysage, la quantité de biomasse est estimée en multipliant la superficie de chaque biome par la quantité de biomasse moyenne par unité de surface pour ce biome. Les quantités moyennes par biome sont elles-mêmes estimées à partir de mesures faites à une échelle plus restreinte \citep{Mounmemi2021}.

\hypertarget{collecte-des-donnuxe9es-destimation-de-la-biomasse-du-bois-mort}{%
\subsection{Collecte des données d'estimation de la biomasse du bois mort}\label{collecte-des-donnuxe9es-destimation-de-la-biomasse-du-bois-mort}}

La collecte des données d'estimation des stocks de carbone dans les débris ligneux étendus au sol s'est faite à partir de la méthode d'intersection des lignes. Elle a consisté à considérer la somme des quatre côtés des parcelles de 2000 m² comme des lignes d'échantillonnages soit 180 m. Par la suite, les diamètres des deux extrémités ainsi que les longueurs des fragments couchés en intersection avec la ligne ont été mesurés lorsque : au moins 50\% du bois était en dehors du sol ; et la ligne d'échantillonnage traversant au moins 50\% du fragment \citep{Kabelong2018, Walker2012}. Pour déterminer la densité des fragments, on a heurté chaque fragment sélectionné avec un outil contondant (machette ou scie) puis on les a rangé dans l'une des trois classes de densité correspondante (dure-intermédiaire-pourri). En effet, si au moment de la percussion l'outil ne pénètre pas l'objet ou rebondi, le fragment est rangé dans la catégorie des dures. Par contre si l'outil pénètre partiellement l'objet et que des débris s'en détachent, le bois est rangé dans la catégorie des intermédiaires. Si enfin l'outil pénètre le bois et qu'il s'émiette ou des fragments importants s'en détachent, le bois est classé comme pourri \citep{Pearson2005}. La collecte des données pour l'estimation de la biomasse des bois morts sur pied s'est faite en mesurant les diamètres et les hauteurs de tous les arbres morts sur pied ainsi que les souches de diamètre supérieur à 5 cm et de longueur supérieure à 50 cm. Par la suite chaque bois mort ou débris a été rangé dans l'une des classes: décomposition peu avancée; décomposition moyenne; et décomposition très avancée selon les conditions de l'aubier \citep{Walker2012}.

\hypertarget{dispositif-de-collecte-des-herbacuxe9es-et-de-la-litiuxe8re}{%
\subsection{Dispositif de collecte des herbacées et de la litière}\label{dispositif-de-collecte-des-herbacuxe9es-et-de-la-litiuxe8re}}

La procédure d'échantillonnage s'est faite en délimitant des quadrats (05) de 1 m² (1 m x 1 m) à l'intérieur des placettes d'échantillonnage de 200 m². Ces quadrats ont ensuite été rasés de toute végétation herbeuse (Figure: \ref{fig:Echantillon-herbaceeslitiere2} a). Les tiges d'herbes et d'arbustes ont été débarrassées de leur feuillage et stockés séparément à l'intérieur des sachets plastiques. Ces échantillons ont été pesés sur place pour déterminer le Poids frais (Pf) en forêt (gramme/m²). Des échantillons composites de 500 grammes ont été constitués séparément pour les feuilles et pour les tiges à partir de la mixture de tous les différents groupes d'échantillons prélevés dans chaque quadrat. Au total, deux échantillons composites (feuilles et tiges) de 500 g chacun par placette d'échantillonnage de 200 m² seront constitués.

\begin{figure}

{\centering \includegraphics[width=0.8\linewidth]{images/herbacees_et_litiere} 

}

\caption{collecte d'échantillons d'herbacées et litière}\label{fig:Echantillon-herbaceeslitiere2}
\end{figure}

Les échantillons de litière grossière (toutes nécromasses ligneuses de diamètre inférieur à 5 cm et d'une longueur inférieure à 50 cm), les matières végétales non décomposées et les résidus des récoltes, et toutes les matières végétales non décomposées seront ramassées à l'intérieur des quadrats de 1 m² et stockées à l'intérieur des sachets plastiques (Figure: \ref{fig:Echantillon-herbaceeslitiere2}b). Ces échantillons ont également été pesés pour déterminer le Poids frais (Pf) en forêt (gramme/m²). Un échantillon composite de 500 grammes y était prélevé après la mixture des échantillons de litière provenant des cinq quadrats de 1 m² chacun. Ces échantillons composites ont été transportés à l'étuve pour être séchés pendant 72 heures à une température de 85°C pour les tiges et 65°C pour les feuilles et la litière. Les mesures du poids sec (Ps) ont suivi après séchage des échantillons à l'étuve.

\hypertarget{dispositif-de-collecte-des-uxe9chantillons-de-sol}{%
\subsection{Dispositif de collecte des échantillons de sol}\label{dispositif-de-collecte-des-uxe9chantillons-de-sol}}

La procédure d'échantillonnage s'est faite premièrement en délimitant un sous-quadrat de 0,25 m² (0,5 m × 0,5 m) dans tous les cinq (05) quadrats de 1 m² à l'intérieur des placettes d'échantillonnage de 200 m² (Figure: \ref{fig:Sous-quadratsol2}).

\begin{figure}

{\centering \includegraphics[width=0.8\linewidth]{images/echantillonnage_des_sols} 

}

\caption{Sous-quadrat de 0,25 m² pour l'échantillonnage des sols. (Adapté de Hairiah et al., (2011)}\label{fig:Sous-quadratsol2}
\end{figure}

Après la mise en place, les couches de litière organique du sol ont été préalablement raclées puis les échantillons de sol prélevés à des profondeurs de 0 à 10 cm, 10 à 20 cm et 20 à 30 cm. Les échantillons de sols en provenance de chaque quadrat ont été mixés entre eux à l'intérieur de petits seaux tout en respectant la profondeur de prélèvement (Figure: \ref{fig:Collecte-echantillonsol2} a). De la mixture obtenue, un échantillon composite de 500 grammes a été prélevé (Figure: \ref{fig:Collecte-echantillonsol2} b). Il en ressort 3 échantillons composites pour les 3 profils de sol par placette de 200 m², lesquels échantillons ont été envoyés au laboratoire d'analyse des sols de l'IITA (Institut Iternational d'Agriculture Tropical) pour la détermination du pourcentage de carbone organique par digestion chromique acide et analyse spectro-photométrique d'après la méthode décrite par Heanes (1984).

\begin{figure}

{\centering \includegraphics[width=0.8\linewidth]{images/profil_sol} 

}

\caption{Mise en place mini profil et collecte échantillon de sol}\label{fig:Collecte-echantillonsol2}
\end{figure}

\hypertarget{analyse-des-donnuxe9es}{%
\subsection{Analyse des données}\label{analyse-des-donnuxe9es}}

\hypertarget{caractuxe9risation-de-la-diversituxe9-floristique}{%
\subsubsection{Caractérisation de la diversité floristique}\label{caractuxe9risation-de-la-diversituxe9-floristique}}

La diversité floristique renvoie à la richesse (nombre d'espèces végétales) et l'équitabilité (égalité en terme de nombre d'individus pour les espèces considérées). L'analyse de la diversité floristique s'est faite à partir de la liste des espèces identifiées dans chaque parcelle. Ces informations ont été regroupées sous forme de matrice de données à partir d'un tableur Excel. Dans cette matrice les lignes correspondent aux espèces et les colonnes aux parcelles d'inventaires.Chaque cellule de la matrice renferme les valeurs d'abondance de chaque espèce dans une parcelle donnée. Nous avons ensuite utilisé le Package \emph{Biodiversity.R} version 2.14-2 qui est une application necessitant l'usage du logiciel R pour l'analyse des données \citep{Kindt_Coe2005} Cette application a la capacité de faire des analyses statistiques de la biodiversité ainsi que les communautés écologiques tout en incluant les courbes d'espèces cumulées, les indices de diversité, les profils Rényi, les analyses d'abondances, etc.

\hypertarget{composition-floristique-et-richesse-spuxe9cifique}{%
\paragraph{Composition floristique et richesse spécifique}\label{composition-floristique-et-richesse-spuxe9cifique}}

La composition floristique a été déterminée à partir de la liste des espèces ligneuses inventoriées dans chaque parcelle. Le nombre d'espèces, de genres et de familles ont été relevés. Les différences entre les parcelles en terme de composition d'espèces ont été investiguées par une valeur statistique issue du calcul des distances écologiques entre les sites. En effet pour des sites qui partagent un grand nombre d'espèces cette distance devrait être petite. Par contre lorsque les sites ont peu d'espèces en commun cette distance devrait être grande. \citet{Magurran2004} a relevé une difference entre la densité des espèces et la richesse spécifique numérique. La densité des espèces est le nombre d'espèces présents dans une unité d'échantillonnage ou surface donnée\citep{Hurlbert1971}. La richesse spécifique numérique est le nombre d'espèces par rapport à un effectif prédéfinit ou unité de biomasse\citep{Kempton1979}. D'après \citet{Borcard2018}, pour garantir la comparaison de deux parcelles d'échantillonnages, \citet{Sanders1968} a proposé une méthode de rarefaction qui estime le nombre d'espèces dans des unités d'échantillonnages renfermant le même nombre d'individus. Etant donné que la rarefaction ne peut être calculé qu'à partir des éffectifs réels(non transformés), la formule de \citet{Sanders1968} a été revue par \citet{Hurlbert1971}. Cette relation estime le nombre \emph{q'} des espèces dans une unité d'échantillonnage standard de \emph{n'} individus reposant sur une unité d'échantillonnage réelle de \emph{n} individus appartenant à \emph{q} espèces et n\_i individus appartenant à l'espèce \emph{i}. L'équation d'Hurlbert \eqref{eq:equation-Hurlbert2} et sa forme simplifiée \eqref{eq:equation-Hurlbertsimplifiee2}.

\begin{equation}
E(q')=\sum_{i=1}^q\left[1-\frac{\binom{n-n_i}{n}}{\binom{n}{n'}}\right]
\label{eq:equation-Hurlbert2}
\end{equation}

avec \(n'\le(n-n_1)\), \(n_1\) étant l'effectif de l'espèce la plus abondante.

\begin{equation}
\binom {n}{n'}=\frac{n!}{n'!(n-n')!};\binom{n-n_i}{n'}=\frac{(n-n_i)!}{n'!(n-n_i-n')!}
\label{eq:equation-Hurlbertsimplifiee2}
\end{equation}

\hypertarget{indices-de-diversituxe9}{%
\paragraph{Indices de diversité}\label{indices-de-diversituxe9}}

La détermination des indices de diversité permet d'obtenir des informations sommaires sur la richesse et l'équitabilité par le calcul d'une valeur statistique unique. Selon \citet{Gosselin_Laroussinie2004} de nombreux indices de diversité existent et leur détermination repose sur des calculs assez complexes. Ces indices sont fréquemment usités en écologie pour la caractérisation des peuplements\citep{Ramade1994}. Cependant l'usage d'un indice de diversité unique n'est pas toujours suffisant pour fournir des informations sur le classement des sites selon leur importance en diversité.Dans le cadre de cet travail, en dehors de certaines techniques de classement des diversités, les indices selectionnés par \citet{Picard_Gourlet-Fleury2008} et \citet{Mounmemi2021} ont été utilisés pour analyser la diversité floristique dans les sites d'inventaires. Il représente la somme des fréquences des differentes espèces identifiées dans les parcelles d'échatillonnages.

\hypertarget{indice-de-diversituxe9-de-shannon}{%
\subparagraph{Indice de diversité de Shannon}\label{indice-de-diversituxe9-de-shannon}}

L'indice de diversité de Shannon est une valeur statistique indépendante de la taille de l'échantillon analysé, par conséquent il est indiqué pour les études comparatives de peuplements forestiers. Tout en permettant de quantifier l'hétérogénéité de la biodiversité d'un milieu, l'indice de Shannon permet de mesurer la composition en espèces en combinant l'abondance relative des espèces à la richesse spécifique. Pour les échantillons de faible taille, l'indice de Shannon varit entre 1 et 4,5 \citep{Riera_Alexandre2004}. Pour des parcelles ayant le même nombre d'espèces, la plus diversifiée sera celle dont les effectifs de chaque catégorie d'espèces seront à peu près les mêmes. la formule mathématique de l'indice de Shannon \citep{Shannon1948} est donnée par la relation \eqref{eq:Indice-deShannon2}.

\begin{equation}
H=-\sum_{i=1}^q\frac{n_i}{N}ln\frac{n_i}{N}
\label{eq:Indice-deShannon2}
\end{equation}

\(H\) : indice de Shannon;
\(q\) : nombre d'espèces;
\(n_i\) : nombre d'individus d'une espèce;
\(N\) : nombre total des individus enregistrés.

Quel que soit le nombre d'individus enregistré dans une unité d'échantillonnage, l'indice de Shannon (\(H\)) atteint son maximum (\(H_max\)) lorsque toutes les espèces sont représentées par le même nombre d'effectif \eqref{eq:max-Shannon2}.

\begin{equation}
H_{max}=-\sum_{i=1}^q\frac{1}{q}ln\frac{1}{q}=lnq
\label{eq:max-Shannon2}
\end{equation}

par conséquent l'équitabilité de Pielou\citep{Pielou1975} peut être définit par la relation \eqref{eq:Equitabilite-dePielou2}:

\begin{equation}
J=\frac{H}{H_{max}}
\label{eq:Equitabilite-dePielou2}
\end{equation}

\(J\) : indice de Pielou

Malgrè le fait qu'il soit fortement dépendant de la richesse spécifique, l'équitabilité de Pielou reste très usité comme indice d'équitabilité. D'autre part l'équitabilité peut être relié à la forme de la courbe d'abondance des espèces, qui est une fonction décrivant l'allure rang/abondance où l'abscisse range les espèces en ordre décroissant d'abondance et l'ordonnée représente le logarithme transformé d'abondance. Les principaux modèles de courbes observés varient des séries géométriques où peu espèces sont dominantes et d'autres quasi rares, vers les modèles de Broken Stick où les espèces ont les mêmes abondances. Lorsque \(J\) tend vers 0, un phénomène de dominance est observé, par contre lorqu'il tend vers 1 on relève une répartition régulière des éffectifs entre les espèces.

\hypertarget{indice-de-diversituxe9-de-simpson}{%
\subparagraph{Indice de diversité de Simpson}\label{indice-de-diversituxe9-de-simpson}}

Il permet de mesurer la probabilité que deux individus sélectionnés de manière aléatoire soient de la même espèce. Sa formule est indiquée comme suit \eqref{eq:Indice-deSimpson}:

\begin{equation}
𝛌=\sum_{i=1}^q\ p_i^2 
\label{eq:Indice-deSimpson}
\end{equation}

\(𝛌\) : indice de Simpson;
\(p_i\) : fréquence relative de l'espèce i;
\(q\) : nombre des espèces.

La valeur de 𝛌 s'accroit lorsque la probabilité que deux individus soient de la même espèce est élevée. Généralement la diversité de Simpson est transformée sous la forme: \(D=1-𝛌\), \(D\) variant de 0 à 1. Plus \(D\) augmente, plus le milieu est hétérogène et la diversité grande. La richesse spécifique, l'entropie de Shannon et la diversité de Simpson ont été regroupées dans la formule de l'entropie généralisée de Rényi \citep{Renyi1961} et reprise par Hill(1973) et Pielou (1975):

\begin{equation}
H_a=\frac{1}{1-a}log\sum_{i=1}^qp_i^a
\label{eq:Equation-deRenyi}
\end{equation}

avec \(a\) étant la mesure de l'ordre de l'entropie(a=0,1,2,\ldots). Hill(1973) a également proposé d'utilisé les \emph{nombres de diversité} \eqref{eq:Nombre-dediversite2}.

\begin{equation}
N_a=e^{H_a}
\label{eq:Nombre-dediversite2}
\end{equation}

le paramètre \(a\) permet de quantifier l'importance des effectifs et par conséquent l'équitabilité: lorsque \(a=0\) la diversité se resume au nombre d'espèces (présence-absence); lorsque \(a\) augmente de plus en plus l'importance est accordée à l'espèce ayant l'effectif le plus élevé. De même avec cette notation, Hill a proposé l'usage du ratio \(E_1={N_1}/{N_0}\) pour l'équitabilité de Shannon et \(E_2={N_2}/{N_0}\) pour l'équitabilité de Simpson. De nos jours la plupart des écologiste préfèrent utilisés les nombres de Hill pour étudier la diversité taxonomique parce que ces nombres parfois appelés \enquote{nombres équivalents} sont faciles à interpreter et représentent le nombre d'espèces équiprobables nécessaires pour produire la même valeur de diversité que la distribution observé. Les nombres de diversité sont également intéressant pour l'interprétation des modèles lineaires parce qu'ils sont plus probables d'être en rélation de colinéarité avec les variables environnementales.

\hypertarget{courbes-cumuluxe9es-des-espuxe8ces}{%
\subparagraph{Courbes cumulées des espèces}\label{courbes-cumuluxe9es-des-espuxe8ces}}

les courbes cumulées des espèces permettent de suivre l'évolution de la richesse spécifique sur des combinaisons de plusieurs sites. Ces courbes représentent la moyenne cumulée de la richesse spécifique lorsque plusieurs ou la totalité des sites étudiés sont mis ensemble \citep{Kindt_Coe2005}. Le calcul de la moyenne cumulée de la richesse spécifique s'explique par le fait que differentes combinaison de sites donnent des richesses spécifiques distinctes. Par contre lorsqu'on considère la moyenne de la richesse spécifique issue d'un éventail de combinaison de sites allant de trois ou plusieurs sites, jusqu'atteindre toutes les tailles de combinaison possible, on est presque certain d'avoir une tendance fidèle de l'évolution de la richesse spécifique dans l'ensemble des sites étudiés. Ces courbes permettent également de comparer la richesse spécifique entre les differentes catégories d'utlisation des terres.

\hypertarget{arbre-de-regression-multivariuxe9}{%
\subparagraph{Arbre de Regression Multivarié}\label{arbre-de-regression-multivariuxe9}}

C'est une méthode d'analyse proposée par \citet{Death2002} avec pour but de ressortir les correlations entre la distribution des espèces et les types d'utlisation de terre. Elle consiste en une technique statistique qui décrit et prédit les préferences écologiques des espèces par rapport à des clusters prédéfinis. Le résultat de l'analyse est un modèle d'arbre dont les \enquote{feuilles} (regroupement de parcelles) sont des ensembles formés par des découpages successifs des données, chaque découpage étant définit par une valeur seuil ou variable explicative basée sur des paramètres environnementaux. Concrètement l'ARM scinde les données en groupes similaires du point de vue de la composition spécifiques basées sur des valeurs environnementales seuils. Pour cela deux procedures sont exécutées simultanément:

\begin{itemize}
\item
  le calcul des contraintes de partitionnement des données;
\item
  et le calcul de l'erreur relative liée aux niveaux de partitionnement successif par validations croisées multiples.
\end{itemize}

Compte tenu de la multitude de potentiels solutions que présente l'algorithme, seules les prédictions les plus précises sont selectionnées. Dans un modèle d'arbre presentant une bonne prédiction, les objets sont assignés à des groupes appropriés. La mesure de l'erreur de prédiction est donnée par le CVRE (Cross Validation Relative Error). Si l'on ignore les validations croisées, le meilleur modèle serait l'arbre dont le nombre de partition minimise la variance non expliquée par l'arbre (l'erreur relative) ou RE (Relative Error). C'est la solution qui maximise le R².

\hypertarget{indices-de-valeur-des-espuxe8ces-ou-espuxe8ces-caractuxe9ristiques}{%
\subparagraph{Indices de valeur des espèces ou espèces caractéristiques}\label{indices-de-valeur-des-espuxe8ces-ou-espuxe8ces-caractuxe9ristiques}}

Pour mesurer l'association entre plusieurs espèces appartenant à un même groupe de parcelles, on a utilisé l'indicateur de valeur \citep{Dufrene_Legendre1997, DeCaceres_Legendre2009}. C'est une méthode qui calcule un indice (IndVal) entre les espèces et les groupes de parcelles auxquelles elles appartiennent et par la suite recherhe le groupe correspondant au test de permutation. Cet indicateur de valeur est un indice qui est le produit de deux composantes \enquote*{A} et \enquote*{B}. la composante \enquote*{A} est l'abondance moyenne de l'espèce dans le groupe cible comparé à son abondance moyenne à travers les autres groupes (spécificité). La composante \enquote*{B} est la proportion de parcelles des groupes où l'espèce est présente (fidélité).

\hypertarget{estimation-de-la-biomasse}{%
\subsubsection{Estimation de la biomasse}\label{estimation-de-la-biomasse}}

\hypertarget{biomasse-ligneuse}{%
\paragraph{Biomasse ligneuse}\label{biomasse-ligneuse}}

Elle correspond à la masse de matière végétale ligneuse sèche par unité de surface. La biomasse totale en bois des ligneux sur pieds répartie en biomasse aérienne (AGB) et en biomasse souterraine (BGB). L'estimation de la biomasse aérienne des arbres s'est faite à partir des équations allométriques élaborée par \citet{Fayolle2018} dans le bassin du Congo. Elle a été utilisée pour les arbres de DHP supérieur à 10 cm. C'est une équation qui prend en compte le diamètre de l'arbre, la densité du bois et l'indice climatique comme prédicteur de la biomasse épigée \eqref{eq:equation-deFayolle2}. La formule mathématique est la suivante :

\begin{equation}
BA=℮^{\left(a+b⨯ln(\rho)+c⨯ln(D)+d⨯ln(D)+e⨯E\right)}
\label{eq:equation-deFayolle2}
\end{equation}

\(BA\) : Biomasse épigée en kg;
\(D\) : Diamètre(cm);
\(ρ\) : densité spécifique d'une espèce exprimée en g/cm\textsuperscript{3} ;
\(E\) : indice climatique ;
\(a, b, c\ et\ d\) sont des estimations du coefficient (\(a\) =-2,3027; \(b\) =1,1599; \(c\) =3,0484; \(d\) =-0,0807) et \(e\) =0,3197.

La biomasse des individus de \(5 cm < dhp < 10 cm\) a eté estimée à partir de l'équation allométrique de \citet{Chave2014} \eqref{eq:equation-Chave2014} dont la relation est la suivante :

\begin{equation}
BA=℮^{\left(-1,803-0,976E+0,976ln(\rho)+2,673ln(D)-0,0299(ln(D))^2\right)}
\label{eq:equation-Chave2014}
\end{equation}

\hypertarget{biomasse-souterrainebgb-des-ligneux-sur-pied}{%
\paragraph{Biomasse souterraine(BGB) des ligneux sur pied}\label{biomasse-souterrainebgb-des-ligneux-sur-pied}}

L'équivalent racinaire de la biomasse des ligneux sur pied \eqref{eq:equival-racinaire2} a été trouvé en multipliant la valeur de la biomasse aérienne(AGB) par un coefficient R dont la valeur est estimée en fonction du niveau de pluviométrie de la zone concernée\citep{GIEC2007}. Pour notre zone d'étude, nous avons considéré \(R=1/4\).

\begin{equation}
BGB=AGB⨯R 
\label{eq:equival-racinaire2}
\end{equation}

\(BGB\): Below Ground Biomass ou Biomasse souterraine (kg);
\(ABG\): Above ground Biomass ou biomasse aérienne (kg);
\(R\): Root to shoot ratio ou coefficient racinaire.

Les stocks de carbone ont été déduits de la biomasse aérienne en la multipliant par la constante 0,475 \citep{Zapfack2013}. La Biomasse Totale (\(TB\)) des ligneux sur pied \eqref{eq:biomasse-totale2} a été alors estimée en sommant les deux valeurs ci-dessus :

\begin{equation}
TB=AGB+BGB
\label{eq:biomasse-totale2}
\end{equation}

\(TB\) : Total Biomass ou Biomasse totale (kg) ;
\(BGB\) : Below Ground Biomass ou Biomasse Souterraine (kg) ;
\(AGB\) : Above Ground Biomass ou Biomasse Aérienne (kg).

Les valeurs de densité du bois des arbres échantillonnés ont été puisées dans la littérature. Pour les espèces dont les valeurs de densité n'ont pas été retrouvées dans la littérature,une valeur par défaut (\(ρ\ defaut = 0,58 g/cm^3\)) pour les forêts tropicales d'Afrique, a été utilisée \citep{Mounmemi2021}.

\hypertarget{biomasse-des-espuxe8ces-arborescentes-sur-pied}{%
\paragraph{Biomasse des espèces arborescentes sur pied}\label{biomasse-des-espuxe8ces-arborescentes-sur-pied}}

Le Tableau: \ref{tab:Biomasse-especesarborescentes2} présente les équations allométriques utilisées pour estimer la biomasse aérienne des palmiers et bananiers.

\begin{table}
\centering
\caption{\label{tab:Biomasse-especesarborescentes2}Equations allométriques de quelques espèces végétales arborescentes.}
\centering
\resizebox{\ifdim\width>\linewidth\linewidth\else\width\fi}{!}{
\begin{threeparttable}
\begin{tabular}[t]{lll}
\toprule
Espèces & Equations & Sources\\
\midrule
Bananiers (Musaceae) & $ABG=0,0303⨯DHP^{2,1345}$ & Arifin (2001)\\
Palmiers (Arecaceae) & $AGB=23,487+41,851⨯(lnH^2)$ & Pearson et Brown (2005)\\
\bottomrule
\end{tabular}
\begin{tablenotes}
\item \textit{Note: } 
\item AGB = Above Ground Biomass (kg) ; DHP = Diamètre à Hauteur de Poitrine (cm) ; H = hauteur du tronc de palmier de sa souche à la base de la partie non visible du tronc (cm)
\end{tablenotes}
\end{threeparttable}}
\end{table}

\hypertarget{biomasse-des-herbacuxe9es-et-de-la-litiuxe8re}{%
\paragraph{Biomasse des herbacées et de la litière}\label{biomasse-des-herbacuxe9es-et-de-la-litiuxe8re}}

Les mesures du poids sec (\(Ps\)) ont suivi après séchage des échantillons à l'étuve. La biomasse a été par la suite déterminée à l'aide de la relation développée par Pearson et Brown \eqref{eq:pearson-Brown2} en 2005:

\begin{equation}
Biomasse=\frac{Pht⨯Ps}{Ph}
\label{eq:pearson-Brown2}
\end{equation}

\(Ps\) = Poids sec (g) de l'échantillon après 3 jours à l'étuve à 65 °C pour les feuilles et la litière, 85°C pour les tiges ;
\(Ph\) = Poids humide (g) de l'échantillon composite soit 500 g;
\(Pht\) = Poids humide total en forêt (somme des poids frais issus des 5 quadrats)

\hypertarget{biomasse-des-arbres-morts}{%
\paragraph{Biomasse des arbres morts}\label{biomasse-des-arbres-morts}}

La biomasse des arbres morts a été estimée selon les formules ci-contre: \eqref{eq:arbres-morts2}

\begin{equation}
\begin{split}
Biomasse=\frac{\pi}{4}⨯D^2⨯L⨯\rho\\
Diamètre\  moyen\ D=\frac{D_1+D_2}{2}
\end{split}
\label{eq:arbres-morts2}
\end{equation}

\(D\) : diamètre moyen de l'arbre mort ( cm);
\(L\) : longueur de l'arbre dans la placette d'échantillonnage (cm);
\(\rho\) : densité spécifique (g/cm\textsuperscript{3}) lié au niveau de décomposition de l'arbre, 0,8: décomposition peu avancée; 0,5: décomposition moyennement avancée; 0,3: décomposition très avancée.

\hypertarget{cartographie-de-la-distribution-spatiale-des-stocks-de-carbone}{%
\subsubsection{Cartographie de la distribution spatiale des stocks de carbone}\label{cartographie-de-la-distribution-spatiale-des-stocks-de-carbone}}

Cette opération s'est appuyé sur la classification des images multibandes correspondantes à nos différentes zones d'intérêts. Le suivi du protocole de classification d'image nous a exigé l'installation des utilitaires et les \emph{plugins} necessaires, sur notre poste de travail notamment Qgis et les Semi automatic Classification Plugins (SCP). Les étapes de ce protocole ont été les suivantes (Figure: \ref{fig:Classification-images2}):

\begin{itemize}
\tightlist
\item
  l'enregistrement des coordonnées géographiques de la zone d'étude;
\item
  le téléchargement des scènes Landsat correspondantes;
\item
  le prétraitement des images (sélection des images, corrections radiométriques et corrections géométriques);
\item
  l'analyse des images par des approches et méthodes appropriées ;
\item
  faire la classification proprement dite des images à partir des algorithmes;
\item
  faire le post-traitement des images à partir des données collectées sur le terrain et valider la carte finale.
\end{itemize}

\begin{figure}

{\centering \includegraphics[width=0.8\linewidth]{images/protocole_classification_image} 

}

\caption{Protocole de classification d'image}\label{fig:Classification-images2}
\end{figure}

La classification est un procédé qui attribue à chaque pixel issu de plusieurs bandes d'images une certaine classe ou thème selon les caractéristiques statistiques (valeur Minimale, valeur Maximale, valeur Moyenne et déviation standard) de la valeur de l'intensité du pixel. Elle vise à créer une corrélation entre les données recueillies sur le terrain (classes d'informations) et le traitement numérique des images à une échelle plus grande (classes spectrales). Il existe plusieurs approches de classification d'image : l'approche basée sur les pixels (chaque pixel de l'image est analysé de manière individuelle selon l'information spectrale qu'il contient et rattaché à la classe correspondante), elle est utile pour analyser les multiples changements dans un espace pendant une période réduite ; l'approche par objet (les pixels mitoyens ayant les mêmes propriétés ou presque (spectrale, spatiale et/ou temporelle) sont regroupés en régions appelées objets ou segments), elle est très usitée pour les images de grande résolution. Deux principales méthodes de classification sont généralement utilisées : la classification supervisée et la classification non supervisée. La classification supervisée s'appuie sur les connaissances de l'analyste, sa familiarité avec les régions géographiques et les types de surfaces présentes dans l'image, pour identifier des échantillons assez homogènes de l'image, représentatifs des types d'occupation de sol (classes d'informations) (Figure: \ref{fig:Classification-supervisee2} ). Ces échantillons constituent un ensemble de données test ou zones de calibrage. Les informations numériques issues des bandes et de chaque pixel de ces zones, sont utilisées par l'ordinateur à travers des algorithmes pour reconnaitre des régions aux caractéristiques similaires sur toute l'étendue de l'image. Les algorithmes utilisés pour ces opérations sont principalement \citep{Anonyme2019e}:

\begin{itemize}
\tightlist
\item
  Maximum de Vraisemblance : algorithme de classification qui développe une règle de décision basée sur la probabilité d'appartenance des valeurs spectrales à une classe donnée;
\item
  Artificial Neural Network : algorithme qui construit des liens de manière autonome entre les inputs et les outputs à partir des données d'entrainements;
\item
  Support vector Machines : algorithme de classification qui crée un hyperplan séparant de manière optimale N-classes dans un espace à N-dimension;
\item
  Arbre de décision : algorithme reposant sur une succession de tests binaires débouchant à la construction d'un arbre de décision visant à créer des régions de plus en plus homogènes;
\item
  Random Forest : amélioration de l'arbre de décisions qui combine différentes classifications de la même donnée pour générer une classification encore plus précise.
\end{itemize}

\begin{figure}

{\centering \includegraphics[width=0.8\linewidth]{images/etapes_classification_supervisee} 

}

\caption{Principales étapes d'une classification supervisée}\label{fig:Classification-supervisee2}
\end{figure}

La classification non supervisée s'appuie uniquement sur l'information numérique des données pour regrouper les pixels en classes d'éléments identiques (classes spectrales) (Figure: \ref{fig:Classification-nonsupervisee2}) . Par la suite, l'analyste, lorsque cela s'avère satisfaisant, peut associer ces classes à des catégories d'informations. Des algorithmes sont utilisés pour isoler des groupes statistiques naturels ou structures de données. Généralement, deux algorithmes sont utilisés :

\begin{itemize}
\tightlist
\item
  K-Means : algorithme qui, à partir des classes localisées aléatoirement, exécute une répétition de regroupement de pixels selon la moyenne la plus proche jusqu'à atteindre une constance entre les itérations;
\item
  et ISODATA : affinement du K-Means qui permet de réorganiser la taille des classes en s'appuyant sur un seuil prédéfini.
\end{itemize}

\begin{figure}

{\centering \includegraphics[width=0.8\linewidth]{images/etapes_classification_non_supervisee} 

}

\caption{Principales étapes d'une classifcation non supervisée}\label{fig:Classification-nonsupervisee2}
\end{figure}

\hypertarget{ruxe9sultats-et-discussions}{%
\chapter{Résultats et Discussions}\label{ruxe9sultats-et-discussions}}

\hypertarget{ruxe9sultats}{%
\section{Résultats}\label{ruxe9sultats}}

\hypertarget{caractuxe9risation-de-la-diversituxe9-floristique-des-foruxeats-communales-de-ndikinimuxe9ki-et-nanga-eboko}{%
\subsection{Caractérisation de la diversité floristique des forêts communales de Ndikiniméki et Nanga-Eboko}\label{caractuxe9risation-de-la-diversituxe9-floristique-des-foruxeats-communales-de-ndikinimuxe9ki-et-nanga-eboko}}

Dans la Forêt Communale de Ndikiniméki notre étude floristique a couvert un ensemble de 32 parcelles d'inventaires (Tableau:\ref{tab:Nomenclatureparcelles-Ndikinimeki3}) réparties sur 05 catégories d'utilisation des terres:

\begin{verbatim}
#> Description of structure(list(x = c("FM", "CAC", "CAC", "FS", "CAC", "FS", "FS",  "FM", "CAC", "CAC", "FM", "FM", "FS", "FM", "JA", "FS", "FMa",  "JA", "CAC", "FS", "FS", "CAC", "FM", "FM", "FM", "CAC", "FMa",  "FM", "FM", "FM", "FM", "FM")), row.names = c(NA, -32L), class = "data.frame")
#> 
#>  Factor 
#>          
#> x            FM CAC    FS  FMa   JA
#>   Count   13.00   8  7.00 2.00 2.00
#>   Percent 40.62  25 21.88 6.25 6.25
#> Mode FM
\end{verbatim}

\begin{itemize}
\tightlist
\item
  13 parcelles de Forêts Matures (FM) représentant 40,62\% des inventaires;
\item
  08 parcelles de Cacaoyères (CAC) comptant pour 25\% des inventaires;
\item
  07 parcelles de Forêts Secondaires (FS) représentant 21,88\% des inventaires;
\item
  02 parcelles de Forêts Marécageuses (FMa) comptant pour 6,25\% des inventaires;
\item
  et 02 parcelles de Jachères (JA) représentant 6,25\% des inventaires.
\end{itemize}

\begin{longtable}[]{@{}l@{}}
\caption{\label{tab:Nomenclatureparcelles-Ndikinimeki3}Nomenclature des parcelles dans la FC de Ndikiniméki. (x: Numéro des parcelles)}\tabularnewline
\toprule()
x \\
\midrule()
\endfirsthead
\toprule()
x \\
\midrule()
\endhead
1 \\
2 \\
3 \\
4 \\
5 \\
6 \\
7 \\
8 \\
9 \\
10 \\
11 \\
12 \\
13 \\
14 \\
15 \\
16 \\
17 \\
18 \\
19 \\
20 \\
21 \\
22 \\
23 \\
24 \\
25 \\
26 \\
27 \\
28 \\
29 \\
30 \\
31 \\
32 \\
\bottomrule()
\end{longtable}

Dans la FC de Nanga-Eboko l'étude floristique a couvert 23 parcelles (Tableau:\ref{tab:NomenclatureparcellesNanga-Eboko3}) réparties en 06 catégories d'utlisation des terres:

\begin{verbatim}
#> Description of structure(list(x = c("FSA", "FSA", "CUL", "SAV", "SAV", "FGA",  "FGA", "FSA", "MAR", "MAR", "FSA", "MAR", "FSA", "FGA", "FSJ",  "SAV", "CUL", "FSA", "FSJ", "FSA", "FSJ", "MAR", "FSJ")), row.names = c(NA,  -23L), class = "data.frame")
#> 
#>  Factor 
#>          
#> x           FSA   FSJ   MAR   FGA   SAV CUL
#>   Count    7.00  4.00  4.00  3.00  3.00 2.0
#>   Percent 30.43 17.39 17.39 13.04 13.04 8.7
#> Mode FSA
\end{verbatim}

\begin{itemize}
\tightlist
\item
  07 parcelles de Forêts Secondaires Adultes (FSA);
\item
  04 parcelles de Forêts Secondaires Jeunes (FSJ);
\item
  04 parcelles de Marécages (MAR);
\item
  03 parcelles de Forêts Galeries (FGA);
\item
  03 parcelles de Savanes (SAV);
\item
  et 02 parcelles de cultures (CUL).
\end{itemize}

\begin{longtable}[]{@{}l@{}}
\caption{\label{tab:NomenclatureparcellesNanga-Eboko3}Nomenclature des parcelles dans la FC de Nanga-Eboko. (x: Numéro des parcelles)}\tabularnewline
\toprule()
x \\
\midrule()
\endfirsthead
\toprule()
x \\
\midrule()
\endhead
1 \\
2 \\
3 \\
4 \\
5 \\
6 \\
7 \\
8 \\
9 \\
10 \\
11 \\
12 \\
13 \\
14 \\
15 \\
16 \\
17 \\
18 \\
19 \\
20 \\
21 \\
22 \\
23 \\
\bottomrule()
\end{longtable}

Dans la FC de Ndikiniméki l'on a identifié 148 espèces (Tableau:\ref{tab:Especes-Ndikinimeki3}) appartenant à 43 familles (Tableau:\ref{tab:Familles-Ndikinimeki3}).

\begin{longtable}[]{@{}l@{}}
\caption{\label{tab:Especes-Ndikinimeki3}Identification des espèces enregistrées dans la FC de Ndikiniméki. (x: Nom des espèces)}\tabularnewline
\toprule()
x \\
\midrule()
\endfirsthead
\toprule()
x \\
\midrule()
\endhead
Afzelia bipindensis \\
Albizia glaberrima \\
Albizia zygia \\
Allophyllus africana \\
Alstonia boonei \\
Aningeria altissima \\
Annona sp. \\
Anonidium mannii \\
Anthonotha macrophylla \\
Antiaris welwitchii \\
Antidesma venosum \\
Antrocaryon klaineanum \\
Baillonnella toxisperma \\
Baphia lepidobotrys \\
Beilschmiedia obscura \\
Blighia welwitschii \\
Bosqueia angolensis \\
Caloncoba welwitschii \\
Canarium schweinfurthii \\
Canthium arnoldianum \\
Carapa procera \\
Carica papaya \\
Ceiba pentandra \\
Celtis adolfi friderici \\
Celtis zenkeiri \\
Citrus reticulata \\
Coelocaryon preussi \\
Coffea sp. \\
Cola acuminata \\
Cola argentea \\
Cola sp. \\
Cordia sp. \\
Corynanthe pachyceras \\
Cyrtogonone argentea \\
Cylicodiscus gabonensis \\
Dacryodes edulis \\
Dacryodes igaganga \\
Dacryodes macrophylla \\
Desbordesia glaucescens \\
Detarium macrocarpum \\
Dialium bipendensis \\
Diospyros crassiflora \\
Diospyros sanza-minika \\
Discoglypremna caloneura \\
Distemonanthus benthamianus \\
Drypetes preussii \\
Drypetes sp. \\
Duboscia macrocarpa \\
Elaeis guineensis \\
Enantia chlorantha \\
Entandrophragma cylindricum \\
Eribroma oblongum \\
Erythrina excelsa \\
Fagara heitzii \\
Ficus exasperata \\
Ficus mucuso \\
Funtumia elastica \\
Gambeya africana \\
Garcinia mannii \\
Grewia coriacia \\
Holoptelea grandis \\
Homalium letestui \\
Hylodendron gabonense \\
Hymenocardia heudelotii \\
Irvingia gabonensis \\
Isolona hexaloba \\
Keayodendron bridelioides \\
Khaya ivorensis \\
Kigelia acutifolia \\
Klainedoxa gabonensis \\
Lannea welwitschii \\
Lasiodiscus marmoratus \\
Lophira alata \\
Lovoa trichilioides \\
Macaranga burifolia \\
Maesopsis eminii \\
Magnistipula zenkeri \\
Malacantha alnifolia \\
Mallotus oppositifolius \\
Mangifera indica \\
Mansonia altissima \\
Mareyopsis longifolia \\
Markhamia tomentosa \\
Massularia acuminata \\
Meiocarpidium lepidotum \\
Milicia excelsa \\
Monodora tenuifolia \\
Musa sp. \\
Musanga cecropioides \\
Myrianthus arboreus \\
Nauclea diderrichii \\
Nauclea pobeguinii \\
Neosloetiopsis sp. \\
Nesogordonia papaverifera \\
Pachypodanthium staudtii \\
Panda oleosa \\
Paraberlinia bifoliolata \\
Parinari glaba \\
Parkia bicolor \\
Pentaclethra macrophylla \\
Persea americana \\
Petersianthus macrocarpus \\
Phyllanthus discoideus \\
Picralima nitida \\
Piptadeniastrum africanum \\
Plagiostyles africana \\
Polyalthia suaveolens \\
Porterandia cladantha \\
Pseudospondias microcarpa \\
Psicotrea \\
Pteleopsis hylodendron \\
Pterocarpus mildbraedii \\
Pterocarpus soyauxii \\
Pterygota macrocarpa \\
Pycnanthus angolensis \\
Raphia sp. \\
Rauvolfia vomitoria \\
Ricinodendron heudelotii \\
Rinorea sp. \\
Samanea dinklagei \\
Santiria trimera \\
Sapium sp. \\
Sorindeia grandifolia \\
Staudtia kamerunensis \\
Sterculia rhinopetala \\
Sterculia tragacantha \\
Strombosia grandifolia \\
Strombosia pustulata \\
Strombosia scheffleri \\
Strombosiopsis tetrandra \\
Symphonia globulifera \\
Syzygium rowlandii \\
Tabernaemontana crassa \\
Tabernaemontana pachysiphon \\
Tetrorchidium didymostemon \\
Terminalia ivorensis \\
Terminalia superba \\
Tetrapleura tetraptera \\
Theobroma cacao \\
Treculia africana \\
Trichilia dregeana \\
Trichilia welwitschii \\
Triplochyton scleroxylon \\
Turreaenthus africanus \\
Uapaca guineensis \\
Vernonia conferta \\
Vitex grandifolia \\
Xylopia aethiopica \\
\bottomrule()
\end{longtable}

\begin{longtable}[]{@{}l@{}}
\caption{\label{tab:Familles-Ndikinimeki3}Identification des familles dans la FC de Ndikiniméki. (x: Nom des familles d'espèces)}\tabularnewline
\toprule()
x \\
\midrule()
\endfirsthead
\toprule()
x \\
\midrule()
\endhead
Anacardiaceae \\
Achariaceae \\
Annonaceae \\
Apocynaceae \\
Arecaceae \\
Asteraceae \\
Bignoniaceae \\
Boraginaceae \\
Burseraceae \\
Cannabaceae \\
Caricaceae \\
Cecropiaceae \\
Chrysobalanaceae \\
Combretaceae \\
Clusiaceae \\
Ebenaceae \\
Euphorbiaceae \\
Fabaceae \\
Flacourtiaceae \\
Irvingiaceae \\
Lauraceae \\
Lecythidaceae \\
Malvaceae \\
Meliaceae \\
Moraceae \\
Musaceae \\
Myristicaceae \\
Myrtaceae \\
Ochnaceae \\
Olacaceae \\
Pandaceae \\
Phyllanthaceae \\
Rhamnaceae \\
Rubiaceae \\
Rutaceae \\
Sapindaceae \\
Sapotaceae \\
Simaroubaceae \\
Sterculiaceae \\
Tiliaceae \\
Ulmaceae \\
Verbenaceae \\
Violaceae \\
\bottomrule()
\end{longtable}

Dans la FC de Nanga-Eboko 122 espèces (Tableau:\ref{tab:EspecesNanga-Eboko3}) rattachées à 39 familles ont été enregistrées (Tableau:\ref{tab:FamillesNanga-Eboko3}).

\begin{longtable}[]{@{}l@{}}
\caption{\label{tab:EspecesNanga-Eboko3}Identification des espèces enregistrées dans la FC de Nanga-Eboko. (x: Nom des espèces)}\tabularnewline
\toprule()
x \\
\midrule()
\endfirsthead
\toprule()
x \\
\midrule()
\endhead
Afzelia pachyloba \\
Albizia ferruginea \\
Albizia glaberrima \\
Albizia gumifera \\
Albizia zygia \\
Annona senegalensis \\
Annona sp. \\
Anopyxis klaineana \\
Anthonotha macrophylla \\
Anthocleista schweinfurthii \\
Antiaris welwitschii \\
Antrocaryon klaineanum \\
Baphia pubescens \\
Beilschmiedia obscura \\
Beilschmiedia sp. \\
Berlinia grandiflora \\
Blighia welwitschii \\
Bosqueia angolensis \\
Bridelia ferruginea \\
Caloncoba gilgiana \\
Caloncoba glauca \\
Canarium schweinfurthii \\
Canthium sp. \\
Ceiba pentandra \\
Celtis adolfi friderici \\
Celtis zenkeiri \\
Cleistanthus polystachyus \\
Coffea sp. \\
Cola argentea \\
Cola lateritia \\
cola sp. \\
Cordia platythyrsa \\
Crossopterix febrifuga \\
Dacryodes macrophylla \\
Desbordesia glaucescens \\
Dialium bipendense \\
Dialium zenkeri \\
Discoglypremna caloneura \\
Drypetes sp. \\
Duboscia macrocarpa \\
Duguetia staudtii \\
Elaeis guineensis \\
Elasophorbia drupifera \\
Enantia chlorantha \\
Eribroma oblonga \\
Erythropleum ivorense \\
Fagara heitzii \\
Fernandoa adolfi friderici \\
Ficus exasperata \\
Ficus mucuso \\
Ficus vogelli \\
Funtumia africana \\
Funtumia elastica \\
Gambeya africana \\
Gambeya beguei \\
Garcinia mannii \\
Hannoa klaineana \\
Holoptelea grandis \\
Hylodendron gabunense \\
Inconnu \\
Irvingia grandifolia \\
Kaela \\
Lannea welwitschii \\
Lophira alata \\
Macaranga burifolia \\
Maesopsis eminii \\
Maesobotrya sp. \\
Malacantha alnifolia \\
Mammea africana \\
Markhamia tomentosa \\
Meiocarpidium lepidotum \\
Milicia excelsa \\
Monodora brevipes \\
Monodora tenuifolia \\
Musanga cecropioides \\
Myrianthus arboreus \\
Nauclea pobeguinii \\
Ongokea gore \\
Ouratea sp. \\
Parkia bicolor \\
Pentaclethra macrophylla \\
Persea americana \\
Petersianthus macrocarpus \\
Phyllanthus discoideus \\
Piptadeniastrum africanum \\
Polyalthia suaveolens \\
Pseudospondias microcarpa \\
Psychotria sp. \\
Pterocarpus soyauxii \\
Pycnanthus angolensis \\
Raphia sp \\
Rauvolfia macrophylla \\
Rauvolfia vomitoria \\
Ricinodendron heudelotii \\
Rinorea sp. \\
Sapium ellipticum \\
Sapium sp. \\
Santiria trimera \\
Sorindeia grandifolia \\
Stachyothyrsus \\
Staudtia kamerunensis \\
Sterculia rhinopetala \\
Sterculia tragacantha \\
Strombosia pustulata \\
Syzygium sp. \\
Tabernaemontana crassa \\
Tabernaemontana pachysiphon \\
Terminalia glaucescens \\
Terminalia superba \\
Tetrapleura tetraptera \\
Tieghemella africana \\
Trichilia dregeana \\
Trichilia tessmannii \\
Trichilia welwitschii \\
Trichoscypha acuminata \\
Triplochyton scleroxylon \\
Uapaca guineensis \\
Vitex grandifolia \\
Xylopia aethiopica \\
Xylopia hypolampra \\
Xylopia staudtii \\
Zanthoxylum tessmanii \\
\bottomrule()
\end{longtable}

\begin{longtable}[]{@{}l@{}}
\caption{\label{tab:FamillesNanga-Eboko3}Identification des familles dans la FC de Nanga-Eboko. (x: Nom des familles d'espèces)}\tabularnewline
\toprule()
x \\
\midrule()
\endfirsthead
\toprule()
x \\
\midrule()
\endhead
Achariaceae \\
Anarcadiaceae \\
Annonaceae \\
Apocynaceae \\
Aptandraceae \\
Arecaceae \\
Bignoniaceae \\
Burseraceae \\
Calophyllaceae \\
Cannabaceae \\
Clusiaceae \\
Combretaceae \\
Ehretiaceae \\
Euphorbiaceae \\
Fabaceae \\
Gentianaceae \\
Irvingiaceae \\
Lamiaceae \\
Lauraceae \\
Lecythidaceae \\
Malvaceae \\
Meliaceae \\
Moraceae \\
Myristicaceae \\
Myrtaceae \\
Ochnaceae \\
Phyllantaceae \\
Putranjivaceae \\
Rhamnaceae \\
Rhizophoraceae \\
Rubiaceae \\
Rutaceae \\
Sapindaceae \\
Sapotaceae \\
Simaroubaceae \\
Strombosiaceae \\
Ulmaceae \\
Urticaceae \\
Violoceae \\
\bottomrule()
\end{longtable}

Soit \(X_i\) le nombre d'individus de \(i\) espèces observées dans l'échantillon, \(i=1,2,...,n;\) \(X_i\) est considéré comme la fréquence des espèces de l'échantillon. Soit \(f_k\) le nombre d'espèces représentées exactement par \(k\) individus dans l'échantillon, \(k=0,1,...,N;\) \(f_k\) est considéré comme la fréquence d'abondance observée. \(f_1=390\ et\ f_1=189\) est le nombre d'espèces représentées exactement par 1 individu (singleton) dans une parcelle unique parmi celles inventoriées respectivement pour les FC de Ndikinikiméki et Nanga-Eboko. \(f_2=93\ et\ f_2=57\) est le nombre d'espèces représentées par 2 individus (doublon) dans l'échantillon d'inventaire respectivement pour les FC de Ndikiniméki et Nanga-Eboko. Pour la FC de Ndikiniméki le nombre d'observation de présence est de 608 tandis que la fréquence non-observable \(f_0=4128\) (Tableau:\ref{tab:ClassesAbondance-Ndikinimeki3}) indique le nombre d'espèces présentes dans l'assemblage des individus mais non observées dans l'échantillon d'inventaire. Ces absences représentent près de 87\% des données de la matrice d'inventaire et la valeur maximale d'abondance a été enregistrée dans une parcelle où une espèce a été observée 29 fois de suite (Figure:\ref{fig:Histogramme-Abondance3}). Ces résultats nous montrent que les espèces identifiées dans ces espaces ont des faibles probabilités d'être revues. En effet, lorsque le nombre de représentant d'une espèce augmente la probabilté d'observation diminue et vis versa.

\begin{longtable}[]{@{}lr@{}}
\caption{\label{tab:ClassesAbondance-Ndikinimeki3}Dénombrement des classes d'abondances dans la FC de Ndikiniméki. (Var1:Classes d'abondances; Freq:Fréquences d'abondances)}\tabularnewline
\toprule()
Var1 & Freq \\
\midrule()
\endfirsthead
\toprule()
Var1 & Freq \\
\midrule()
\endhead
0 & 4128 \\
1 & 390 \\
2 & 93 \\
3 & 46 \\
4 & 22 \\
5 & 12 \\
6 & 9 \\
7 & 7 \\
8 & 10 \\
9 & 4 \\
10 & 2 \\
12 & 1 \\
13 & 1 \\
14 & 1 \\
15 & 2 \\
17 & 2 \\
18 & 1 \\
22 & 1 \\
23 & 1 \\
26 & 2 \\
29 & 1 \\
\bottomrule()
\end{longtable}

\begin{longtable}[]{@{}lr@{}}
\caption{\label{tab:ClassesAbondanceNanga-Eboko3}Dénombrement des classes d'abondances dans la FC de Nanga-Eboko. (Var1:Classes d'abondances; Freq:Fréquences d'abondances)}\tabularnewline
\toprule()
Var1 & Freq \\
\midrule()
\endfirsthead
\toprule()
Var1 & Freq \\
\midrule()
\endhead
0 & 2473 \\
1 & 189 \\
2 & 57 \\
3 & 31 \\
4 & 13 \\
5 & 9 \\
6 & 11 \\
7 & 4 \\
8 & 2 \\
9 & 2 \\
10 & 3 \\
11 & 2 \\
12 & 3 \\
13 & 1 \\
16 & 2 \\
18 & 1 \\
23 & 2 \\
25 & 1 \\
\bottomrule()
\end{longtable}

\begin{figure}

{\centering \includegraphics[width=0.8\linewidth]{03-resultats_et_discussions_files/figure-latex/Histogramme-Abondance3-1} 

}

\caption{Histogrammes des effectifs des classes d'abondance des espèces.}\label{fig:Histogramme-Abondance3}
\end{figure}

Pour la FC de Nanga-Eboko la fréquence non-observable est de 2473 tandis que la valeur maximale d'abondance est de 25 (Tableau:\ref{tab:ClassesAbondanceNanga-Eboko3}).

Une projection des Coordonnées géographiques des parcelles (données GPS recueillies sur le terrain) permet de confirmer la localisation des parcelles d'inventaires dans la zone de transition Forêts-Savane (Figures:\ref{fig:Coordonnees-parcelles3}, \ref{fig:Localisation-occupationsolsNdiki3} et \ref{fig:Localisation-occupationsolsNanga3}).

\begin{figure}

{\centering \includegraphics[width=0.8\linewidth]{03-resultats_et_discussions_files/figure-latex/Coordonnees-parcelles3-1} 

}

\caption{Localisation géographique des parcelles d'inventaires }\label{fig:Coordonnees-parcelles3}
\end{figure}

\begin{verbatim}
#> Linking to GEOS 3.9.3, GDAL 3.5.2, PROJ 8.2.1; sf_use_s2()
#> is TRUE
\end{verbatim}

\begin{figure}

{\centering \includegraphics[width=0.8\linewidth]{03-resultats_et_discussions_files/figure-latex/Localisation-occupationsolsNdiki3-1} 

}

\caption{Localisation des sites par catégories d'occupation des terres dans FC de Ndikiniméki}\label{fig:Localisation-occupationsolsNdiki3}
\end{figure}

\begin{figure}

{\centering \includegraphics[width=0.8\linewidth]{03-resultats_et_discussions_files/figure-latex/Localisation-occupationsolsNanga3-1} 

}

\caption{Localisation des sites par catégories d'occupation des terres dans FC de Nanga-Eboko}\label{fig:Localisation-occupationsolsNanga3}
\end{figure}

En examinant l'allure des diagrammes de répartition des fréquences d'occurrences des espèces (Figure:\ref{fig:Frequences-occurences3}), classées en rang de fréquences croissantes, on observe que le nombre d'espèces est inversement proportionnel à la fréquence d'occurrence. Lorsque la fréquence d'occurrence augmente, le nombre d'espèces diminue cela pourrait suggerer que ces forêts renferment de nombreuses espèces rares qui n'apparaissent que quelques fois et qu'il y aurait très peu d'espèces communes qui apparaitraient très souvent (Figure:\ref{fig:Nombres-occurences3}). En effet 3/4 des espèces repertoriées ont une fréquence d'occurence inférieure à 10\% ce qui montre que la majorité des espèces sont rares et n'apparaissent que de manière épars.

\begin{figure}

{\centering \includegraphics[width=0.8\linewidth]{03-resultats_et_discussions_files/figure-latex/Nombres-occurences3-1} 

}

\caption{Histogrammes de répartition des espèces}\label{fig:Nombres-occurences3}
\end{figure}

\begin{figure}

{\centering \includegraphics[width=0.8\linewidth]{03-resultats_et_discussions_files/figure-latex/Frequences-occurences3-1} 

}

\caption{Fréquences d'occurences des espèces.}\label{fig:Frequences-occurences3}
\end{figure}

En considérant les 32 parcelles prises chacune séparément dans la FC de Ndikiniméki, la richesse spécifique observée la plus faible a été retrouvée dans la parcelle (P19) avec 05 espèces enregitrées tandis que la plus élévée est de 33 espèces et se situe dans la parcelle (P25) (Figure:\ref{fig:Richesses-specifiquesparcelles3}). 75\% des parcelles inventoriées comptent 22 espèces et la moyenne générale se situe autour de 19 espèces (Figure:\ref{fig:Cartogrphiepools-specifiques3}). Dans la FC de Nanga-Eboko, la richesse spécifique observée la plus faible a été enregistrée dans la parcelle (P3) avec une seule espèce tandis que la richesse spécifique maximale a été relevée dans la parcelle (P11) avec 32 espèces.

\begin{longtable}[]{@{}lr@{}}
\caption{\label{tab:NdikiNombre-especesparcelles3}Classement de la richesse spécifiques par parcelles dans la FC de Ndikiniméki.(x:Richesse spécifiques par parcelles)}\tabularnewline
\toprule()
& x \\
\midrule()
\endfirsthead
\toprule()
& x \\
\midrule()
\endhead
19 & 5 \\
5 & 8 \\
22 & 9 \\
3 & 10 \\
2 & 12 \\
9 & 12 \\
18 & 12 \\
17 & 13 \\
26 & 13 \\
27 & 14 \\
10 & 15 \\
14 & 15 \\
6 & 16 \\
13 & 16 \\
12 & 18 \\
11 & 19 \\
15 & 20 \\
16 & 20 \\
28 & 20 \\
4 & 21 \\
1 & 22 \\
7 & 22 \\
8 & 22 \\
29 & 22 \\
20 & 24 \\
23 & 25 \\
31 & 28 \\
30 & 29 \\
24 & 30 \\
32 & 31 \\
21 & 32 \\
25 & 33 \\
\bottomrule()
\end{longtable}

\begin{longtable}[]{@{}lr@{}}
\caption{\label{tab:NangaNombre-especesparcelles3}Classement de la richesse spécifiques par parcelles dans la FC de Nanga-Eboko. (x:Richesse spécifiques par parcelles)}\tabularnewline
\toprule()
& x \\
\midrule()
\endfirsthead
\toprule()
& x \\
\midrule()
\endhead
3 & 1 \\
16 & 3 \\
17 & 3 \\
4 & 4 \\
5 & 6 \\
9 & 7 \\
12 & 8 \\
10 & 10 \\
14 & 13 \\
19 & 13 \\
8 & 14 \\
7 & 15 \\
2 & 17 \\
6 & 17 \\
23 & 18 \\
22 & 19 \\
1 & 20 \\
13 & 22 \\
15 & 22 \\
18 & 22 \\
21 & 23 \\
20 & 24 \\
11 & 32 \\
\bottomrule()
\end{longtable}

\begin{figure}

{\centering \includegraphics[width=0.8\linewidth]{03-resultats_et_discussions_files/figure-latex/Richesses-specifiquesparcelles3-1} 

}

\caption{Nombres d'espèces par parcelles}\label{fig:Richesses-specifiquesparcelles3}
\end{figure}

Une projection de la richesse spécifique de ces parcelles d'échantillonnage selon leur coordonnées géographiques nous montre que la richesse la richesse spécifique semble indépendante de la longitude. En la cartographie de la richesse spécifique dans les deux forêts communales montre qu'elle est plus importante au delà d'un seuil. Dans la FC Ndikiniméki les sites ayant une richesse spécifique importante sont situés au dessus de 4°75'. Dans la FC de Nanga-Eboko cette limite se situe au delà de 4°8' (Figure:\ref{fig:Cartogrphiepools-specifiques3}).

\begin{figure}

{\centering \includegraphics[width=0.8\linewidth]{03-resultats_et_discussions_files/figure-latex/Cartogrphiepools-specifiques3-1} 

}

\caption{Cartographie des richesses spécifiques}\label{fig:Cartogrphiepools-specifiques3}
\end{figure}

Nous observons que sur les 1348 individus constituant la taille de l'échantillon dans la FC de Ndikiniméki, la parcelle (P27) enregistre l'effectif le plus faible (27 individus) et présente une richesse spécifiques de 14 espèces:

\begin{Shaded}
\begin{Highlighting}[]
\CommentTok{\# nombre d\textquotesingle{}abondance par parcelles dans FC de Ndikiniméki}
\NormalTok{(Ndiki\_abund}\OtherTok{\textless{}{-}}\FunctionTok{rowSums}\NormalTok{(Ndiki\_spec))}
\CommentTok{\#\textgreater{}  1  2  3  4  5  6  7  8  9 10 11 12 13 14 15 16 17 18 19 20 }
\CommentTok{\#\textgreater{} 42 50 35 50 39 37 34 52 35 46 33 46 43 29 35 37 42 38 34 38 }
\CommentTok{\#\textgreater{} 21 22 23 24 25 26 27 28 29 30 31 32 }
\CommentTok{\#\textgreater{} 47 39 45 51 56 42 27 41 58 47 39 61}
\CommentTok{\# Abondance totale des individus dans la FC de Ndikiniméki}
\FunctionTok{sum}\NormalTok{(Ndiki\_abund)}
\CommentTok{\#\textgreater{} [1] 1348}
\CommentTok{\# Minimum et Maximum d\textquotesingle{}abondance par parcelles dans la FC de Ndikiniméki}
\FunctionTok{range}\NormalTok{(Ndiki\_abund)}
\CommentTok{\#\textgreater{} [1] 27 61}
\CommentTok{\# Abondance dans le site ayant la plus petite richesse spécifique observée dans la FC de Ndikiniméki}
\NormalTok{Ndiki\_abund[Ndiki\_nbsp}\SpecialCharTok{==}\FunctionTok{min}\NormalTok{(Ndiki\_nbsp)]}
\CommentTok{\#\textgreater{} 19 }
\CommentTok{\#\textgreater{} 34}
\CommentTok{\# Abondance dans le site ayant la plus grande richesse spécifique observee dans la FC de Ndikiniméki}
\NormalTok{Ndiki\_abund[Ndiki\_nbsp}\SpecialCharTok{==}\FunctionTok{max}\NormalTok{(Ndiki\_nbsp)]}
\CommentTok{\#\textgreater{} 25 }
\CommentTok{\#\textgreater{} 56}
\end{Highlighting}
\end{Shaded}

La parcelle (P19) a un éffectif presque équivalent (34 individus), cependant elle présente la richesse spécifique observée la plus faible (5 espèces). La parcelle (P32) est celle présentant la plus forte abondance (61 individus) avec une richesse spécifiques observée de 31 espèces. Dans la FC de Nanga-Eboko 843 individus ont constitué la taille de l'échantillon. La parcelle (P3) a affiché l'effectif le plus réduit (01 individu) et la richesse spécifique la plus faible (01 espèce). Par contre la parcelle (P20) présente l'abondance la plus importante (65 individus) pour une richesse spécifiques de 24 espèces:

\begin{Shaded}
\begin{Highlighting}[]
\NormalTok{(Nanga\_abund}\OtherTok{\textless{}{-}}\FunctionTok{rowSums}\NormalTok{(Nanga\_spec))}
\CommentTok{\#\textgreater{}  1  2  3  4  5  6  7  8  9 10 11 12 13 14 15 16 17 18 19 20 }
\CommentTok{\#\textgreater{} 50 64  1  5 16 39 32 49 33 36 47 32 48 41 41 11  4 47 30 65 }
\CommentTok{\#\textgreater{} 21 22 23 }
\CommentTok{\#\textgreater{} 43 46 63}
\CommentTok{\# Abondance totale des individus dans la FC de Nanga{-}Eboko}
\FunctionTok{sum}\NormalTok{(Nanga\_abund)}
\CommentTok{\#\textgreater{} [1] 843}
\CommentTok{\#Minimum et Maximum d\textquotesingle{}abondance par parcelles dans la FC de Nanga{-}Eboko}
\FunctionTok{range}\NormalTok{(Nanga\_abund)}
\CommentTok{\#\textgreater{} [1]  1 65}
\CommentTok{\# Abondance dans la parcelle ayant la plus petite richesse spécifique observée dans la FC de Nanga{-}Eboko}
\NormalTok{Nanga\_abund[Nanga\_nbsp}\SpecialCharTok{==}\FunctionTok{min}\NormalTok{(Nanga\_nbsp)]}
\CommentTok{\#\textgreater{} 3 }
\CommentTok{\#\textgreater{} 1}
\CommentTok{\# Abondance dans la parcelle ayant la plus grande richesse spécifique observée dans la FC de Nanga{-}Eboko}
\NormalTok{Nanga\_abund[Nanga\_nbsp}\SpecialCharTok{==}\FunctionTok{max}\NormalTok{(Nanga\_nbsp)]}
\CommentTok{\#\textgreater{} 11 }
\CommentTok{\#\textgreater{} 47}
\end{Highlighting}
\end{Shaded}

Paradoxallement la parcelle (P11) compte 47 individus pour une richesse spécifique de 32 espèces ce qui fait d'elle la parcelle la plus importante de toute de toute FC de Nanga-Eboko en terme de richesse spécifique (Figure:\ref{fig:Abondance-individus3}).

\begin{figure}

{\centering \includegraphics[width=0.8\linewidth]{03-resultats_et_discussions_files/figure-latex/Abondance-individus3-1} 

}

\caption{Nombres d'individus par parcelle}\label{fig:Abondance-individus3}
\end{figure}

Pour analyser la rarefaction des espèces par rapport au nombre d'individus échantillonnés par parcelles, nous avons standardisé les unités d'échantillonnages de la FC de Ndikiniméki à un effectif de 47 individus correspondant au nombre moyen d'individus enregistrés dans les 3/4 des parcelles inventoriées et appartenant à 25\% des espèces les plus communes:

\begin{Shaded}
\begin{Highlighting}[]
\CommentTok{\# Analyse de la rarefaction au $3\^{}ème\^{}$ quartile dans la FC de Ndikiniméki}
\FunctionTok{summary}\NormalTok{(Ndiki\_abund)}
\CommentTok{\#\textgreater{}    Min. 1st Qu.  Median    Mean 3rd Qu.    Max. }
\CommentTok{\#\textgreater{}   27.00   36.50   41.50   42.12   47.00   61.00}
\NormalTok{Ndiki\_rare47}\OtherTok{\textless{}{-}}\FunctionTok{rarefy}\NormalTok{(Ndiki\_spec,}\AttributeTok{sample =} \DecValTok{47}\NormalTok{)}
\CommentTok{\# Classement des sites selon l\textquotesingle{}abondance observée dans la FC de Ndikiniméki.}
\FunctionTok{sort}\NormalTok{(Ndiki\_nbsp)}
\CommentTok{\#\textgreater{} 19  5 22  3  2  9 18 17 26 27 10 14  6 13 12 11 15 16 28  4 }
\CommentTok{\#\textgreater{}  5  8  9 10 12 12 12 13 13 14 15 15 16 16 18 19 20 20 20 21 }
\CommentTok{\#\textgreater{}  1  7  8 29 20 23 31 30 24 32 21 25 }
\CommentTok{\#\textgreater{} 22 22 22 22 24 25 28 29 30 31 32 33}
\CommentTok{\# Classement des sites selon la rarefaction au $3\^{}ème\^{}$ quartile dans la FC de Ndikiniméki}
\FunctionTok{sort}\NormalTok{(}\FunctionTok{round}\NormalTok{(Ndiki\_rare47))}
\CommentTok{\#\textgreater{} 19  5 22  3  2  9 18 17 26 27 10 14  6 13 12 11 29  4 15 16 }
\CommentTok{\#\textgreater{}  5  8  9 10 12 12 12 13 13 14 15 15 16 16 18 19 19 20 20 20 }
\CommentTok{\#\textgreater{} 28  8  1  7 20 23 32 31 24 25 30 21 }
\CommentTok{\#\textgreater{} 20 21 22 22 24 25 26 28 29 29 29 32}
\CommentTok{\# Site ayant le Min de richesses spécifiques estimées dans la FC de Ndikiniméki}
\NormalTok{Ndiki\_rare47[Ndiki\_rare47}\SpecialCharTok{==}\FunctionTok{min}\NormalTok{(Ndiki\_rare47)]}
\CommentTok{\#\textgreater{} 19 }
\CommentTok{\#\textgreater{}  5}
\CommentTok{\# site ayant le Max de richesses spécifiques observées dans la FC de Ndikiniméki}
\NormalTok{Ndiki\_rare47[Ndiki\_rare47}\SpecialCharTok{==}\FunctionTok{max}\NormalTok{(Ndiki\_rare47)]}
\CommentTok{\#\textgreater{} 21 }
\CommentTok{\#\textgreater{} 32}
\end{Highlighting}
\end{Shaded}

On a constaté dans ces conditions que la richesse spécifique varie de 5 à 32 espèces respectivement pour l'unité la plus pauvre et l'unité la plus riche. On a observé également que l'unité la plus pauvre en terme de richesse spécifiques observées (5 espèces) demeure la parcelle (P19), cependant la plus riche correspond désormais à la parcelle (P21) avec 32 espèces trouvées. De même dans la FC de Nanga-Eboko les unités ont été ramenées à un effectif de 48 individus par unité d'échantillonnage:

\begin{Shaded}
\begin{Highlighting}[]
\CommentTok{\# Analyse de la rarefaction au $3\^{}ème\^{}$ quartile dans la FC de Nanga{-}Eboko}
\FunctionTok{summary}\NormalTok{(Nanga\_abund)}
\CommentTok{\#\textgreater{}    Min. 1st Qu.  Median    Mean 3rd Qu.    Max. }
\CommentTok{\#\textgreater{}    1.00   31.00   41.00   36.65   47.50   65.00}
\NormalTok{Nanga\_rare48}\OtherTok{\textless{}{-}}\FunctionTok{rarefy}\NormalTok{(Nanga\_spec,}\AttributeTok{sample =} \DecValTok{48}\NormalTok{)}
\CommentTok{\# Classement des sites selon l\textquotesingle{}abondance observée dans la FC de Nanga{-}Eboko}
\FunctionTok{sort}\NormalTok{(Nanga\_nbsp)}
\CommentTok{\#\textgreater{}  3 16 17  4  5  9 12 10 14 19  8  7  2  6 23 22  1 13 15 18 }
\CommentTok{\#\textgreater{}  1  3  3  4  6  7  8 10 13 13 14 15 17 17 18 19 20 22 22 22 }
\CommentTok{\#\textgreater{} 21 20 11 }
\CommentTok{\#\textgreater{} 23 24 32}
\CommentTok{\# Classement des sites selon la rarefaction au $3\^{}ème\^{}$ quartile dans la FC de Nanga{-}Eboko}
\FunctionTok{sort}\NormalTok{(}\FunctionTok{round}\NormalTok{(Nanga\_rare48))}
\CommentTok{\#\textgreater{}  3 16 17  4  5  9 12 10 14 19  2  8  7 23  6 20 22  1 13 15 }
\CommentTok{\#\textgreater{}  1  3  3  4  6  7  8 10 13 13 14 14 15 16 17 19 19 20 22 22 }
\CommentTok{\#\textgreater{} 18 21 11 }
\CommentTok{\#\textgreater{} 22 23 32}
\CommentTok{\# Site ayant le Min de richesses spécifiques estimées dans la FC de Nanga{-}Eboko}
\NormalTok{Nanga\_rare48[Nanga\_rare48}\SpecialCharTok{==}\FunctionTok{min}\NormalTok{(Nanga\_rare48)]}
\CommentTok{\#\textgreater{} 3 }
\CommentTok{\#\textgreater{} 1}
\CommentTok{\# Site ayant le Max de richesses spécifiques estimées dans la FC de Nanga{-}Eboko}
\NormalTok{Nanga\_rare48[Nanga\_rare48}\SpecialCharTok{==}\FunctionTok{max}\NormalTok{(Nanga\_rare48)]}
\CommentTok{\#\textgreater{} 11 }
\CommentTok{\#\textgreater{} 32}
\end{Highlighting}
\end{Shaded}

Les résultats obtenus en terme de richesse spécifiques ont été les mêmes que ceux observés dans les échantillonnages à base de parcelles pour ce qui concerne le minimum et le maximum de richesse : ce sont respectivement les parcelles (P3) et (P11).

Pour quantifier la force et la direction des relations entre les différents indices de diversité et d'équité, nous avons utilisé un graphique matriciel contenant des nuages de points et des coefficients de correlation de Pearson établis à partir des paires d'indice. Ces paramètres sont les suivant:

\begin{itemize}
\tightlist
\item
  N0: Richesse spécifique;
\item
  N1: Diversité de Shannon (base e);
\item
  N1b2: Diversité de Shannon (base 2);
\item
  N2: Diversité de Simpson;
\item
  H: Entropie de Shannon (base e);
\item
  Hb2: Entropie de Shannon (base 2);
\item
  J: Equitabilité de Pielou;
\item
  E10: Equitabilité de Shannon (Ratio de Hill);
\item
  E20: Equitabilité de Simpson (Ratio de Hill
\end{itemize}

La matrice de correlation des indices de diversité et d'équitabilité dans la FC de Ndikiniméki (Figure: \ref{fig:Matrice-corrNdiki3}) nous montre une très forte correlation entre la richesse spécifique (N0) et la diversité de Shannon (N1) corr=0,971. Ceci laisse penser qu'il y aurait un équilibre entre le nombre d'espèces et leur abondance relative, traduisant ainsi une meilleure resilience et une stabilité de l'écosystème face aux perturbations environnementales. De même la très forte correlation (corr=1) entre les deux méthodes de mesure de la diversité de Shannon (N1b2 et N1) permet d'affirmer que les différences observées dans les valeurs des indices sont effectivement significatives et non liées aux variations de méthodes. On relève également une très forte correlation entre la diversité de Simpson (N2) et la diversité de Shannon (N1), corr=0,949. Ce résultat signifit que lorsque l'indice de Simpson augmente (diminution de la dominance des espèces), l'indice de Shannon a tendance également à augmenter (augmentation de l'équitabilité). Par contre on note une faible correlation entre l'entropie de Shannon et l'équitabilité de Pielou (corr=0,117) ce qui signifit que ces deux mesures de diversité ne sont pas liées et capturent chacun des aspects différents de la biodiversité. En effet l'entropie de Shannon évalue la biodiversté en tenant compte à la fois de la richesse spécifique et l'équité, en revanche l'équitabilité de Pielou mesure la répartition des individus entre les différentes espèces, en fonction de la proportion de chaque espèce.

\begin{verbatim}
#> Registered S3 method overwritten by 'GGally':
#>   method from   
#>   +.gg   ggplot2
\end{verbatim}

\begin{figure}

{\centering \includegraphics[width=0.8\linewidth]{03-resultats_et_discussions_files/figure-latex/Matrice-corrNdiki3-1} 

}

\caption{Matrice de correlation des indices de diversité et d'équitabilité dans la FC de Ndikiniméki}\label{fig:Matrice-corrNdiki3}
\end{figure}

Dans la matrice de correlation de la FC de Nanga-Eboko, les annotations suivantes ont été utilisées:

\begin{itemize}
\tightlist
\item
  NO\_ng: Richesse spécifique;
\item
  N1\_ng: Diversité de Shannon (base e);
\item
  N1b2\_ng: Diversité de Shannon (base 2);
\item
  N2\_ng: Diversité de Simpson;
\item
  H\_ng: Entropie de Shannon (base e);
\item
  Hb2\_ng: Entropie de Shannon (base 2);
\item
  J\_ng: Equitabilité de Pielou;
\item
  E10\_ng: Equitabilité de Shannon (Ratio de Hill);
\item
  E20\_ng: Equitabilité de Simpson (Ratio de Hill).
\end{itemize}

Cette matrice montre une forte correlation entre la richesse spécifiques (N0\_ng) et la diversité de Shannon (N1\_ng) traduisant ainsi un équilibre entre le nombre d'espèces et l'abondance relative (corr=0,994). On note également une forte correlation entre la diversité de Simpson et la diversité de Shannon (corr=0,996) cela suggère une grande diversité et une distribution plus équitable des individus parmi les espèces(Figure: \ref{fig:Matrice-corrNanga3}).

\begin{figure}

{\centering \includegraphics[width=0.8\linewidth]{03-resultats_et_discussions_files/figure-latex/Matrice-corrNanga3-1} 

}

\caption{Matrice de correlation des indices de diversité et d'équité dans la FC de Nanga-Eboko}\label{fig:Matrice-corrNanga3}
\end{figure}

Si l'on se refère aux résultats du Tableau:\ref{tab:Comparaisonrichesse-observeeNdikinimeki3} pour analyser la diversité des espèces dans les différentes catégories d'utilisation des terres, on pourrait conclure que les Forêts Marécageuses (FMa) ont la richesse spécifiques la plus faible puisque les observations nous montrent qu'elles comptent 22 espèces tandis que le résultat est plus élevée pour les autres catégories d'utilisation des terres (plus d'une centaine d'espèces pour les Forêts Matures). Notons également dans le même tableau que, \(n\) représente le nombre d'unités d'échantillonnages pour chaque catégorie d'utilisation des terres. En regardant les différences de valeurs attribuées à \(n\) dans les différentes catégories, on constate que les Forêts marécageuses et les Jachères ont les valeurs les plus faibles. Par contre dans la forêt communale de Nanga-Eboko,les Cultures et les Savanes enregistrent le plus faible nombre d'espèces avec respectivement 3 et 7 espèces observées (Tableau:\ref{tab:ComparaisonrichesseNanga-Eboko3}).

\begin{table}
\centering
\caption{\label{tab:Comparaisonrichesse-observeeNdikinimeki3}Richesse spécifique observée des differentes catégories d'utilisation des terres dans la FC de Ndikiniméki.}
\centering
\resizebox{\ifdim\width>\linewidth\linewidth\else\width\fi}{!}{
\begin{threeparttable}
\begin{tabular}[t]{lrr}
\toprule
  & n & richness\\
\midrule
CAC & 8 & 48\\
FM & 13 & 112\\
FMa & 2 & 22\\
FS & 7 & 77\\
JA & 2 & 28\\
\bottomrule
\end{tabular}
\begin{tablenotes}
\item \textit{Note: } 
\item n: nombre de parcelle; richness: richesse spécifique; FM: Forêt Mature; FS: Forêt Secondaire; FMa: Forêt Marécageuse; JA: Jachère; CAC: Cacaoyère
\end{tablenotes}
\end{threeparttable}}
\end{table}

\begin{table}
\centering
\caption{\label{tab:ComparaisonrichesseNanga-Eboko3}Richesse spécifique observée des differentes catégories d'utilisation des terres dans la FC de Nanga-Eboko.}
\centering
\resizebox{\ifdim\width>\linewidth\linewidth\else\width\fi}{!}{
\begin{threeparttable}
\begin{tabular}[t]{lrr}
\toprule
  & n & richness\\
\midrule
CUL & 2 & 3\\
FGA & 3 & 38\\
FSA & 7 & 86\\
FSJ & 4 & 47\\
MAR & 4 & 30\\
\addlinespace
SAV & 3 & 7\\
\bottomrule
\end{tabular}
\begin{tablenotes}
\item \textit{Note: } 
\item n: nombre de parcelle; richness: richesse spécifique; FSA: Forêts Secondaires Adultes; FGA: Forêts Galeries; FSJ: Forêts Secondaires Jeunes; MAR: Marécages; SAV: Savanes; CUL: Cultures.
\end{tablenotes}
\end{threeparttable}}
\end{table}

Etant donné que la richesse spécifique est dépendante de la taille de l'échantillon, il n'est pas aisé de comparer la richesse spécifique entre des catégories d'occupation des terres de tailles différentes. Ces différences pouvant être liées soit à la taille de l'échantillon ou à la différence de richesse spécifique entre les catégories d'utilisation des terres proprement dites. Il faut également noter que la somme des espèces lorsque plus d'un site sont combinés n'est pas toujours égale à la somme arithmétique des espèces de cette combinaison. Ceci ne peut être le cas que si les sites combinés partagent exactement le même nombre d'espèces. Ainsi, pour déterminer le nombre d'espèces dans des combinaisons de sites, on a calculé la moyenne cumulée de la richesse spécifique des différentes combinaison de sites. A cause du fait qu'il n'existe pas une valeur unique pour la richesse spécifique des sites, nous avons calculé la valeur moyenne. La richesse spécifique moyenne pour un site est de 19 espèces pour la FC de Ndiniméki:

Richesse spécifique moyennes issues de la combinaison des sites dans la FC de Ndikiniméki

\begin{verbatim}
#> Species Accumulation Curve
#> Accumulation method: exact
#> Call: specaccum(comm = x, method = method, permutations = permutations,      conditioned = conditioned, gamma = gamma, digits = 2) 
#> 
#>                                                          
#> Sites     1.00000  2.000000  3.000000  4.000000  5.000000
#> Richness 19.00000 33.808468 45.913105 56.084789 64.802583
#> sd        7.30582  8.539923  8.916379  8.952395  8.827465
#>                                                          
#> Sites     6.000000  7.000000  8.000000  9.00000 10.000000
#> Richness 72.392614 79.086427 85.052813 90.41776 95.277545
#> sd        8.621373  8.373942  8.106497  7.83050  7.552434
#>                                                    
#> Sites    11.000000  12.000000  13.000000  14.000000
#> Richness 99.707430 103.767500 107.506603 110.965074
#> sd        7.275762   7.002243   6.732646   6.467143
#>                                                     
#> Sites     15.000000  16.000000  17.000000  18.000000
#> Richness 114.176631 117.169746 119.968655 122.594117
#> sd         6.205533   5.947362   5.691989   5.438611
#>                                                    
#> Sites     19.000000  20.000000  21.000000  22.00000
#> Richness 125.064008 127.393796 129.596916 131.68509
#> sd         5.186281   4.933863   4.680072   4.42331
#>                                                     
#> Sites     23.000000  24.000000  25.000000  26.000000
#> Richness 133.668561 135.556350 137.356393 139.075708
#> sd         4.161733   3.893036   3.614284   3.321622
#>                                                     
#> Sites     27.000000  28.000000  29.000000  30.000000
#> Richness 140.720513 142.296329 143.808065 145.260081
#> sd         3.009619   2.670734   2.292332   1.847428
#>                        
#> Sites     31.000000  32
#> Richness 146.656250 148
#> sd         1.265143   0
\end{verbatim}

Pour la FC de Nanga-Eboko cette moyenne est 14 espèces:

Richesse spécifiques moyennes issues de la combinaison des sites dans la FC de Nanga-Eboko

\begin{verbatim}
#> Species Accumulation Curve
#> Accumulation method: exact
#> Call: specaccum(comm = x, method = method, permutations = permutations,      conditioned = conditioned, gamma = gamma) 
#> 
#>                                                           
#> Sites     1.000000  2.000000  3.000000  4.000000  5.000000
#> Richness 14.478261 26.387352 36.407679 45.007792 52.516687
#> sd        7.982969  9.506421 10.018949 10.136834 10.076833
#>                                                           
#> Sites     6.000000  7.000000  8.000000  9.000000 10.000000
#> Richness 59.170406 65.142031 70.561091 75.526289 80.114017
#> sd        9.928819  9.732144  9.503957  9.251164  8.975833
#>                                                           
#> Sites    11.000000 12.000000 13.000000 14.000000 15.000000
#> Richness 84.384209 88.384467 92.153049 95.721073 99.114178
#> sd        8.677656  8.355357  8.006420  7.628199  7.216774
#>                                                    
#> Sites     16.000000  17.000000  18.00000  19.000000
#> Richness 102.353769 105.457983 108.44242 111.320723
#> sd         6.767327   6.273166   5.72600   5.110619
#>                                              
#> Sites     20.000000  21.000000  22.000000  23
#> Richness 114.105025 116.806324 119.434783 122
#> sd         4.402379   3.541358   2.428563   0
\end{verbatim}

La courbe cumulative des espèces a permis de calculer la richesse spécifique moyenne pour une large combinaison de sites en considérant la présence et non l'éffectif des individus constituant les espèces. Le résultat nous a montré que la richesse spécifique moyenne croît rapidement avec le nombre de sites combinés mais lorsque la somme des sites combinés atteint un certain seuil ce taux d'accroissement devient faible. Dans la FC de Ndikiniméki par exemple, la richesse spécifique moyenne pour toutes les combinaisons possibles de 23 sites est de 133,66 tandis que pour une combinaison de 25 sites elle est de 137,55. \(sd\) indique l'écart type observé lorsque par chance certaines combinaisons de sites présentent un nombre élevé d'espèces tandis que d'autres en contiennent un faible nombre. Au regard de la courbe cumulée des espèces, l'axe horizontal représente l'ensemble des 32 et 23 sites cumulés respectivement pour les FC de Ndikiniméki (Figure:\ref{fig:Courbeaccumulation-Ndikimeki3}) et Nanga-Eboko (Figure:\ref{fig:CourbeaccumulationNanga-Eboko3}) tandis que l'axe vertical représente l'ensemble des 148 et 122 espèces cumulées respectivement pour chacune des forêts communales.

\begin{figure}

{\centering \includegraphics[width=0.8\linewidth]{03-resultats_et_discussions_files/figure-latex/Courbeaccumulation-Ndikimeki3-1} 

}

\caption{Courbe d'accumulation des espèces de la FC de Ndikiniméki.}\label{fig:Courbeaccumulation-Ndikimeki3}
\end{figure}

\begin{figure}

{\centering \includegraphics[width=0.8\linewidth]{03-resultats_et_discussions_files/figure-latex/CourbeaccumulationNanga-Eboko3-1} 

}

\caption{Courbe d'accumulation des espèces de la FC de Nanga-Eboko.}\label{fig:CourbeaccumulationNanga-Eboko3}
\end{figure}

Il a également été important de rechercher et de comparer le nombre moyen d'espèces par site et selon les différentes catégories d'utilisation des terres afin de resoudre le problème de la taille des échantillons. Pour déterminer si la différence de richesse spécifique entre les différentes catégories pourraient être liées à la taille des échantillons, on a construit une courbe cumulée des espèces pour chaque catégorie d'utilisation des terres et on les a comparé par la suite. Ces courbes présentent la richesse spécifique moyenne cumulée pour une large combinaison de sites dans chaque catégorie d'utilisation des terres jusqu'à ce que toutes les combinaisons de tailles possibles d'échantillons soient épuisées. En se référant à la moyenne cumulée de la richesse spécifique dans chaque catégorie d'utilisation des terres de la Forêt Communale de Ndikiniméki, il a été possible de comparer la richesse spécifique sur des échantillons de tailles comparables. Par exemple à l'échelle de 2 parcelles, on peut observer que les Forêts Matures (FM) ont une plus grande richesse spécifique (41 espèces) que les quatre autres catégories (36, 28, 22 et 18 espèces). De plus, il y a désormais une légère évidence que les Forêts marécageuses (FMa) ont une richesse spécifique plus importante (22 espèces) que les Cacaoyères (18 espèces) bien que le nombre total d'espèces qui avaient été observé plus haut était élevé (Figure:\ref{fig:Courbeaccumulation-comparaisonNdiki3}). Par contre dans la FC de Nanga-Eboko, la différence de richesse spécifique entre les types d'utilisation des terres n'a pas été influencé par la taille de l'échantillon, les différences observées sont semblables aux différences issues des calculs (figure:\ref{fig:Courbeaccumulation-comparaisonNanga3}).

\begin{figure}

{\centering \includegraphics[width=0.8\linewidth]{03-resultats_et_discussions_files/figure-latex/Courbeaccumulation-comparaisonNdiki3-1} 

}

\caption{Comparaison de la richesse spécifique entre les catégories d'utilisation des terres dans la Forêt Communale de Ndikiniméki.}\label{fig:Courbeaccumulation-comparaisonNdiki3}
\end{figure}

\begin{figure}

{\centering \includegraphics[width=0.8\linewidth]{03-resultats_et_discussions_files/figure-latex/Courbeaccumulation-comparaisonNanga3-1} 

}

\caption{Comparaison de la richesse spécifique entre les catégories d'utilisation des terres dans la Forêt Communale de Nanga-Eboko.}\label{fig:Courbeaccumulation-comparaisonNanga3}
\end{figure}

Pour déterminer l'ensemble des espèces qui pourraient être identifiées sur toute l'étendue de la zone d'étude, l'on a fait des extrapolations à partir de la courbe d'accumulation des espèces jusqu'atteindre le point de l'axe des abcisses correspondant à l'entièreté de la zone d'étude. Certaines prédictions du nombre total des espèces de la zone d'étude ont été faites à partir des formules de Chao, bootstrap et Jackknife au premier et second ordre. Ces estimations nous ont révélé que le nombre total des espèces dans l'ensemble de zone d'étude de la forêt communale de Ndikiniméki pourrait varier entre 168 et 206, suggérant ainsi que l'échantillonnage choisi en amont aurait capturé près de 75\% des espèces de la zone (Tableau:\ref{tab:Estimationespeces-Ndikinimeki3}). Dans la forêt communale de Nanga-Eboko il pourrait osciller entre 146 et 226 espèces suggérant une capture inférieure aux 3/4 des espèces potentielles contenues dans la zone d'étude (Tableau:\ref{tab:EstimationespecesNanga-Eboko3}).

\begin{longtable}[]{@{}lrrrr@{}}
\caption{\label{tab:Estimationespeces-Ndikinimeki3}Nombre total estimé des espèces dans la Forêt Communale de Ndikiniméki.}\tabularnewline
\toprule()
& jack1 & jack2 & chao & boot \\
\midrule()
\endfirsthead
\toprule()
& jack1 & jack2 & chao & boot \\
\midrule()
\endhead
pooled & 189.66 & 206.38 & 182.45 & 167.87 \\
\bottomrule()
\end{longtable}

\begin{longtable}[]{@{}lrrrr@{}}
\caption{\label{tab:EstimationespecesNanga-Eboko3}Nombre total estimé des espèces dans la Forêt Communale de Nanga-Eboko.}\tabularnewline
\toprule()
& jack1 & jack2 & chao & boot \\
\midrule()
\endfirsthead
\toprule()
& jack1 & jack2 & chao & boot \\
\midrule()
\endhead
pooled & 178.43 & 218.36 & 226.05 & 145.98 \\
\bottomrule()
\end{longtable}

La construction des courbes d'abondance a permis d'analyser la structure de la diversité dans les deux forêts communales. Après avoir déterminé le nombre total d'individus par catégories (espèces et familles) et les ranger suivant un ordre décroissant (Tableaux:\ref{tab:Tableauabondances-Ndikinimeki3} et \ref{tab:TableauabondancesNanga-Eboko3}),les courbes d'abondance présentent un étalement suivant l'axe des abscisses indicant ainsi une richesse spécifique importante. La raideur des pentes des courbes sur l'une des portions indique une faible équitabilité dans la distribution des espèces. \emph{Theobroma cacao} est l'espèce classée en première position en terme d'abondance (165 individus enregistrés) dans la FC de Ndikiniméki (Figure:\ref{fig:Courbesabondance-Ndikinimeki3}) tandis que l'espèce \emph{Raphia sp.} est en pool position dans la FC de Nanga-Eboko (Figure:\ref{fig:CourbesabondanceNanga-Eboko3}). Dans la FC de Ndikiniméki la famille la plus abondante est celle des Sterculiaceae suivit des Myristicaceae et des Fabaceae (Figure:\ref{fig:Courbesabondance-Ndikinimeki3}). Dans la FC de Nanga-Eboko la famille la plus importante en terme d'effectif est celle des Annonaceae suivit des Aracaceae et des Malvaceae (Figure:\ref{fig:CourbesabondanceNanga-Eboko3}).

\begin{table}
\centering
\caption{\label{tab:Tableauabondances-Ndikinimeki3}Les 10 plus abondantes espèces de la Forêt Communale de Ndikiniméki.}
\centering
\resizebox{\ifdim\width>\linewidth\linewidth\else\width\fi}{!}{
\begin{threeparttable}
\begin{tabular}[t]{lrrrrrrrr}
\toprule
  & rank & abundance & proportion & plower & pupper & accumfreq & logabun & rankfreq\\
\midrule
Theobroma cacao & 1 & 165 & 12.2 & 4.1 & 20.4 & 12.2 & 2.2 & 0.7\\
Pycnanthus angolensis & 2 & 142 & 10.5 & 6.8 & 14.3 & 22.8 & 2.2 & 1.4\\
Musa sp. & 3 & 63 & 4.7 & 1.0 & 8.3 & 27.4 & 1.8 & 2.0\\
Sterculia rhinopetala & 4 & 58 & 4.3 & 2.0 & 6.6 & 31.8 & 1.8 & 2.7\\
Petersianthus macrocarpus & 5 & 46 & 3.4 & 1.4 & 5.4 & 35.2 & 1.7 & 3.4\\
\addlinespace
Elaeis guineensis & 6 & 43 & 3.2 & 1.3 & 5.1 & 38.4 & 1.6 & 4.1\\
Uapaca guineensis & 7 & 38 & 2.8 & 1.0 & 4.6 & 41.2 & 1.6 & 4.7\\
Hylodendron gabonense & 8 & 28 & 2.1 & -0.2 & 4.3 & 43.2 & 1.4 & 5.4\\
Rinorea sp. & 9 & 27 & 2.0 & 0.3 & 3.7 & 45.3 & 1.4 & 6.1\\
Coelocaryon preussi & 10 & 23 & 1.7 & 0.5 & 2.9 & 47.0 & 1.4 & 6.8\\
\bottomrule
\end{tabular}
\begin{tablenotes}
\item \textit{Note: } 
\item rank: rang; abundance: abondance; plower: limite inferieure de confiance pour le rang de l'espèce; pupper: limite supérieure de confiance; accumfreq: fréquence cumulée; logabun: logarithme de l'abondance; rankfreq: fréquence du rang
\end{tablenotes}
\end{threeparttable}}
\end{table}

\begin{figure}

{\centering \includegraphics[width=0.8\linewidth]{03-resultats_et_discussions_files/figure-latex/Courbesabondance-Ndikinimeki3-1} 

}

\caption{Courbes d'abondance des espèces et familles dans la FC de Ndikiniméki.}\label{fig:Courbesabondance-Ndikinimeki3}
\end{figure}

\begin{table}
\centering
\caption{\label{tab:TableauabondancesNanga-Eboko3}Les 10 plus abondantes espèces de la FC de Nanga-Eboko.}
\centering
\resizebox{\ifdim\width>\linewidth\linewidth\else\width\fi}{!}{
\begin{threeparttable}
\begin{tabular}[t]{lrrrrrrrr}
\toprule
  & rank & abundance & proportion & plower & pupper & accumfreq & logabun & rankfreq\\
\midrule
Raphia sp & 1 & 89 & 10.6 & 0.0 & 21.1 & 10.6 & 1.9 & 0.8\\
Bosqueia angolensis & 2 & 49 & 5.8 & 2.9 & 8.8 & 16.4 & 1.7 & 1.6\\
Musanga cecropioides & 3 & 49 & 5.8 & 2.5 & 9.1 & 22.2 & 1.7 & 2.5\\
Monodora brevipes & 4 & 47 & 5.6 & 1.4 & 9.8 & 27.8 & 1.7 & 3.3\\
Uapaca guineensis & 5 & 37 & 4.4 & -0.1 & 8.9 & 32.1 & 1.6 & 4.1\\
\addlinespace
Sterculia rhinopetala & 6 & 34 & 4.0 & -0.1 & 8.1 & 36.2 & 1.5 & 4.9\\
Lannea welwitschii & 7 & 26 & 3.1 & 0.6 & 5.5 & 39.3 & 1.4 & 5.7\\
Trichilia tessmannii & 8 & 26 & 3.1 & 0.5 & 5.7 & 42.3 & 1.4 & 6.6\\
Piptadeniastrum africanum & 9 & 20 & 2.4 & 0.5 & 4.3 & 44.7 & 1.3 & 7.4\\
Beilschmiedia sp. & 10 & 19 & 2.3 & -0.6 & 5.1 & 47.0 & 1.3 & 8.2\\
\bottomrule
\end{tabular}
\begin{tablenotes}
\item \textit{Note: } 
\item rank: rang; abundance: abondance; plower: limite inferieure de confiance pour le rang de l'espèce; pupper: limite supérieure de confiance; accumfreq: fréquence cumulée; logabun: logarithme de l'abondance; rankfreq: fréquence du rang
\end{tablenotes}
\end{threeparttable}}
\end{table}

\begin{figure}

{\centering \includegraphics[width=0.8\linewidth]{03-resultats_et_discussions_files/figure-latex/CourbesabondanceNanga-Eboko3-1} 

}

\caption{Courbes d'abondance des espèces et familles dans la FC de Nanga-Eboko.}\label{fig:CourbesabondanceNanga-Eboko3}
\end{figure}

Afin de prédire la présence ou l'abondance des espèces en fonction des variables environnementales telles que le type d'utilisation des terres, l'altitude (\emph{elevation}), le stock de carbone des arbres et le carbone du sol, l'Abre de Regression Multivarié (ARM) a été d'une importance déterminante. Les statistiques à la base de l'ARM sont respectivement: l'erreur résiduel (\(Error\)), l'erreur due à la validation croisée (\(CV\ Error\)) et l'erreur standard (\(SE\)). A partir de la Figure:\ref{fig:ChoixpluspetitARM-Ndikinimeki3} nous constatons que la matrice des espèces dans la FC de Ndikiniméki a été partitionnée suivant les Types d'utilisation des terres (TUT) qui sont des catégories bien établies et distinctes. Cet arbre présente 2 groupes. Le premier regroupe les cacaoyères (CAC) et les jachères (JA) tandis que le second rassemble les forêts matures (FM), les forêts marécageuses (FMa) et les forêts secondaires (FS). Chaque groupe est caractérisé par un diagramme à bâtonnets qui indique l'abondance des espèces dans les groupes constitués, le nombre de parcelles (\(n\)) et l'erreur relative du groupe. \(L'erreur\ résiduel = 0,754\) ce qui signifit que le modèle R² (\(1-0,754\)) explique 24,6\% de la variance de la matrice des espèces.

\begin{verbatim}
#> Registered S3 methods overwritten by 'mvpart':
#>   method            from 
#>   labels.rpart      rpart
#>   model.frame.rpart rpart
#>   plot.rpart        rpart
#>   predict.rpart     rpart
#>   print.rpart       rpart
#>   residuals.rpart   rpart
#>   summary.rpart     rpart
#>   text.rpart        rpart
#> X-Val rep : 1  2  3  4  5  6  7  8  9  10  11  12  13  14  15  16  17  18  19  20  21  22  23  24  25  26  27  28  29  30  31  32  33  34  35  36  37  38  39  40  41  42  43  44  45  46  47  48  49  50  51  52  53  54  55  56  57  58  59  60  61  62  63  64  65  66  67  68  69  70  71  72  73  74  75  76  77  78  79  80  81  82  83  84  85  86  87  88  89  90  91  92  93  94  95  96  97  98  99  100
#> Minimum tree sizes
#> tabmins
#>   2 
#> 100
\end{verbatim}

\begin{figure}

{\centering \includegraphics[width=0.8\linewidth]{03-resultats_et_discussions_files/figure-latex/ChoixpluspetitARM-Ndikinimeki3-1} 

}

\caption{Plus petit arbre entre une erreur standard et l'erreur due à la validation croisée pour la FC de Ndikiniméki.}\label{fig:ChoixpluspetitARM-Ndikinimeki3}
\end{figure}

Le plus grand arbre de regression multivariée est difficile à interpreter. Du fait qu'il intègre au moins 4 éléments discriminants, il en découle un nombre important de groupes bien qu'il offre une puissance explicative élevée (R²=60\%) et superieur à celle de la solution à 2 groupes (Figure:\ref{fig:ChoixplusgrandARM-Ndikinimeki3}).

\begin{verbatim}
#> X-Val rep : 1  2  3  4  5  6  7  8  9  10  11  12  13  14  15  16  17  18  19  20  21  22  23  24  25  26  27  28  29  30  31  32  33  34  35  36  37  38  39  40  41  42  43  44  45  46  47  48  49  50  51  52  53  54  55  56  57  58  59  60  61  62  63  64  65  66  67  68  69  70  71  72  73  74  75  76  77  78  79  80  81  82  83  84  85  86  87  88  89  90  91  92  93  94  95  96  97  98  99  100
#> Minimum tree sizes
#> tabmins
#>   2 
#> 100
\end{verbatim}

\begin{figure}

{\centering \includegraphics[width=0.8\linewidth]{03-resultats_et_discussions_files/figure-latex/ChoixplusgrandARM-Ndikinimeki3-1} 

}

\caption{Plus grand arbre de regression multivariée pour la FC de Ndikiniméki.}\label{fig:ChoixplusgrandARM-Ndikinimeki3}
\end{figure}

Entre les deux solutions précédentes, une répartition en 4 groupes produit un graphe qui pourrait être interpreté comme des differences de composition en espèces à travers les groupes constitués, le R² de cet arbre (\(1-erreur\ relative\)) est 0,367 (Figure:\ref{fig:Choixarbreapproprie-Ndikinimeki3}). Le premier découpage s'est fait sur la base du type d'utilisation des sols (TUT): il a permis d'isoler un groupe de 10 parcelles dominé par les Cacaoyères (CAC) et un second groupe de 22 parcelles retrouvées dans les Forêts marécageuses (Fma), les Forêts matures (FM) et les Forêts secondaires (FS). Sur la base des stocks de carbone, le premier groupe de 10 parcelles a de nouveau été scindé en deux sous-groupes:le premier regroupe 7 parcelles établies dans les Cacaoyères (groupe 3) dont les stocks de carbone sont inférieurs à 20,85t/ha (Tableau:\ref{tab:Groupe3-Ndikinimeki}) et un second de 3 parcelles (groupe 4) éssentiellement dans les Jachères avec des stocks de carbone superieurs à 20,85 t/ha (Tableau:\ref{tab:Groupe4-Ndikinimeki}. Le bloc de 22 parcelles a permis d'isolé un sous-groupe de 2 parcelles (groupe 6) appartenant aux forêts marécageuses (Tableau:\ref{tab:Groupe6-Ndikinimeki}). Les 20 parcelles restantes forment un bloc (groupe 7) réparti entre les Forêts matures et les forêts secondaires (Tableau:\ref{tab:Groupe7-Ndikinimeki}).

\begin{verbatim}
#> X-Val rep : 1  2  3  4  5  6  7  8  9  10  11  12  13  14  15  16  17  18  19  20  21  22  23  24  25  26  27  28  29  30  31  32  33  34  35  36  37  38  39  40  41  42  43  44  45  46  47  48  49  50  51  52  53  54  55  56  57  58  59  60  61  62  63  64  65  66  67  68  69  70  71  72  73  74  75  76  77  78  79  80  81  82  83  84  85  86  87  88  89  90  91  92  93  94  95  96  97  98  99  100
#> Minimum tree sizes
#> tabmins
#>   2 
#> 100
\end{verbatim}

\begin{figure}

{\centering \includegraphics[width=0.8\linewidth]{03-resultats_et_discussions_files/figure-latex/Choixarbreapproprie-Ndikinimeki3-1} 

}

\caption{Arbre approprié pour une classification supervisée dans la FC de Ndikiniméki.}\label{fig:Choixarbreapproprie-Ndikinimeki3}
\end{figure}

\begin{table}
\centering
\caption{\label{tab:Groupe3-Ndikinimeki}Composition environnementale du groupe 3 dans la FC de Ndikiniméki.}
\centering
\resizebox{\ifdim\width>\linewidth\linewidth\else\width\fi}{!}{
\begin{threeparttable}
\begin{tabular}[t]{llrrr}
\toprule
  & TUT & carbone\_sol & carbone & elevation\\
\midrule
2 & CAC & 5.72 & 18.50 & 493.87569\\
3 & CAC & 6.36 & 8.31 & 483.93507\\
5 & CAC & 5.52 & 13.20 & 494.64359\\
9 & CAC & 6.70 & 15.80 & 517.48100\\
10 & CAC & 4.69 & 16.50 & 506.98081\\
\addlinespace
19 & CAC & NA & 9.98 & 766.63084\\
22 & CAC & NA & 15.90 & 739.12294\\
\bottomrule
\end{tabular}
\begin{tablenotes}
\item \textit{Note: } 
\item TUT: Type d'utilisation des terres; carbone\_sol: stocks de carbone du sol (tC/ha); carbone: stock de carbone de ligneux (tC/ha); elevation: altitude; CAC; Cacaoyères; NA: Valeur manquante.
\end{tablenotes}
\end{threeparttable}}
\end{table}

\begin{table}
\centering
\caption{\label{tab:Groupe4-Ndikinimeki}Composition environnementale du groupe 4 dans la FC de Ndikiniméki.}
\centering
\resizebox{\ifdim\width>\linewidth\linewidth\else\width\fi}{!}{
\begin{threeparttable}
\begin{tabular}[t]{llrrr}
\toprule
  & TUT & carbone\_sol & carbone & elevation\\
\midrule
15 & JA & 4.43 & 118.0 & 744.89082\\
18 & JA & 3.76 & 23.2 & 803.40607\\
26 & CAC & NA & 58.6 & 734.87373\\
\bottomrule
\end{tabular}
\begin{tablenotes}
\item \textit{Note: } 
\item TUT: Type d'utilisation des terres; carbone\_sol: stocks de carbone du sol (tC/ha); carbone: stock de carbone de ligneux (tC/ha); elevation: altitude; JA: Jachères; CAC: Cacaoyères; NA: Valeur manquante.
\end{tablenotes}
\end{threeparttable}}
\end{table}

\begin{table}
\centering
\caption{\label{tab:Groupe6-Ndikinimeki}Composition environnementale du groupe 6 dans la FC de Ndikiniméki.}
\centering
\resizebox{\ifdim\width>\linewidth\linewidth\else\width\fi}{!}{
\begin{threeparttable}
\begin{tabular}[t]{llrrr}
\toprule
  & TUT & carbone\_sol & carbone & elevation\\
\midrule
17 & FMa & 4.52 & 77.8 & 782.93816\\
27 & FMa & 10.10 & 52.1 & 490.23347\\
\bottomrule
\end{tabular}
\begin{tablenotes}
\item \textit{Note: } 
\item TUT: Type d'utilisation des terres; carbone\_sol: stocks de carbone du sol (tC/ha); carbone: stock de carbone de ligneux (tC/ha); elevation: altitude; FMa: Forêts marécageuses.
\end{tablenotes}
\end{threeparttable}}
\end{table}

\begin{table}
\centering
\caption{\label{tab:Groupe7-Ndikinimeki}Composition environnementale du groupe 7 dans la FC de Ndikiniméki.}
\centering
\resizebox{\ifdim\width>\linewidth\linewidth\else\width\fi}{!}{
\begin{threeparttable}
\begin{tabular}[t]{llrrr}
\toprule
  & TUT & carbone\_sol & carbone & elevation\\
\midrule
1 & FM & 6.78 & 22.3 & 509.85335\\
4 & FS & 6.66 & 48.0 & 504.75392\\
6 & FS & 4.94 & 25.7 & 564.92455\\
7 & FS & NA & 22.6 & 635.51778\\
8 & FM & 10.10 & 20.2 & 706.91655\\
\addlinespace
11 & FM & 6.44 & 17.2 & 558.23267\\
12 & FM & 5.42 & 19.6 & 505.22692\\
13 & FS & 5.50 & 41.6 & 500.72012\\
14 & FM & 3.91 & 11.4 & 738.59781\\
16 & FS & 5.38 & 26.1 & 803.69778\\
\addlinespace
20 & FS & 6.97 & 36.0 & 831.12353\\
21 & FS & NA & 51.9 & 818.22151\\
23 & FM & 4.56 & 27.5 & 491.68903\\
24 & FM & 6.14 & 39.2 & 727.15077\\
25 & FM & NA & 39.4 & 770.75311\\
\addlinespace
28 & FM & 5.26 & 16.5 & 510.65417\\
29 & FM & 6.38 & 31.4 & 544.31595\\
30 & FM & NA & 29.6 & 489.74141\\
31 & FM & NA & 21.0 & 511.10347\\
32 & FM & NA & 41.9 & 493.47825\\
\bottomrule
\end{tabular}
\begin{tablenotes}
\item \textit{Note: } 
\item TUT: Type d'utilisation des terres; carbone\_sol: stocks de carbone du sol (tC/ha); carbone: stock de carbone de ligneux (tC/ha); elevation: altitude; NA: Valeur manquante; FM: Forêts matures; FS: Forêts secondaires.
\end{tablenotes}
\end{threeparttable}}
\end{table}

Parmi les 148 espèces identifiées dans la FC de Ndikiniméki, 7 ont été retenues comme représentives ou caractéristiques de certains groupes de parcelles appartenant à des Types t'Utilisation de Terres (TUT).Parmi ces espèces selectionnées, 3 sont exclusivement associées à un groupe, 4 sont associées à 2 groupes et aucune espèce n'a été associée à 3 groupes à la fois:

Indicateurs des groupes identifiés liés aux variables explicatives dans la FC de Ndikiniméki

\begin{verbatim}
#> 
#>  Multilevel pattern analysis
#>  ---------------------------
#> 
#>  Association function: IndVal.g
#>  Significance level (alpha): 0.05
#> 
#>  Total number of species: 148
#>  Selected number of species: 7 
#>  Number of species associated to 1 group: 3 
#>  Number of species associated to 2 groups: 4 
#>  Number of species associated to 3 groups: 0 
#> 
#>  List of species associated to each combination: 
#> 
#>  Group 4  #sps.  1 
#>                             A      B  stat p.value  
#> Klainedoxa gabonensis  0.8696 0.6667 0.761   0.033 *
#> 
#>  Group 6  #sps.  2 
#>                         A      B  stat p.value    
#> Nauclea pobeguinii 1.0000 1.0000 1.000   0.001 ***
#> Vernonia conferta  0.9545 1.0000 0.977   0.001 ***
#> 
#>  Group 3+4  #sps.  3 
#>                       A      B  stat p.value    
#> Theobroma cacao  0.9451 1.0000 0.972   0.001 ***
#> Musa sp.         1.0000 0.8000 0.894   0.009 ** 
#> Dacryodes edulis 1.0000 0.6000 0.775   0.035 *  
#> 
#>  Group 4+7  #sps.  1 
#>                            A      B  stat p.value  
#> Pycnanthus angolensis 0.7607 1.0000 0.872   0.018 *
#> ---
#> Signif. codes:  0 '***' 0.001 '**' 0.01 '*' 0.05 '.' 0.1 ' ' 1
\end{verbatim}

L'espèce \emph{Klainedoxa gabonensis} (\(p.value=0,027\)) s'est revelé comme un bon indicateur du groupe 4 (Jachères) avec une composante spécifique élevée (\(A=0,8696\)) et une fidélité au dessus de la moyenne (\(B=0,6667\)). Les espèces telles que: \emph{Nauclea pobeguinii}(\(A=1,0000\ et\ B=1,0000\)) et \emph{Vernonia conferta} (\(A=0,9545\ et\ B=1,0000\)) sont de très bons indicateurs (\(p.value=0,002\)) du groupe 6 (Forêts marécageuses) parce qu'elles apparaissent dans toutes les parcelles qui constituent ce groupe (\(B=1,0000\)).Avec une composante de spécificité \(A=1,0000\), ces espèces ne sont rencontrées presque exclusivement que dans les Forêts marécageuses. les espèces telles que: \emph{Theobroma cacao}, \emph{Musa sp.} et \emph{Dacryodes edulis} sont des indicateurs des groupe 3+4 (Cacaoyères et Jachères). Dans cette association de deux groupes constituée de 10 parcelles au total, \emph{Theobroma} est rencontrée dans presque tous les sites et nulle part ailleur (\(A=0,9451\ et\ B=1,000\)) cependant \emph{Musa sp.} et \emph{Dacryodes edulis} ne sont pas présentent dans toutes les parcelles qui constituent ce regroupement. Une association des groupes 4 et 7 (Jachères, Forêts matures et Forêts secondaires) a pour indicateur l'espèce \emph{Pycnanthus angolensis}. Cette espèce est présente dans toutes les parcelles constituant ce regroupement (\(B=1,0000\)) mais cependant on pourrait la retrouver dans des parcelles en dehors de ce groupe (\(A=0,7606\)).

Dans le cas de la FC de Nanga-Eboko, lorsque nous observons le plus petit arbre qui se situe dans les limites de l'erreur standard, nous constatons que cet modèle comporte 3 groupes. Le premier regroupe les cultures (CUL) et les savanes (SAV), le second représente les marécages (MAR) tandis que le troisième rassemble les forêts galeries (FGA), les forêts secondaires adultes (FSA) et les forêts secondaires jeunes (FSJ). Son erreur résiduel est égale à 0,568 et signifit que ce modèle pourrait expliquer 43,7\% (\(1-0,568\)) de la variance de la matrice des espèces (Figure:\ref{fig:ChoixpluspetitARMNanga-Eboko3}).

\begin{verbatim}
#> X-Val rep : 1  2  3  4  5  6  7  8  9  10  11  12  13  14  15  16  17  18  19  20  21  22  23  24  25  26  27  28  29  30  31  32  33  34  35  36  37  38  39  40  41  42  43  44  45  46  47  48  49  50  51  52  53  54  55  56  57  58  59  60  61  62  63  64  65  66  67  68  69  70  71  72  73  74  75  76  77  78  79  80  81  82  83  84  85  86  87  88  89  90  91  92  93  94  95  96  97  98  99  100
#> Minimum tree sizes
#> tabmins
#>   3 
#> 100
\end{verbatim}

\begin{figure}

{\centering \includegraphics[width=0.8\linewidth]{03-resultats_et_discussions_files/figure-latex/ChoixpluspetitARMNanga-Eboko3-1} 

}

\caption{Plus petit arbre entre une erreur standard et l'erreur due à la validation croisée pour la FC de Nanga-Eboko.}\label{fig:ChoixpluspetitARMNanga-Eboko3}
\end{figure}

Le modèle d'arbre le plus grand est celui à 7 groupes. Il a une puissance explicative élevée (\(R²=0,661\)), cependant la plurarité des facteurs explicatifs rend difficile son interpretation (Figure:\ref{fig:ChoixplusgrandARMNanga-Eboko3}).

\begin{verbatim}
#> X-Val rep : 1  2  3  4  5  6  7  8  9  10  11  12  13  14  15  16  17  18  19  20  21  22  23  24  25  26  27  28  29  30  31  32  33  34  35  36  37  38  39  40  41  42  43  44  45  46  47  48  49  50  51  52  53  54  55  56  57  58  59  60  61  62  63  64  65  66  67  68  69  70  71  72  73  74  75  76  77  78  79  80  81  82  83  84  85  86  87  88  89  90  91  92  93  94  95  96  97  98  99  100
#> Minimum tree sizes
#> tabmins
#>   3 
#> 100
\end{verbatim}

\begin{figure}

{\centering \includegraphics[width=0.8\linewidth]{03-resultats_et_discussions_files/figure-latex/ChoixplusgrandARMNanga-Eboko3-1} 

}

\caption{Plus grand arbre de regression multivariée pour la FC de Nanga-Eboko.}\label{fig:ChoixplusgrandARMNanga-Eboko3}
\end{figure}

Un modèle à 4 groupes intermediaire entre le minimum et le maximum de groupes possibles semble plus facile à interpreter. Sa puissance explicative est de \(R²=0,502\) ce qui suggère qu'au moins 50\% de la variance de la matrice des espèces pourrait trouver une justification dans ce modèle (Figure:\ref{fig:ChoixarbreapproprieNanga-Eboko3}). La première discrimination est faite sur la base des catégories d'utilisation des terres et permet d'isoler: un groupe de 5 parcelles (Tableau:\ref{tab:Groupe2-Nanga}) appartenant aux Cultures (CUL) et aux Savanes (SAV) et un second groupe de 4 parcelles (Tableau:\ref{tab:Groupe4-Nanga}) dans les Marécages (MAR) et un troisième groupe de 14 parcelles dans les Forêts galeries (FGA), les Forêts secondaires adultes (FSA) et les Forêts secondaires jeunes (FSJ). Ce dernier groupe sur la base d'un autre paramètre environnemental tel que les stocks de carbone dans le sol, a été découpé en deux sous-groupes: le premier regroupe 3 parcelles dont les stocks de carbone du sol sont inférieurs à 3,795t/ha (Tableau:\ref{tab:Groupe6-Nanga}) et le second est constitué de 11 parcelles dont les stocks de carbone du sol sont supérieurs ou égales à 3,795t/ha (Tableau:\ref{tab:Groupe7-Nanga}).

\begin{verbatim}
#> X-Val rep : 1  2  3  4  5  6  7  8  9  10  11  12  13  14  15  16  17  18  19  20  21  22  23  24  25  26  27  28  29  30  31  32  33  34  35  36  37  38  39  40  41  42  43  44  45  46  47  48  49  50  51  52  53  54  55  56  57  58  59  60  61  62  63  64  65  66  67  68  69  70  71  72  73  74  75  76  77  78  79  80  81  82  83  84  85  86  87  88  89  90  91  92  93  94  95  96  97  98  99  100
#> Minimum tree sizes
#> tabmins
#>   3 
#> 100
\end{verbatim}

\begin{figure}

{\centering \includegraphics[width=0.8\linewidth]{03-resultats_et_discussions_files/figure-latex/ChoixarbreapproprieNanga-Eboko3-1} 

}

\caption{Arbre approprié pour une classification supervisée dans la FC de Nanga-Eboko.}\label{fig:ChoixarbreapproprieNanga-Eboko3}
\end{figure}

\begin{table}
\centering
\caption{\label{tab:Groupe2-Nanga}Composition environnementale du groupe 2 dans la FC de Nanga-Eboko.}
\centering
\resizebox{\ifdim\width>\linewidth\linewidth\else\width\fi}{!}{
\begin{threeparttable}
\begin{tabular}[t]{llrrr}
\toprule
  & TUT & Carbone\_sol & Carbone & Elevation\\
\midrule
3 & CUL & 2.9971985 & 8.61 & 643.49524\\
4 & SAV & 4.3988540 & 13.80 & 636.91013\\
5 & SAV & 3.2247951 & 23.90 & 617.23598\\
16 & SAV & 6.8328961 & 16.30 & 641.93456\\
17 & CUL & NA & 6.20 & 658.38135\\
\bottomrule
\end{tabular}
\begin{tablenotes}
\item \textit{Note: } 
\item TUT: Type d'utilisation des terres; Carbone\_sol: stocks de carbone du sol (tC/ha); Carbone: stock de carbone de ligneux (tC/ha); Elevation: altitude; CUL: Cultures; SAV: Savanes; NA: Valeur manquante.
\end{tablenotes}
\end{threeparttable}}
\end{table}

\begin{table}
\centering
\caption{\label{tab:Groupe4-Nanga}Composition environnementale du groupe 4 dans la FC de Nanga-Eboko.}
\centering
\resizebox{\ifdim\width>\linewidth\linewidth\else\width\fi}{!}{
\begin{threeparttable}
\begin{tabular}[t]{llrrr}
\toprule
  & TUT & Carbone\_sol & Carbone & Elevation\\
\midrule
9 & MAR & 5.0195438 & 11.1 & 612.84773\\
10 & MAR & 2.9388678 & 18.4 & 599.46846\\
12 & MAR & 4.7552576 & 40.5 & 617.62700\\
22 & MAR & 1.7002256 & 14.6 & 623.66549\\
\bottomrule
\end{tabular}
\begin{tablenotes}
\item \textit{Note: } 
\item TUT: Type d'utilisation des terres; Carbone\_sol: stocks de carbone du sol (tC/ha); Carbone: stock de carbone de ligneux (tC/ha); Elevation: altitude; MAR: Marécages.
\end{tablenotes}
\end{threeparttable}}
\end{table}

\begin{table}
\centering
\caption{\label{tab:Groupe6-Nanga}Composition environnementale du groupe 6 dans la FC de Nanga-Eboko.}
\centering
\resizebox{\ifdim\width>\linewidth\linewidth\else\width\fi}{!}{
\begin{threeparttable}
\begin{tabular}[t]{llrrr}
\toprule
  & TUT & Carbone\_sol & Carbone & Elevation\\
\midrule
7 & FGA & 3.6746280 & 75.4 & 613.10142\\
18 & FSA & 3.4476684 & 35.0 & 606.53717\\
20 & FSA & 3.5581825 & 63.9 & 618.81868\\
\bottomrule
\end{tabular}
\begin{tablenotes}
\item \textit{Note: } 
\item TUT: Type d'utilisation des terres; Carbone\_sol: stocks de carbone du sol (tC/ha); Carbone: stock de carbone de ligneux (tC/ha); Elevation: altitude; FGA: Forêts Galeries; FSA: Forêts Secondaires Adultes.
\end{tablenotes}
\end{threeparttable}}
\end{table}

\begin{table}
\centering
\caption{\label{tab:Groupe7-Nanga}Composition environnementale du groupe 7 dans la FC de Nanga-Eboko.}
\centering
\resizebox{\ifdim\width>\linewidth\linewidth\else\width\fi}{!}{
\begin{threeparttable}
\begin{tabular}[t]{llrrr}
\toprule
  & TUT & Carbone\_sol & Carbone & Elevation\\
\midrule
1 & FSA & 4.3658873 & 42.6 & 608.76556\\
2 & FSA & 6.0144750 & 65.2 & 629.89594\\
6 & FGA & 5.3303448 & 57.8 & 620.93674\\
8 & FSA & 4.2106535 & 22.8 & 644.97739\\
11 & FSA & 7.2681204 & 22.8 & 639.67198\\
\addlinespace
13 & FSA & 4.5925335 & 59.6 & 649.85674\\
14 & FGA & 5.0794911 & 69.8 & 592.46553\\
15 & FSJ & 6.1766307 & 52.7 & 625.65163\\
19 & FSJ & 4.8375422 & 18.4 & 622.03717\\
21 & FSJ & 3.9155740 & 47.3 & 619.59523\\
\addlinespace
23 & FSJ & 5.1846006 & 70.8 & 638.45004\\
\bottomrule
\end{tabular}
\begin{tablenotes}
\item \textit{Note: } 
\item TUT: Type d'utilisation des terres; Carbone\_sol: stocks de carbone du sol (tC/ha); Carbone: stock de carbone de ligneux (tC/ha); Elevation: altitude; FSA: Forêts Secondaires Adultes; FGA: Forêts Galeries; FSJ: Forêts Secondaires Jeunes.
\end{tablenotes}
\end{threeparttable}}
\end{table}

D'après le modèle d'analyse de la composition spécifique des différents groupes constitués dans la FC de Nanga-Eboko, il ressort des 122 espèces identifiées que, 12 ont été associées à la représentation d'un groupe. 3 espèces ont été identifiées comme caractéristiques de 2 groupes et aucune espèce n'a été associée à la caractérisation de 3 groupes:

Indicateurs des groupes identifiés liés aux variables explicatives dans la FC de Nanga-Eboko

\begin{verbatim}
#> 
#>  Multilevel pattern analysis
#>  ---------------------------
#> 
#>  Association function: IndVal.g
#>  Significance level (alpha): 0.05
#> 
#>  Total number of species: 122
#>  Selected number of species: 15 
#>  Number of species associated to 1 group: 13 
#>  Number of species associated to 2 groups: 2 
#>  Number of species associated to 3 groups: 0 
#> 
#>  List of species associated to each combination: 
#> 
#>  Group 2  #sps.  3 
#>                          A   B  stat p.value    
#> Annona senegalensis    1.0 1.0 1.000   0.001 ***
#> Terminalia glaucescens 1.0 0.8 0.894   0.001 ***
#> Inconnu                1.0 0.6 0.775   0.019 *  
#> 
#>  Group 4  #sps.  3 
#>                           A      B  stat p.value    
#> Raphia sp            1.0000 1.0000 1.000   0.001 ***
#> Berlinia grandiflora 1.0000 0.7500 0.866   0.002 ** 
#> cola sp.             0.9429 0.5000 0.687   0.046 *  
#> 
#>  Group 6  #sps.  6 
#>                        A      B  stat p.value   
#> Beilschmiedia sp. 0.9322 1.0000 0.966   0.004 **
#> Uapaca guineensis 0.8627 1.0000 0.929   0.003 **
#> Milicia excelsa   0.6197 1.0000 0.787   0.043 * 
#> Afzelia pachyloba 0.9167 0.6667 0.782   0.017 * 
#> Sapium ellipticum 0.9167 0.6667 0.782   0.012 * 
#> Ceiba pentandra   0.5867 1.0000 0.766   0.032 * 
#> 
#>  Group 7  #sps.  1 
#>                           A      B  stat p.value  
#> Trichilia tessmannii 1.0000 0.6364 0.798   0.021 *
#> 
#>  Group 6+7  #sps.  2 
#>                          A      B  stat p.value   
#> Bosqueia angolensis 0.9571 0.8571 0.906   0.002 **
#> Lannea welwitschii  1.0000 0.6429 0.802   0.040 * 
#> ---
#> Signif. codes:  0 '***' 0.001 '**' 0.01 '*' 0.05 '.' 0.1 ' ' 1
\end{verbatim}

L'espèce \emph{Annona senegalensis} s'est revelée comme un très bon indicateur du \enquote{groupe 2} constitué éssentiellement des savanes (\(p.value=0,001\)) suivie de \emph{Terminalia glaucescens} (\(p.value=0,003\)). Le \enquote{groupe 4} est constitué de 4 parcelles appartenant aux zones de marécage. Il est marqué par 4 espèces dont l'une d'elle \emph{Raphia sp.} (\(p.value=0,001\)) est fortement liée aux conditions écologiques de ce milieu. Les autres espèces caractéristiques de ce groupe sont: \emph{Berlinia grandiflora} (\(p.value=0,005\)) et \emph{Cola sp.} (\(p.value=0,040\)). Le \enquote{groupe 6} est marqué par les stocks de carbone du sol inferiers à 3,7 tC/ha localisés dans les forêts secondaires adultes (FSA) et les forêts galeries (FGA). Il est marqué par 6 espèces dont les plus caractéristiques sont \emph{Beilschmiedia sp.} (\(p.value=0,005\)) et \emph{Uapaca guineensis} (\(p.value=0,005\)). le \enquote{groupe 7} est constitué de 11 parcelles dont les stocks de carbone du sol sont supériers ou égale à 3,7 tC/ha. L'espèce \emph{Tchichilia tessmanni} (\(p.value=0,027\)) peut être considéré comme un indicateur de ce groupe. Dans le cas où l'on associerait les groupes 6 et 7, l'espèce \emph{Bosqueia angolensis} serait un très bon indicateur à la fois pour les forêts secondaires adultes, forêts secondaires jeunes et les forêts galeries (\(p.value=0,001\)).

\hypertarget{estmation-des-stocks-de-carbone}{%
\subsection{Estmation des stocks de carbone}\label{estmation-des-stocks-de-carbone}}

\hypertarget{estimation-des-stocks-de-carbone-des-ligneux}{%
\subsubsection{Estimation des stocks de carbone des ligneux}\label{estimation-des-stocks-de-carbone-des-ligneux}}

L'estimation des stocks de carbone des ligneux à partir des mesures dendrométriques a portée sur :

\begin{itemize}
\tightlist
\item
  1228 individus dans la FC de Ndikiniméki dont les diamètres varient de 5 à 120 cm pour une moyenne de 33,16 cm dont 51\% ont un diamètre inférieur à 40 cm:
\end{itemize}

\begin{verbatim}
#>    Strates           diametre_cl  StC_total_arbre.t.C.ha
#>  Length:1228        [30,40):235   Min.   : 0.011717     
#>  Class :character   [5,10) :213   1st Qu.: 0.155157     
#>  Mode  :character   [10,15):187   Median : 0.366378     
#>                     [15,20):109   Mean   : 0.827977     
#>                     [40,50): 99   3rd Qu.: 1.006809     
#>                     [60,70): 94   Max.   :18.773502     
#>                     (Other):291
\end{verbatim}

\begin{itemize}
\tightlist
\item
  763 individus dans la FC de Nanga-Eboko dont les diamètres varient de 5 à 185 cm pour une moyenne de 30,46 cm et 61\% ont un diamètre inférieur à 40 cm:
\end{itemize}

\begin{verbatim}
#>    Strates           diametre_cl  StC_total_arbre.t.C.ha
#>  Length:763         [30,40):230   Min.   : 0.014918     
#>  Class :character   [5,10) :148   1st Qu.: 0.262412     
#>  Mode  :character   [10,15): 88   Median : 0.530625     
#>                     [15,20): 73   Mean   : 1.124347     
#>                     [40,50): 67   3rd Qu.: 1.219292     
#>                     [50,60): 34   Max.   :17.286403     
#>                     (Other):123
\end{verbatim}

La distribution des stocks de carbone entre les individus des différents Types d'Utilisation de Terres dans la FC de Ndikiniméki, varie entre 0,01 et 18,77 tC/ha pour une moyenne de 0,82 ± 0,04 tC/ha. Les Forêts matures (FM) occupent le haut de l'échelle avec un total cumulé de 337,1 ± 0,60 tC/ha pour une moyenne de 0,57 tC/ha suivit des Forêts secondaires (FS) 251,96 ± 1,01 tC/ha cumulés dans l'ensemble pour une moyenne de 0,91 tC/ha. A la suite des Cacaoyères (CAC) et des Jachères (JA), les Forêts marécageuses (FMa) ferment la queue avec un total cumulé de 129,88 ± 2,75 tC/ha pour une moyenne de 2,54 tC/ha (Tableau:\ref{tab:Stockscarboneligneux-Ndikinimeki3}). L'estimation total des stocks de carbone dans l'ensemble des 5 catégories d'utilisation des terres dans la FC de Ndikiniméki s'élève à 1016,76 ± 89,02 tC/ha pour une moyenne de 203,35 tC/ha (Figure:\ref{fig:Barplotcarboneligneux-Ndikinimeki3}). Bien que les individus de diamètre inferieur à 40 cm constituent plus de la moitié de l'effectif échantillonné, environ 54\% des stocks de carbone des ligneux proviennent des individus de diamètre compris entre 60 et 100 cm de diamètre (environ 30\% des effectifs) Tableau:\ref{tab:Stockscarboneclassesdiametre-Ndikinimeki3} et cette tendance s'observe dans toutes les catégories d'utilisation des terres dans la FC de Ndikiniméki (Figures:\ref{fig:Carboneligneux-Ndiki3} et \ref{fig:Carboneligneux-Ndikinimeki3}).

\begin{table}
\centering
\caption{\label{tab:Stockscarboneligneux-Ndikinimeki3}Stocks de carbone des ligneux dans les différentes catégories d'utilisation de terres dans la FC de Ndikiniméki.}
\centering
\resizebox{\ifdim\width>\linewidth\linewidth\else\width\fi}{!}{
\begin{threeparttable}
\begin{tabular}[t]{lr}
\toprule
Strates & carbone.total\\
\midrule
FM & 337.10\\
FS & 251.96\\
CAC & 156.73\\
JA & 141.09\\
FMa & 129.88\\
\bottomrule
\end{tabular}
\begin{tablenotes}
\item \textit{Note: } 
\item length: taille de l'échantillon; mean: moyenne; sd: écart type; range: amplitude; FM: Forêt Mature; FS: Forêt Secondaire; FMa: Forêt Marécageuse; JA: Jachère; CAC: Cacaoyère; carbone\_total:stock de carbone total en tC/ha.
\end{tablenotes}
\end{threeparttable}}
\end{table}

\begin{table}
\centering
\caption{\label{tab:Stockscarboneclassesdiametre-Ndikinimeki3}Répartion des stocks de carbone suivant les classes de diamètre dans la FC de Ndikiniméki.}
\centering
\resizebox{\ifdim\width>\linewidth\linewidth\else\width\fi}{!}{
\begin{threeparttable}
\begin{tabular}[t]{lr}
\toprule
diametre\_cl & carbone.total\\
\midrule
{}[80,90) & 197.31\\
{}[60,70) & 145.61\\
{}[70,80) & 117.50\\
{}[90,100) & 87.31\\
{}[30,40) & 86.94\\
\addlinespace
{}[40,50) & 63.76\\
{}[15,20) & 60.66\\
{}[25,30) & 57.05\\
{}[50,60) & 49.16\\
{}[20,25) & 47.64\\
\addlinespace
{}[10,15) & 42.26\\
{}[100,110) & 32.43\\
{}[110,120) & 14.98\\
{}[5,10) & 11.91\\
{}[120,130) & 2.23\\
\bottomrule
\end{tabular}
\begin{tablenotes}
\item \textit{Note: } 
\item diametre\_cl: classe de diamètre; carbone.total: stock de carbone total en tC/ha.
\end{tablenotes}
\end{threeparttable}}
\end{table}

\begin{figure}

{\centering \includegraphics[width=0.8\linewidth]{03-resultats_et_discussions_files/figure-latex/Barplotcarboneligneux-Ndikinimeki3-1} 

}

\caption{Répartition des stocks de carbone des ligneux dans les TUT de la FC de Ndikiniméki.}\label{fig:Barplotcarboneligneux-Ndikinimeki3}
\end{figure}

\begin{figure}

{\centering \includegraphics[width=0.8\linewidth]{03-resultats_et_discussions_files/figure-latex/Carboneligneux-Ndiki3-1} 

}

\caption{Répartition des stocks de carbone suivant les classes de diamètre dans la FC de Ndikiniméki.}\label{fig:Carboneligneux-Ndiki3}
\end{figure}

\begin{figure}

{\centering \includegraphics[width=0.8\linewidth]{03-resultats_et_discussions_files/figure-latex/Carboneligneux-Ndikinimeki3-1} 

}

\caption{Répartition des stocks de carbone suivant les classes de diamètre et les TUT dans la FC de Ndikiniméki.}\label{fig:Carboneligneux-Ndikinimeki3}
\end{figure}

L'analyse de la distribution des stocks de carbone dans les 5 catégories d'utilisation de terres dans la FC de Ndikiniméki revèle dans l'ensemble, des différences significatives dans la distribution des stocks de carbone des ligneux. Seules les Forêts marécageuses (FMa)-Jachères (JA) d'une part et les Forêts matures (FM)-Cacaoyères (CAC) d'autre part ne présentent pas de différences significatives dans la répartition des stocks de carbone (Figure:\ref{fig:ttestligneux-Ndiki3}). Les forêts marécageuses et les jachères enregistrent un faible nombre d'observation et des écarts types particulièrement élevés (Tableau:\ref{tab:Stocksmoyensligneux-Ndiki3}). Cela peut se justifier d'une part, par le faible nombre de parcelles échantillonnées pour ces deux catégories (2 parcelles chacune) et d'autre part, par une grande dispersion des données (écart type élevé). En effet, c'est dans ces parcelles qu'ont été enregistrées les espèces ayant les plus fortes quantités de carbone de toute la distribution. De même les amplitudes de distribution des stocks de carbone de ces deux catégories sont particulièrement grandes témoignant ainsi, un écart important entre les valeurs de stock les plus faibles et les valeurs les plus élevées observées dans toute la distribution des stocks de carbone des ligneux. Par exemple dans les jachères nous observons un effectif de 54 individus pour une amplitude d'environ 18,73 tC/ha tandis les Forêts matures enregistrent un effectif de 586 individus pour une amplitude de 4,03 tC/ha. Ces résultats montrent qu'une perturbation du milieu dans les jachères, notamment des activités agricoles pourraient justifier cet état de fait. Par contre, dans les forêts marécageuses, cette situation pourrait s'expliquer par le faible taux d'échantillonnage couplé à la faible densité des espèces liée aux conditions édaphiques.

\begin{figure}

{\centering \includegraphics[width=0.8\linewidth]{03-resultats_et_discussions_files/figure-latex/ttestligneux-Ndiki3-1} 

}

\caption{t.test des stocks de carbone des ligneux dans la FC de Ndikiniméki.}\label{fig:ttestligneux-Ndiki3}
\end{figure}

\begin{table}
\centering
\caption{\label{tab:Stocksmoyensligneux-Ndiki3}Calcul des stocks moyens de carbone des ligneux dans les TUT de la FC de Ndikiniméki.}
\centering
\resizebox{\ifdim\width>\linewidth\linewidth\else\width\fi}{!}{
\begin{threeparttable}
\begin{tabular}[t]{lrrrr}
\toprule
Strates & length & mean & sd & range\\
\midrule
FM & 586 & 0.58 & 0.60 & 4.04\\
FS & 276 & 0.91 & 1.02 & 6.72\\
CAC & 261 & 0.60 & 1.06 & 10.11\\
JA & 54 & 2.61 & 4.44 & 18.74\\
FMa & 51 & 2.55 & 2.75 & 12.14\\
\bottomrule
\end{tabular}
\begin{tablenotes}
\item \textit{Note: } 
\item length: taille de l'échantillon; mean: moyenne; sd: écart type; range: amplitude; FM: Forêt Mature; FS: Forêt Secondaire; FMa: Forêt Marécageuse; JA: Jachère; CAC: Cacaoyère
\end{tablenotes}
\end{threeparttable}}
\end{table}

La distribution des stocks de carbone des ligneux entre les individus des différents Types d'Utilisation de Terres dans la FC de Nanga-Eboko, varie entre 0,01 et 17,28 tC/ha pour une moyenne de 1,12 ± 1,17 tC/ha. Les Forêts secondaires adultes (FSA) occupent le haut du classement avec un total cumulé de 312,06 ± 1,26 tC/ha pour une moyenne de 0,83 tC/ha suivit des Forêts galeries (FGA) 203,11 ± 2,51 tC/ha cumulés dans l'ensemble pour une moyenne de 1,84 tC/ha. A la suite, nous avons les Forêts secondaires jeunes (FSJ) qui cumulent 189,21 ± 1,52 tC/ha pour une moyenne de 1,03 tC/ha, les zones marécageuses (MAR) totalisent 84,65 ± 2,58 tC/ha avec une moyenne de 1,41 tC/ha, les Savanes quant à elles cumulent 54,04 ± 1,65 tC/ha pour une moyenne de 1,68 tC/ha. Les Cultures (CUL) ferment la queue du classement avec un total cumulé de 14,81 ± 3,22 tC/ha pour une moyenne de 2,96 tC/ha (Tableau:\ref{tab:Stockscarboneligneux-Nanga3}). L'estimation total des stocks de carbone des ligneux dans l'ensemble des 6 catégories d'occupation des terres dans la FC de Nanga-Eboko s'élève à 857,88 ± 111,42 tC/ha pour une moyenne de 142,98 tC/ha (Figure:\ref{fig:Barplotcarboneligneux-Nanga3}). Contrairement à la FC de Ndikiniméki, l'éssentiel des stocks de carbone des ligneux est localisé chez les individus de diamètre inférieur à 40 cm (environ 45\% des stocks enregistrés) tandis que les individus de diamètre compris entre 60 et 100 cm représentent 30\% des stocks enregistrés (Tableau:\ref{tab:Stockscarboneclassesdiametre-Nanga3}).

\begin{table}
\centering
\caption{\label{tab:Stockscarboneligneux-Nanga3}Stocks de carbone des ligneux dans les différentes catégories d'utilisation de terres dans la FC de Nanga-Eboko.}
\centering
\resizebox{\ifdim\width>\linewidth\linewidth\else\width\fi}{!}{
\begin{threeparttable}
\begin{tabular}[t]{lr}
\toprule
Strates & carbone\_total\\
\midrule
FSA & 312.06\\
FGA & 203.11\\
FSJ & 189.21\\
MAR & 84.65\\
SAV & 54.04\\
\addlinespace
CUL & 14.81\\
\bottomrule
\end{tabular}
\begin{tablenotes}
\item \textit{Note: } 
\item FSA: Forêts Secondaires Adultes; FGA: Forêts Galeries; FSJ: Forêts Secondaires Jeunes; MAR: Marecages; SAV: Savanes; CUL: Cultures; carbone\_total: stock de carbone total en tC/ha
\end{tablenotes}
\end{threeparttable}}
\end{table}

\begin{table}
\centering
\caption{\label{tab:Stockscarboneclassesdiametre-Nanga3}Répartion des stocks de carbone suivant les classes de diamètre dans la FC de Nanga-Eboko.}
\centering
\resizebox{\ifdim\width>\linewidth\linewidth\else\width\fi}{!}{
\begin{threeparttable}
\begin{tabular}[t]{lr}
\toprule
diametre\_cl & carbone\_total\\
\midrule
{}[30,40) & 120.36\\
{}[15,20) & 86.68\\
{}[60,70) & 72.74\\
{}[80,90) & 72.52\\
{}[70,80) & 69.31\\
\addlinespace
{}[25,30) & 58.93\\
{}[20,25) & 58.83\\
{}[50,60) & 56.12\\
{}[40,50) & 53.87\\
{}[10,15) & 44.78\\
\addlinespace
{}[90,100) & 42.47\\
{}[110,120) & 37.43\\
{}[120,130) & 36.28\\
{}[140,190) & 22.34\\
{}[5,10) & 15.58\\
\addlinespace
{}[100,110) & 9.63\\
\bottomrule
\end{tabular}
\begin{tablenotes}
\item \textit{Note: } 
\item diametre\_cl: classe de diametre; carbone\_total: stock de carbone total en tC/ha.
\end{tablenotes}
\end{threeparttable}}
\end{table}

\begin{figure}

{\centering \includegraphics[width=0.8\linewidth]{03-resultats_et_discussions_files/figure-latex/Barplotcarboneligneux-Nanga3-1} 

}

\caption{Répartition des stocks de carbone des ligneux dans les TUT de la FC de Nanga-Eboko.}\label{fig:Barplotcarboneligneux-Nanga3}
\end{figure}

L'analyse de la distribution des stocks de carbone des ligneux dans les 6 catégories d'occupation de terres dans la FC de Nanga-Eboko revèle dans l'ensemble peu de différences significatives dans la distribution des stocks de carbone. En effet sur les 15 couples de comparaison entre les catégories d'utilisation des terres, seules 04 présentent des différences significatives: on a relevé une différénce significative entre les Forêts secondaires adultes (FSA)-Forêts galeries (FGA) et Forêts secondaires adultes (FSA)-Savanes (SAV) d'une part, Forêts secondaires jeunes (FSJ)-Forêts galeries (FGA) et Forêts secondaires jeunes (FSJ)-Savanes (SAV) d'autre part (Figure:\ref{fig:ttestligneuxNanga-Eboko3}). Les zones savanicoles se distinguent des autres catégories à cause de la taille réduite des ligneux qui les constituent. En effet les plus grands arbres ne dépassent guère 40cm de diamètre. Quant aux forêts galeries la différence de stocks de carbone entre les plus petits et les plus grands arbres est supérieure à celle des Forêts secondaires adultes et des forêts secondaires jeunes (\ref{tab:StocksmoyensligneuxNanga-Eboko3}).

\begin{figure}

{\centering \includegraphics[width=0.8\linewidth]{03-resultats_et_discussions_files/figure-latex/ttestligneuxNanga-Eboko3-1} 

}

\caption{t.test des stocks de carbone des ligneux dans la FC de Nanga-Eboko.}\label{fig:ttestligneuxNanga-Eboko3}
\end{figure}

\begin{table}
\centering
\caption{\label{tab:StocksmoyensligneuxNanga-Eboko3}Calcul des stocks moyens de carbone des ligneux dans les TUT de la FC de Nanga-Eboko.}
\centering
\resizebox{\ifdim\width>\linewidth\linewidth\else\width\fi}{!}{
\begin{threeparttable}
\begin{tabular}[t]{lrrrr}
\toprule
Strates & length & mean & sd & range\\
\midrule
FSA & 374 & 0.83 & 1.27 & 11.51\\
FSJ & 182 & 1.04 & 1.52 & 10.78\\
FGA & 110 & 1.85 & 2.51 & 16.62\\
MAR & 60 & 1.41 & 2.59 & 17.23\\
SAV & 32 & 1.69 & 1.66 & 5.67\\
\addlinespace
CUL & 5 & 2.96 & 3.22 & 8.13\\
\bottomrule
\end{tabular}
\begin{tablenotes}
\item \textit{Note: } 
\item length: taille de l'échantillon; mean: moyenne; sd: écart type; range: amplitude; FSA: Forêts Secondaires Adultes; FGA: Forêts Galeries; FSJ: Forêts Secondaires Jeunes; MAR: Marécages; SAV: Savanes; CUL: Cultures.
\end{tablenotes}
\end{threeparttable}}
\end{table}

\hypertarget{estimation-des-stocks-de-carbone-dans-le-bois-mort}{%
\subsubsection{Estimation des stocks de carbone dans le bois mort}\label{estimation-des-stocks-de-carbone-dans-le-bois-mort}}

L'analyse de la distribution des stocks de carbone dans le bois mort de la FC de Ndikiniméki, indique que la moyenne des stocks se situe autour de 0,09 ± 0,19 tC/ha. Les mesures ont porté sur 364 individus repartis entre les 05 catégories d'utlisation des terres. Environ 40\% des mesures ont été faites dans les Forêts matures, 29\% dans les cacaoyères, 25\% dans les Forêts secondaires et 6\% dans les Jachères et les Forêts marécageuses:

\begin{verbatim}
#> Description of Ndiki_StC.BM
#> 
#>  Numeric 
#>             mean   sd median min  max valid.n
#> StC.BM.t.ha 0.09 0.19   0.02   0 1.78     364
#> 
#>  Factor 
#>          
#> Strates       FM    CAC FS   JA   FMa
#>   Count   146.00 105.00 91 12.0 10.00
#>   Percent  40.11  28.85 25  3.3  2.75
#> Mode FM
\end{verbatim}

En contraste avec les effectifs mesurés dans les différentes catégories, les stocks les plus importants ont été enregistrés dans les cacaoyères (18,20 tC/ha) (Tableau:\ref{tab:StockscarboneBM-Ndikinimeki3}) tandis que les Forêts matures arrivent en troisième position avec 5,21 tC/ha (Figure:\ref{fig:BarplotcarboneBM-Ndikinimeki3}). De même le t.test de Student a revelé qu'il existe une différence très significative dans la distribution des stocks entre les Cacaoyères et les Forêts matures. Le même test a également révélé des differences significatives entre Cacaoyères-Forêts secondaires d'une part et Forêts matures-Forêts secondaires d'autre part (Figure:\ref{fig:ttestBM-Ndikinimeki3}). Un écart type élevé dans les cacaoyères par rapport aux autres catégories traduit une grande dispersion des données, c'est également dans les cacaoyères qu'on a enregistré les bois morts ayant les tailles les plus imposantes (Tableau:\ref{tab:StocksmoyensBM-Ndikinimeki3}). En effet, les bois morts qu'on y trouvent présentent des allures très diversifiées en ce qui concerne leur diamètre, leur longueur et surtout leur densité spécifique liée au degré de décomposition du bois. Les cacaoyères rencontrées sont pour la plupart jeunes, on y rencontre davantage des bois morts avec des densités spécifiques élevées, la superposition des graphiques de densité des stocks relève bien l'importance de la moyenne des stocks de carbone des cacaoyères par rapport aux autres catégories. Cette situation n'est pas comparable aux forêts matures où les bois morts rencontrés sont en état de décomposition avancé. Les forêts secondaires, les forêts marécageuses et les jachères se situent à des niveaux intermédiaires où l'activité anthropique est moins marquée (Tableau:\ref{tab:StocksmoyensBM-Ndikinimeki3}).

\begin{figure}

{\centering \includegraphics[width=0.8\linewidth]{03-resultats_et_discussions_files/figure-latex/GraphiquedensiteBM-Ndikinimeki3-1} 

}

\caption{Graphiques de densité des stocks de carbone des bois morts dans la FC de Ndikiniméki.}\label{fig:GraphiquedensiteBM-Ndikinimeki3}
\end{figure}

\begin{table}[H]
\centering
\caption{\label{tab:StocksmoyensBM-Ndikinimeki3}Calcul des stocks moyens de carbone des bois morts dans les TUT de la FC de Ndikiniméki.}
\centering
\begin{threeparttable}
\begin{tabular}[t]{lrrrr}
\toprule
Strates & length & mean & sd & range\\
\midrule
FM & 146 & 0.04 & 0.07 & 0.37\\
CAC & 105 & 0.17 & 0.29 & 1.78\\
FS & 91 & 0.09 & 0.16 & 0.94\\
JA & 12 & 0.12 & 0.14 & 0.47\\
FMa & 10 & 0.08 & 0.18 & 0.59\\
\bottomrule
\end{tabular}
\begin{tablenotes}
\item \textit{Note: } 
\item length: taille de l'échantillon; mean: moyenne; sd: écart type; range: amplitude; FM: Forêt Mature; FS: Forêt Secondaire; FMa: Forêt Marécageuse; JA: Jachère; CAC: Cacaoyère; carbone\_total:stock de carbone total en tC/ha.
\end{tablenotes}
\end{threeparttable}
\end{table}

\begin{table}[H]
\centering
\caption{\label{tab:StockscarboneBM-Ndikinimeki3}Stocks de carbone des ligneux dans les différentes catégories d'utilisation de terres dans la FC de Ndikiniméki.}
\centering
\begin{threeparttable}
\begin{tabular}[t]{lr}
\toprule
Strates & Carbone.total\\
\midrule
CAC & 18.20\\
FS & 8.18\\
FM & 5.21\\
JA & 1.49\\
FMa & 0.77\\
\bottomrule
\end{tabular}
\begin{tablenotes}
\item \textit{Note: } 
\item FM: Forêt Mature; FS: Forêt Secondaire; FMa: Forêt Marécageuse; JA: Jachère; CAC: Cacaoyère; Carbone.total:stock de carbone total en tC/ha.
\end{tablenotes}
\end{threeparttable}
\end{table}

\begin{figure}

{\centering \includegraphics[width=0.8\linewidth]{03-resultats_et_discussions_files/figure-latex/ttestBM-Ndikinimeki3-1} 

}

\caption{t.test des stocks de carbone des bois morts dans la FC de Ndikiniméki.}\label{fig:ttestBM-Ndikinimeki3}
\end{figure}

\begin{figure}

{\centering \includegraphics[width=0.8\linewidth]{03-resultats_et_discussions_files/figure-latex/BarplotcarboneBM-Ndikinimeki3-1} 

}

\caption{Répartition des stocks de carbone des bois morts dans les TUT de la FC de Ndikiniméki.}\label{fig:BarplotcarboneBM-Ndikinimeki3}
\end{figure}

La distribution des stocks de carbone dans les bois morts de la FC de Nanga-Eboko indique une moyenne de 0,16 ± 0,41 tC/ha correspondant pratiquement au double de la moyenne des stocks dans la FC de Ndikiniméki:

\begin{verbatim}
#> Description of Nanga_StC.BM
#> 
#>  Numeric 
#>             mean   sd median min max valid.n
#> StC.BM.t.ha 0.16 0.41   0.03   0 3.9     166
#> 
#>  Factor 
#>          
#> Strates     FSA   MAR   FSJ   FGA   CUL  SAV
#>   Count   78.00 27.00 25.00 20.00 10.00 6.00
#>   Percent 46.99 16.27 15.06 12.05  6.02 3.61
#> Mode FSA
\end{verbatim}

Les forêts secondaires adultes représentent pratiquement 47\% des individus identifiés mais ne renferment que 23\% des stocks de carbone soit 6,33 tC/ha (Figure:\ref{fig:BarplotcarboneBMNanga-Eboko3}). Les Forêts galeries (FGA) par contre totalisent 9,49 tC/ha avec seulement 12\% des individus identifiés (Figure:\ref{tab:StockscarboneBMNanga-Eboko3}). Les Savanes (SAV) et les Cultures présentent les taux de carbone les plus faibles de toute la forêt communale. 0,38 tC/ha pour les Cultures et 0,05 tC/ha pour les Savanes (Figure:\ref{fig:GraphiquedensiteBMNanga-Eboko3}). Les cultures sont pratiquées sur des parcelles forestières presque dépourvues d'arbres. Ces cultures sont éssentiellement vivrières et les espèces ligneuses y sont peu nombreuses. Les zones savanicoles étant généralement exposées aux feux de brousse, il serait assez rare de trouver de grandes quantités de bois morts étalées sur le sol d'où le très faible taux de carbone enregistré. Le t.test de Student a montré une différence significative entre les Savanes (SAV) et les 05 autres catégories d'utilisation des terres en ce qui concerne la distribution des stocks de carbone (Figure:\ref{fig:ttestBMNanga-Eboko3}).

\begin{figure}

{\centering \includegraphics[width=0.8\linewidth]{03-resultats_et_discussions_files/figure-latex/GraphiquedensiteBMNanga-Eboko3-1} 

}

\caption{Graphiques de densité des stocks de carbone des bois morts dans la FC de Nanga-Eboko.}\label{fig:GraphiquedensiteBMNanga-Eboko3}
\end{figure}

\begin{table}[H]
\centering
\caption{(\#tab:Stocks moyens BM FC Nanga-Eboko)Calcul des stocks moyens de carbone des bois morts dans les TUT de la FC de Nanga-Eboko.}
\centering
\begin{threeparttable}
\begin{tabular}[t]{lrrrr}
\toprule
Strates & length & mean & sd & range\\
\midrule
FSA & 78 & 0.08 & 0.17 & 0.94\\
MAR & 27 & 0.20 & 0.44 & 2.01\\
FSJ & 25 & 0.22 & 0.37 & 1.36\\
FGA & 20 & 0.47 & 0.89 & 3.89\\
CUL & 10 & 0.04 & 0.07 & 0.22\\
\addlinespace
SAV & 6 & 0.01 & 0.01 & 0.02\\
\bottomrule
\end{tabular}
\begin{tablenotes}
\item \textit{Note: } 
\item length: taille de l'échantillon; mean: moyenne; sd: écart type; range: amplitude; FSA: Forêts Secondaires Adultes; FGA: Forêts Galeries; FSJ: Forêts Secondaires Jeunes; MAR: Marecages; SAV: Savanes; CUL: Cultures.
\end{tablenotes}
\end{threeparttable}
\end{table}

\begin{table}[H]
\centering
\caption{\label{tab:StockscarboneBMNanga-Eboko3}Stocks de carbone des bois morts dans les différentes catégories d'utilisation de terres dans la FC de Nanga-Eboko.}
\centering
\begin{threeparttable}
\begin{tabular}[t]{lr}
\toprule
Strates & Carbone.total\\
\midrule
FGA & 9.49\\
FSA & 6.33\\
FSJ & 5.44\\
MAR & 5.40\\
CUL & 0.38\\
\addlinespace
SAV & 0.05\\
\bottomrule
\end{tabular}
\begin{tablenotes}
\item \textit{Note: } 
\item length: taille de l'échantillon; mean: moyenne; sd: écart type; range: amplitude; FSA: Forêts Secondaires Adultes; FGA: Forêts Galeries; FSJ: Forêts Secondaires Jeunes; MAR: Marecages; SAV: Savanes; CUL: Cultures.
\end{tablenotes}
\end{threeparttable}
\end{table}

\begin{figure}

{\centering \includegraphics[width=0.8\linewidth]{03-resultats_et_discussions_files/figure-latex/ttestBMNanga-Eboko3-1} 

}

\caption{t.test des stocks de carbone des bois morts dans la FC de Nanga-Eboko.}\label{fig:ttestBMNanga-Eboko3}
\end{figure}

\begin{figure}

{\centering \includegraphics[width=0.8\linewidth]{03-resultats_et_discussions_files/figure-latex/BarplotcarboneBMNanga-Eboko3-1} 

}

\caption{Répartition des stocks de carbone des bois morts dans les TUT de la FC de Nanga-Eboko.}\label{fig:BarplotcarboneBMNanga-Eboko3}
\end{figure}

\hypertarget{estimation-des-stcoks-de-carbone-des-herbacuxe9es-dans-la-fc-de-ndikinimuxe9ki}{%
\subsection{Estimation des stcoks de carbone des herbacées dans la FC de Ndikiniméki}\label{estimation-des-stcoks-de-carbone-des-herbacuxe9es-dans-la-fc-de-ndikinimuxe9ki}}

L'estimation des stocks de carbone des herbacées dans la FC de Ndikiniméki a porté sur 160 quadrats de 1m² chacun repartis dans les 5 catégories d'utilisation des terres. Les stocks de carbone sont estimés à 10,58 tC/ha avec une moyenne de 0,06 tC/ha :

\begin{verbatim}
#> Description of Ndiki_StC.Herb
#> 
#>  Numeric 
#>               mean   sd median min  max valid.n
#> StC.Herb.t.ha 0.07 0.08   0.03   0 0.63     160
#> 
#>  Factor 
#>          
#> Strates      FM CAC    FS   FMa    JA
#>   Count   65.00  40 35.00 10.00 10.00
#>   Percent 40.62  25 21.88  6.25  6.25
#> Mode FM
\end{verbatim}

Les quantités de carbone les plus importantes ont été enregistrées dans les Cacaoyères avec 2,96 tC/ha représentant environ 30 \% des stocks de carbone (Figure:\ref{fig:Barplotcarboneherbacees-Ndikinimeki3}). Les Cacaoyères étant pour la plupart des agroforêts, une canopée ouverte favorise le développement des herbacées qui, à la faveur de la lumière prospèrent abondamment surtout avant les opérations de désherbage (Figure:\ref{fig:Graphiquedensiteherbacees-Ndikinimeki3}). Les Forêts marécageuses par contre détiennent le plus faible taux de carbone dans les herbacées (1,03 tC/ha). Cela peut se justifier par les conditions du milieu, en effet le sol étant gorgé d'eau, il ne favorise pas le développement de la plupart des espèces forestières (Tableau:\ref{tab:Stockscarboneherbacees-Ndikinimeki3}). L'analyse statistique de la distribution des stocks de carbone des herbacées montre qu'il existe une différence très significative entre les Forêts matures et les Forêts marécageuses, et significative entre Forêts matures-Forêts secondaires et Forêts matures-Cacaoyères (Figure:\ref{fig:ttestherbacees-Ndikinimeki3}).

\begin{figure}

{\centering \includegraphics[width=0.8\linewidth]{03-resultats_et_discussions_files/figure-latex/Graphiquedensiteherbacees-Ndikinimeki3-1} 

}

\caption{Graphiques de densité des stocks de carbone des herbacées dans la FC de Ndikiniméki.}\label{fig:Graphiquedensiteherbacees-Ndikinimeki3}
\end{figure}

\begin{table}[H]
\centering
\caption{\label{tab:Stocksmoyensherbacees-FCNdiki3}Calcul des stocks moyens de carbone des herbacées dans les TUT de la FC de Ndikiniméki.}
\centering
\begin{threeparttable}
\begin{tabular}[t]{lrrrr}
\toprule
Strates & length & mean & sd & range\\
\midrule
FM & 65 & 0.04 & 0.04 & 0.25\\
CAC & 40 & 0.07 & 0.08 & 0.33\\
FS & 35 & 0.07 & 0.06 & 0.21\\
FMa & 10 & 0.10 & 0.04 & 0.11\\
JA & 10 & 0.14 & 0.17 & 0.57\\
\bottomrule
\end{tabular}
\begin{tablenotes}
\item \textit{Note: } 
\item length: taille de l'échantillon; mean: moyenne; sd: écart type; range: amplitude; FM: Forêt Mature; FS: Forêt Secondaire; FMa: Forêt Marécageuse; JA: Jachère; CAC: Cacaoyère.
\end{tablenotes}
\end{threeparttable}
\end{table}

\begin{table}[H]
\centering
\caption{\label{tab:Stockscarboneherbacees-Ndikinimeki3}Stocks de carbone des herbacées dans les différentes catégories d'utilisation de terres dans la FC de Ndikiniméki.}
\centering
\begin{threeparttable}
\begin{tabular}[t]{lr}
\toprule
Strates & Carbone.total\\
\midrule
CAC & 2.97\\
FM & 2.64\\
FMa & 1.03\\
FS & 2.52\\
JA & 1.43\\
\bottomrule
\end{tabular}
\begin{tablenotes}
\item \textit{Note: } 
\item FM: Forêt Mature; FS: Forêt Secondaire; FMa: Forêt Marécageuse; JA: Jachère; CAC: Cacaoyère; Carbone.total:stock de carbone total en tC/ha.
\end{tablenotes}
\end{threeparttable}
\end{table}

\begin{figure}

{\centering \includegraphics[width=0.8\linewidth]{03-resultats_et_discussions_files/figure-latex/ttestherbacees-Ndikinimeki3-1} 

}

\caption{t.test des stocks de carbone des herbacées dans la FC de Ndikiniméki.}\label{fig:ttestherbacees-Ndikinimeki3}
\end{figure}

\begin{figure}

{\centering \includegraphics[width=0.8\linewidth]{03-resultats_et_discussions_files/figure-latex/Barplotcarboneherbacees-Ndikinimeki3-1} 

}

\caption{Répartition des stocks de carbone des herbacées dans les TUT de la FC de Ndikiniméki.}\label{fig:Barplotcarboneherbacees-Ndikinimeki3}
\end{figure}

\hypertarget{estimation-des-stcoks-de-carbone-des-herbacuxe9es-dans-la-fc-de-nanga-eboko}{%
\subsection{Estimation des stcoks de carbone des herbacées dans la FC de Nanga-Eboko}\label{estimation-des-stcoks-de-carbone-des-herbacuxe9es-dans-la-fc-de-nanga-eboko}}

L'estimation des stocks de carbone des herbacées dans la FC de Nanga-Eboko a porté sur 115 quadrats repartis entre les 06 catégories d'utilisation des terres:

\begin{verbatim}
#> Description of Nanga_StC.Herb
#> 
#>  Numeric 
#>               mean   sd median min  max valid.n
#> StC.Herb.t.ha 0.03 0.06   0.01   0 0.37     115
#> 
#>  Factor 
#>          
#> Strates     FSA   FSJ   MAR   FGA   SAV  CUL
#>   Count   35.00 20.00 20.00 15.00 15.00 10.0
#>   Percent 30.43 17.39 17.39 13.04 13.04  8.7
#> Mode FSA
\end{verbatim}

Les stocks de carbone ont été estimés à 3,84 tC/ha pour une moyenne de 0,03 ± 0,06 tC/ha ce qui représente à peu près le tiers des stocks enregistrés dans la FC de Ndikiniméki (Tableau:\ref{tab:StockscarboneherbaceesNanga-Eboko3}). Les espaces savanicoles totalisent 2,33 tC/ha soit 60\% des stocks de carbone enregistrés sur l'ensemble des herbacées de la Forêt communale de Nanga-Eboko (Figure:\ref{fig:BarplotcarboneherbaceesNanga-Eboko3}). Les 05 autres catégories se partagent les 40\% restant. Le t.test de Student révèle en effet qu'il existe une difference très significative entre Savane-Forêts galeries, Savanes-Forêts secondaires adultes, Savanes-Forêts secondaires jeunes et Savanes-Cultures (Figure:\ref{fig:ttestherbaceesNanga-Eboko3}). Ces savanes appartiennent à la catégories des savanes herbeuses où les arbres sont distants et de faibles envergures ce qui favorise le développement d'une végétation herbeuse très abondante.

\begin{table}[H]
\centering
\caption{\label{tab:StocksmoyensherbaceesFCNanga-Eboko3}Calcul des stocks moyens de carbone des herbacées dans les TUT de la FC de Nanga-Eboko.}
\centering
\begin{threeparttable}
\begin{tabular}[t]{lrrrr}
\toprule
Strates & length & mean & sd & range\\
\midrule
FSA & 35 & 0.01 & 0.01 & 0.06\\
FSJ & 20 & 0.01 & 0.01 & 0.05\\
MAR & 20 & 0.01 & 0.02 & 0.06\\
FGA & 15 & 0.03 & 0.03 & 0.14\\
SAV & 15 & 0.16 & 0.10 & 0.34\\
\addlinespace
CUL & 10 & 0.05 & 0.04 & 0.12\\
\bottomrule
\end{tabular}
\begin{tablenotes}
\item \textit{Note: } 
\item length: taille de l'échantillon; mean: moyenne; sd: écart type; range: amplitude; FSA: Forêts Secondaires Adultes; FGA: Forêts Galeries; FSJ: Forêts Secondaires Jeunes; MAR: Marecages; SAV: Savanes; CUL: Cultures.
\end{tablenotes}
\end{threeparttable}
\end{table}

\begin{table}[H]
\centering
\caption{\label{tab:StockscarboneherbaceesNanga-Eboko3}Stocks de carbone des herbacées dans les différentes catégories d'utilisation de terres dans la FC de Nanga-Eboko.}
\centering
\begin{threeparttable}
\begin{tabular}[t]{lr}
\toprule
Strates & Carbone.total\\
\midrule
SAV & 2.33\\
CUL & 0.46\\
FGA & 0.44\\
MAR & 0.28\\
FSA & 0.24\\
\addlinespace
FSJ & 0.11\\
\bottomrule
\end{tabular}
\begin{tablenotes}
\item \textit{Note: } 
\item FSA: Forêts Secondaires Adultes; FGA: Forêts Galeries; FSJ: Forêts Secondaires Jeunes; MAR: Marécages; SAV: Savanes; CUL: Cultures; Carbone.total: stock de carbone total en tC/ha.
\end{tablenotes}
\end{threeparttable}
\end{table}

\begin{figure}

{\centering \includegraphics[width=0.8\linewidth]{03-resultats_et_discussions_files/figure-latex/ttestherbaceesNanga-Eboko3-1} 

}

\caption{t.test des stocks de carbone des herbacées dans la FC de Nanga-Eboko.}\label{fig:ttestherbaceesNanga-Eboko3}
\end{figure}

\begin{figure}

{\centering \includegraphics[width=0.8\linewidth]{03-resultats_et_discussions_files/figure-latex/BarplotcarboneherbaceesNanga-Eboko3-1} 

}

\caption{Répartition des stocks de carbone des herbacées dans les TUT de la FC de Nanga-Eboko.}\label{fig:BarplotcarboneherbaceesNanga-Eboko3}
\end{figure}

\hypertarget{estimation-des-stocks-de-carbone-dans-la-litiuxe8re-de-la-fc-de-ndikinimuxe9ki}{%
\subsection{Estimation des stocks de carbone dans la litière de la FC de Ndikiniméki}\label{estimation-des-stocks-de-carbone-dans-la-litiuxe8re-de-la-fc-de-ndikinimuxe9ki}}

La distribution des stocks de carbone de la litière dans la FC de Ndikiniméki varit de 0,01 à 0,12 tC/ha pour une moyenne de 0,04 ± 0,02 tC/ha:

\begin{verbatim}
#> Description of Ndiki_StC.Lit
#> 
#>  Numeric 
#>              mean   sd median  min  max valid.n
#> StC_Lit.t.ha 0.04 0.02   0.04 0.01 0.12     161
#> 
#>  Factor 
#>          
#> Strates      FM   CAC    FS   FMa    JA
#>   Count   66.00 40.00 35.00 10.00 10.00
#>   Percent 40.99 24.84 21.74  6.21  6.21
#> Mode FM
\end{verbatim}

Le stock total de carbone dans les 5 catégories d'utilisation des terres est de 6,65 tC/ha. Les Cacaoyères (CAC) cumulent 1,81 tC/ha avec seulement 25\% des quadrats établis tandis que les Forêts matures (FM) totalisent 1,57 tC/ha pour 41\% des quadrats (Figure:\ref{fig:Barplotcarbonelitiere-Ndikinimeki3}). Bien qu'ayant la moyenne des stocks la plus élevée, les Jachères (JA) ont le total de stocks de carbone de litière le plus faible (0,68 tC/ha) Figure:\ref{fig:Graphiquedensitelitiere-Ndikinimeki3}. Les tests statistques ont également montrés qu'il existe une différence très significative dans la distribution des stocks entre les Jachères et les Forêts matures, les Forêts secondaires et les Cacaoyères. Il existe également une différence très significative entre les Forêts matures et les Forêts secondaires d'une part et les Forêts matures et les Cacaoyères d'autre part. Cependant il n'existe pas de différence significative entre les Cacaoyères et les Forêts marécageuses d'une part, les Cacaoyères et les Forêts secondaires d'autre part. De même il n'existe pas de différence significative entre les Forêts secondaires et les Forêts marécageuses (Figure:\ref{fig:ttestlitiereNdikinimeki3}).

\begin{figure}

{\centering \includegraphics[width=0.8\linewidth]{03-resultats_et_discussions_files/figure-latex/Graphiquedensitelitiere-Ndikinimeki3-1} 

}

\caption{Graphiques de densité des stocks de carbone de la litière dans la FC de Ndikiniméki.}\label{fig:Graphiquedensitelitiere-Ndikinimeki3}
\end{figure}

\begin{figure}

{\centering \includegraphics[width=0.8\linewidth]{03-resultats_et_discussions_files/figure-latex/ttestlitiereNdikinimeki3-1} 

}

\caption{t.test des stocks de carbone de la litière dans la FC de Ndikiniméki.}\label{fig:ttestlitiereNdikinimeki3}
\end{figure}

\begin{figure}

{\centering \includegraphics[width=0.8\linewidth]{03-resultats_et_discussions_files/figure-latex/Barplotcarbonelitiere-Ndikinimeki3-1} 

}

\caption{Répartition des stocks de carbone de la litière dans les TUT de la FC de Ndikiniméki.}\label{fig:Barplotcarbonelitiere-Ndikinimeki3}
\end{figure}

\hypertarget{estimation-des-stocks-de-carbone-dans-la-litiuxe8re-de-la-fc-de-nanga-eboko}{%
\subsection{Estimation des stocks de carbone dans la litière de la FC de Nanga-Eboko}\label{estimation-des-stocks-de-carbone-dans-la-litiuxe8re-de-la-fc-de-nanga-eboko}}

La distribution des stocks de carbone de la litière dans la FC de Nanga-Eboko est à peu près équivalente à celle de la FC de Ndikiniméki avec des valeurs oscillant entre 0,01 et 0,03 tC/ha pour une moyenne de 0,03 ± 0,02 tC/ha (Figure:\ref{fig:GraphiquedensitlitiereNanga-Eboko3}). Les mesures de stocks éffectuées ont porté sur 115 quadrats de 1m² chacun dont 30\% appartenaient aux Forêts secondaires adultes (FSA), 17\% aux Forêts secondaires jeunes (FSJ), 17\% aux zones marécageuses (MAR), 17\% aux Forêts galeries (FGA), 13\% aux savanes et 9\% aux Cultures (CUL):

\begin{verbatim}
#> Description of Nanga_StC.Lit
#> 
#>  Numeric 
#>              mean   sd median  min  max valid.n
#> StC_Lit.t.ha 0.03 0.02   0.03 0.01 0.13     115
#> 
#>  Factor 
#>          
#> Strates     FSA   FSJ   MAR   FGA   SAV  CUL
#>   Count   35.00 20.00 20.00 15.00 15.00 10.0
#>   Percent 30.43 17.39 17.39 13.04 13.04  8.7
#> Mode FSA
\end{verbatim}

Bien que les zones de cultures ont enregistré le stock moyen le plus élevé (0,04 ± 0,03 tC/ha),elles ont eu le stock total le plus faible et ce sont les Forêts secondaires adultes qui ont obtenues le stock cumulé le plus élevé (0,82 tC/ha) avec l'une des moyennes les plus faibles (Figure:\ref{fig:BarplotcarbonelitiereNanga-Eboko3}). Du point de vue de l'analyse statistique, en dehors des differences significatives entre Forêts secondaires adultes-Forêts galeries et assez significatives entre Forêts secondaires adultes-Forêts secondaires jeunes, il n'a pas été relevé de differences significatives dans la distribution des stocks de carbone de la litière dans les 06 catégories d'utilisation des terres de la FC de Nanga-Eboko (Figure:\ref{fig:ttestlitiereNanga-Eboko3}).

\begin{figure}

{\centering \includegraphics[width=0.8\linewidth]{03-resultats_et_discussions_files/figure-latex/GraphiquedensitlitiereNanga-Eboko3-1} 

}

\caption{Graphiques de densité des stocks de carbone de la litière dans la FC de Nanga-Eboko.}\label{fig:GraphiquedensitlitiereNanga-Eboko3}
\end{figure}

\begin{table}[H]
\centering
\caption{\label{tab:StockscarbonelitiereNanga-Eboko3}Stocks de carbone de la litière dans les différentes catégories d'utilisation de terres dans la FC de Nanga-Eboko.}
\centering
\begin{threeparttable}
\begin{tabular}[t]{lr}
\toprule
Strates & Carbone.total\\
\midrule
FSA & 0.82\\
FSJ & 0.64\\
MAR & 0.58\\
FGA & 0.50\\
CUL & 0.41\\
\addlinespace
SAV & 0.37\\
\bottomrule
\end{tabular}
\begin{tablenotes}
\item \textit{Note: } 
\item FSA: Forêts Secondaires Adultes; FGA: Forêts Galeries; FSJ: Forêts Secondaires Jeunes; MAR: Marécages; SAV: Savanes; CUL: Cultures; Carbone.total: stock de carbone total en tC/ha.
\end{tablenotes}
\end{threeparttable}
\end{table}

\begin{figure}

{\centering \includegraphics[width=0.8\linewidth]{03-resultats_et_discussions_files/figure-latex/ttestlitiereNanga-Eboko3-1} 

}

\caption{t.test des stocks de carbone de la litière dans la FC de Nanga-Eboko.}\label{fig:ttestlitiereNanga-Eboko3}
\end{figure}

\begin{figure}

{\centering \includegraphics[width=0.8\linewidth]{03-resultats_et_discussions_files/figure-latex/BarplotcarbonelitiereNanga-Eboko3-1} 

}

\caption{Répartition des stocks de carbone de la litière dans les TUT de la FC de Nanga-Eboko.}\label{fig:BarplotcarbonelitiereNanga-Eboko3}
\end{figure}

\hypertarget{estimation-des-stocks-de-carbone-du-sol-dans-la-fc-de-ndikinimuxe9ki}{%
\subsection{Estimation des stocks de carbone du sol dans la FC de Ndikiniméki}\label{estimation-des-stocks-de-carbone-du-sol-dans-la-fc-de-ndikinimuxe9ki}}

L'analyse de la distribution des stocks de carbone du sol dans la FC de Ndikiniméki s'est faite de manière générale sur un total de 69 échantillons de terre prélevés dans chacun des trois horizons du sol de chacune des 23 parcelles d'échantillonnages (0-10cm ; 10-20cm et 20-30cm), dont 39\% dans les Forêts matures (FM), 44\% dans les Forêts secondaires (FS) et les Cacaoyères (CAC), et 18\% dans les Forêts marécageuses (FMa) et les Jachères (JA). Les résultats montrent que les stocks de carbone varient de 10,46 à 48,38 tC/ha avec une moyenne de 19,75 ± 7,19 tC/ha:

\begin{verbatim}
#> Description of Ndiki_StC.Sol
#> 
#>  Numeric 
#>          mean   sd median   min   max valid.n
#> MCsolHa 19.75 7.19  18.21 10.46 48.38      69
#> 
#>  Factor 
#>          
#> Depth      0_10 10_20 20_30
#>   Count   23.00 23.00 23.00
#>   Percent 33.33 33.33 33.33
#> Mode >1 mode 
#>          
#> Strates      FM   CAC    FS FMa  JA
#>   Count   27.00 15.00 15.00 6.0 6.0
#>   Percent 39.13 21.74 21.74 8.7 8.7
#> Mode FM
\end{verbatim}

Les couches les plus proches de la surface ont enregistré les stocks de carbone les plus importants avec une moyenne de 24,48 ± 5,44 tC/ha pour les profondeurs de 0 à 10 cm. La moyenne des stocks de carbone dans les couches moyennes (10 à 20 cm) et profondes (20 à 30 cm) sont respectivement de 18,13 ± 3 tC/ha et 15,80 ± 12,71 tC/ha:

\begin{table}[H]
\centering
\caption{\label{tab:Descriptionsommairesol-Ndikinimeki3} Répartition des stocks de carbone selon les profondeurs du sol dans les TUT de la FC de Ndikiniméki.}
\centering
\begin{threeparttable}
\begin{tabular}[t]{lrrr}
\toprule
  & 0\_10 & 10\_20 & 20\_30\\
\midrule
FMa & 30.81 & 20.93 & 21.24\\
FM & 26.76 & 18.47 & 15.93\\
CAC & 25.93 & 17.64 & 14.44\\
FS & 22.64 & 20.28 & 16.00\\
JA & 16.25 & 13.32 & 11.41\\
\bottomrule
\end{tabular}
\begin{tablenotes}
\item \textit{Note: } 
\item 0\_10: Profodeur de 0 à 10 cm; 10\_20: Profondeur de 10 à 20 cm; 20\_30 cm: Profondeur de 20 à 30 cm; FM: Forêt Mature; FS: Forêt Secondaire; FMa: Forêt Marécageuse; JA: Jachère; CAC: Cacaoyère.
\end{tablenotes}
\end{threeparttable}
\end{table}

\begin{verbatim}
#> Description of Ndiki_Depth.Sol
#> 
#>  Numeric 
#>        mean median   var   sd valid.n
#> 0_10  24.48  25.93 29.63 5.44       5
#> 10_20 18.13  18.47  8.99 3.00       5
#> 20_30 15.80  15.93 12.71 3.56       5
\end{verbatim}

En terme de proportion, les forêts marécageuses présentent la moyenne des stocks de carbone la plus importante (24,32 ± 11,39 tC/ha) avec une courbe de distribution étalée qui traduit une répartition équilibrée des stocks de carbone dans les différentes profondeurs. Les jachères ont la moyenne de stocks la moins importante (13,65 ± 2,51 tC/ha) avec une courbe de distribution unimodale autour de la moyenne, ce qui laisse penser à une répartition des stocks de carbone fortement marquée par un paramètre du milieu (l'épaisseur des couches constitutives du sol). Les courbes de répartition des stocks de carbone des forêts matures et des cacaoyères semblent obéir à la même logique que celle des jachères, cependant cette influence y est moins marquée. La courbe de répartition des stocks carbone du sol dans les forêts secondaires par contre est bimodale, ce qui pourrait supposer que les stocks de carbone dans cette catégorie seraient influencés par au moins deux paramètres (l'épaisseur des couches du sol et l'altitude) Figure:\ref{fig:Graphiquedensitesol-Ndikinimeki3}. Le t.test de Student ne montre pas de différence significative dans la répartition des stocks de carbone des 05 catégories d'utilisation des terres de la FC de Ndikiniméki sauf pour les jachères. En effet le test indique qu'il existe des differences significatives entre: Jachères-Cacaoyères, Jachères-Forêts secondaires et Jachères-Forêts matures (Figure:\ref{fig:ttestsol-Ndikinimeki3}).

\begin{figure}

{\centering \includegraphics[width=0.8\linewidth]{03-resultats_et_discussions_files/figure-latex/Graphiquedensitesol-Ndikinimeki3-1} 

}

\caption{Graphiques de densité des stocks de carbone de la sol dans la FC de Ndikiniméki.}\label{fig:Graphiquedensitesol-Ndikinimeki3}
\end{figure}

\begin{table}
\centering
\caption{\label{tab:Stocksmoyenssol-Ndikinimeki3}Calcul des stocks moyens de carbone du sol dans les TUT de la FC de Ndikiniméki.}
\centering
\begin{threeparttable}
\begin{tabular}[t]{lrrrr}
\toprule
Strates & length & mean & sd & range\\
\midrule
FM & 27 & 20.39 & 8.07 & 37.39\\
CAC & 15 & 19.34 & 5.78 & 20.06\\
FS & 15 & 19.64 & 4.75 & 12.88\\
FMa & 6 & 24.33 & 11.39 & 28.84\\
JA & 6 & 13.66 & 2.51 & 6.76\\
\bottomrule
\end{tabular}
\begin{tablenotes}
\item \textit{Note: } 
\item length: taille de l'échantillon; mean: moyenne; sd: écart type; range: amplitude; FM: Forêt Mature; FS: Forêt Secondaire; FMa: Forêt Marécageuse; JA: Jachère; CAC: Cacaoyère.
\end{tablenotes}
\end{threeparttable}
\end{table}

Lorsqu'on observe les courbes de distribution des stocks de carbone dans le sol en fonction de la profondeur, on constate de manière sommaire, que les stocks de carbone diminuent avec la profondeur dans toutes les catégories d'utilisation des terres sauf dans les forêts marécageuses où ils tendent à se stabiliser à partir de 20 cm de profondeur (Figure:\ref{fig:Gradientcarbonesol-Ndikinimeki3} et Tableau:\ref{tab:Descriptionsommairesol-Ndikinimeki3}). En terme de stoks cumulés, les quantités les plus importantes sont localisées dans les forêts marécageuses avec une estimation de 72,97 tC/ha, suivies des forêts matures (61,15 tC/ha), les forêts secondaires (58,9 tC/ha) et les Cacaoyères avec environ (58 tC/ha). Les stocks de carbone les plus faibles sont localisés dans les jachères (40,95 tC/ha) Figures:\ref{fig:Barplotcarbonesol-Ndikinimeki3} et \ref{fig:Barplotcarboneprofondeursol-Ndikinimeki3}.

\begin{table}[H]
\centering
\caption{\label{tab:Stockscarbonesol-Ndikinimeki3}Stocks de carbone du sol dans les différentes catégories d'utilisation de terres dans la FC de Ndikiniméki.}
\centering
\begin{threeparttable}
\begin{tabular}[t]{lr}
\toprule
Strates & Carbone.total\\
\midrule
FMa & 24.33\\
FM & 20.39\\
FS & 19.64\\
CAC & 19.34\\
JA & 13.66\\
\bottomrule
\end{tabular}
\begin{tablenotes}
\item \textit{Note: } 
\item FM: Forêt Mature; FS: Forêt Secondaire; FMa: Forêt Marécageuse; JA: Jachère; CAC: Cacaoyère; Carbone.total:stock de carbone total en tC/ha.
\end{tablenotes}
\end{threeparttable}
\end{table}

\begin{figure}

{\centering \includegraphics[width=0.8\linewidth]{03-resultats_et_discussions_files/figure-latex/ttestsol-Ndikinimeki3-1} 

}

\caption{t.test des stocks de carbone du sol dans la FC de Ndikiniméki.}\label{fig:ttestsol-Ndikinimeki3}
\end{figure}

\begin{figure}

{\centering \includegraphics[width=0.8\linewidth]{03-resultats_et_discussions_files/figure-latex/Barplotcarbonesol-Ndikinimeki3-1} 

}

\caption{Répartition des stocks de carbone du sol dans les TUT de la FC de Ndikiniméki.}\label{fig:Barplotcarbonesol-Ndikinimeki3}
\end{figure}

\begin{figure}

{\centering \includegraphics[width=0.8\linewidth]{03-resultats_et_discussions_files/figure-latex/Barplotcarboneprofondeursol-Ndikinimeki3-1} 

}

\caption{Répartition des stocks de carbone du sol suivant la profondeur dans les TUT de la FC de Ndikiniméki.}\label{fig:Barplotcarboneprofondeursol-Ndikinimeki3}
\end{figure}

\begin{figure}

{\centering \includegraphics[width=0.8\linewidth]{03-resultats_et_discussions_files/figure-latex/Gradientcarbonesol-Ndikinimeki3-1} 

}

\caption{Gradient de distribution des stocks de carbone  du sol dans la FC de Ndikiniméki.}\label{fig:Gradientcarbonesol-Ndikinimeki3}
\end{figure}

\hypertarget{stocks-de-carbone-du-sol-dans-les-profondeurs-de-0-uxe0-10-cm}{%
\subsubsection{Stocks de carbone du sol dans les profondeurs de 0 à 10 cm}\label{stocks-de-carbone-du-sol-dans-les-profondeurs-de-0-uxe0-10-cm}}

A la profondeur de 0 à 10 cm, 23 échantillons de terre ont été prélevés dans l'ensemble des 05 catégories d'utilisation des terres. L'analyse des échantillons montre que la distribution des stocks de carbone varie de 13,86 à 48,38 tC/ha pour une moyenne de 25,12 ± 7,75tC/ha:

\begin{verbatim}
#> Description of Ndiki_StC_Sol.0_10
#> 
#>  Numeric 
#>          mean   sd median   min   max valid.n
#> MCsolHa 25.12 7.75  23.98 13.86 48.38      23
#> 
#>  Factor 
#>          
#> Depth     0_10 10_20 20_30
#>   Count     23     0     0
#>   Percent  100     0     0
#> Mode 0_10 
#>          
#> Strates      FM   CAC    FS FMa  JA
#>   Count    9.00  5.00  5.00 2.0 2.0
#>   Percent 39.13 21.74 21.74 8.7 8.7
#> Mode FM
\end{verbatim}

C'est l'horizon le plus proche de la surface où abondent les débris de matières organiques telles que les feuilles et les fragments de bois morts. La moyenne la plus élevée a été enregistrée dans les Forêts marécageuses (30,8 ± 12 tC/ha) qui sont des écosystèmes de basses altitudes et réceptacle des débris organiques transportés par l'eau ou le vent (Tableau:\ref{tab:Stockscarbonesolde0-10cmNdikinimeki3}). Au niveau de cette couche superficielle, le test de Student indique une difference signicative dans la distribution des stocks de carbone entre les Cacaoyères et les Jachères (Figure:\ref{fig:ttestsol0-10cmNdikinimeki3}).

\begin{table}[H]
\centering
\caption{\label{tab:Stockscarbonesolde0-10cmNdikinimeki3}Stocks de carbone du sol de 0 à 10 cm dans les différentes catégories d'utilisation de terres dans la FC de Ndikiniméki.}
\centering
\begin{threeparttable}
\begin{tabular}[t]{lr}
\toprule
Strates & total.Carbone\\
\midrule
FMa & 30.81\\
FM & 26.76\\
CAC & 25.93\\
FS & 22.64\\
JA & 16.25\\
\bottomrule
\end{tabular}
\begin{tablenotes}
\item \textit{Note: } 
\item FM: Forêt Mature; FS: Forêt Secondaire; FMa: Forêt Marécageuse; JA: Jachère; CAC: Cacaoyère; Carbone.total:stock de carbone total en tC/ha.
\end{tablenotes}
\end{threeparttable}
\end{table}

\begin{table}[H]
\centering
\caption{\label{tab:Stocksmoyenssol0-10cmNdikinimki3}Calcul des stocks moyens de carbone du sol de 0 à 10 cm dans les TUT de la FC de Ndikiniméki.}
\centering
\begin{threeparttable}
\begin{tabular}[t]{lrrrr}
\toprule
Strates & length & mean & sd & range\\
\midrule
FM & 9 & 26.762579 & 9.9530530 & 34.5221563\\
CAC & 5 & 25.932698 & 4.0806765 & 11.0552267\\
FS & 5 & 22.637030 & 2.9252690 & 7.1349337\\
FMa & 2 & 30.809216 & 12.0079011 & 16.9817366\\
JA & 2 & 16.248602 & 1.4087725 & 1.9923053\\
\bottomrule
\end{tabular}
\begin{tablenotes}
\item \textit{Note: } 
\item length: taille de l'échantillon; mean: moyenne; sd: écart type; range: amplitude; FM: Forêt Mature; FS: Forêt Secondaire; FMa: Forêt Marécageuse; JA: Jachère; CAC: Cacaoyère.
\end{tablenotes}
\end{threeparttable}
\end{table}

\begin{figure}

{\centering \includegraphics[width=0.8\linewidth]{03-resultats_et_discussions_files/figure-latex/ttestsol0-10cmNdikinimeki3-1} 

}

\caption{t.test des stocks de carbone du sol de 0 à 10 cm dans la FC de Ndikiniméki.}\label{fig:ttestsol0-10cmNdikinimeki3}
\end{figure}

\hypertarget{stocks-de-carbone-du-sol-dans-les-profondeurs-de-10-uxe0-20cm}{%
\subsubsection{Stocks de carbone du sol dans les profondeurs de 10 à 20cm}\label{stocks-de-carbone-du-sol-dans-les-profondeurs-de-10-uxe0-20cm}}

Dans cet horizon moyen l'écart entre le maximum et le minimum de stock enregistré est plus faible que celui de l'horizon superficiel. De même, c'est un horizon moins riche en carbone, les distributions des stocks de carbone vont de 11,88 à 29,47 tC/ha pour une moyenne de 18,45 ± 4,82 tC/ha:

\begin{verbatim}
#> Description of Ndiki_StC_Sol.10_20
#> 
#>  Numeric 
#>          mean   sd median   min   max valid.n
#> MCsolHa 18.45 4.82   18.4 11.88 29.47      23
#> 
#>  Factor 
#>          
#> Depth     10_20 0_10 20_30
#>   Count      23    0     0
#>   Percent   100    0     0
#> Mode 10_20 
#>          
#> Strates      FM   CAC    FS FMa  JA
#>   Count    9.00  5.00  5.00 2.0 2.0
#>   Percent 39.13 21.74 21.74 8.7 8.7
#> Mode FM
\end{verbatim}

Les Forêts marécageuses et les Forêts secondaires enregistrent les moyennes de stock les plus élevées (plus de 20 tC/ha). Les jachères demeurent les plus faibles en stocks de carbone dans cet horizon (Tableau:\ref{tab:Stockscarbonesolde10-20cmNdikinimeki3}). Cependant les analyses statisques ne montrent aucune difference significative dans la distribution des stocks de carbone dans cet horizon (Figure:\ref{fig:ttestsol10-20cmNdikinimeki3}).

\begin{table}[H]
\centering
\caption{\label{tab:Stockscarbonesolde10-20cmNdikinimeki3}Stocks de carbone du sol de 10 à 20 cm dans les différentes catégories d'utilisation de terres dans la FC de Ndikiniméki.}
\centering
\begin{threeparttable}
\begin{tabular}[t]{lr}
\toprule
Strates & total.Carbone\\
\midrule
FMa & 20.93\\
FS & 20.28\\
FM & 18.47\\
CAC & 17.64\\
JA & 13.32\\
\bottomrule
\end{tabular}
\begin{tablenotes}
\item \textit{Note: } 
\item FM: Forêt Mature; FS: Forêt Secondaire; FMa: Forêt Marécageuse; JA: Jachère; CAC: Cacaoyère; Carbone.total:stock de carbone total en tC/ha.
\end{tablenotes}
\end{threeparttable}
\end{table}

\begin{table}[H]
\centering
\caption{\label{tab:Stocksmoyenssol10-20cmFCNdikinimeki}Calcul des stocks moyens de carbone du sol de 10 à 20 cm dans les TUT de la FC de Ndikiniméki.}
\centering
\begin{threeparttable}
\begin{tabular}[t]{lrrrr}
\toprule
Strates & length & mean & sd & range\\
\midrule
FM & 9 & 18.47 & 3.24 & 10.65\\
CAC & 5 & 17.64 & 3.07 & 7.17\\
FS & 5 & 20.28 & 6.24 & 12.53\\
FMa & 2 & 20.93 & 12.08 & 17.08\\
JA & 2 & 13.32 & 2.04 & 2.88\\
\bottomrule
\end{tabular}
\begin{tablenotes}
\item \textit{Note: } 
\item length: taille de l'échantillon; mean: moyenne; sd: écart type; range: amplitude; FM: Forêt Mature; FS: Forêt Secondaire; FMa: Forêt Marécageuse; JA: Jachère; CAC: Cacaoyère.
\end{tablenotes}
\end{threeparttable}
\end{table}

\begin{figure}

{\centering \includegraphics[width=0.8\linewidth]{03-resultats_et_discussions_files/figure-latex/ttestsol10-20cmNdikinimeki3-1} 

}

\caption{t.test des stocks de carbone du sol de 10 à 20 cm dans la FC de Ndikiniméki.}\label{fig:ttestsol10-20cmNdikinimeki3}
\end{figure}

\hypertarget{stocks-de-carbone-du-sol-dans-les-profondeurs-de-20-uxe0-30cm}{%
\subsubsection{Stocks de carbone du sol dans les profondeurs de 20 à 30cm}\label{stocks-de-carbone-du-sol-dans-les-profondeurs-de-20-uxe0-30cm}}

Dans les profondeurs de 20 à 30 cm les stocks oscillent entre 10,46 à 32,02 tC/ha avec une moyenne de 15,69 ± 5,24 tC/ha:

\begin{verbatim}
#> Description of Ndiki_StC_Sol.20_30
#> 
#>  Numeric 
#>          mean   sd median   min   max valid.n
#> MCsolHa 15.69 5.24  14.97 10.46 32.02      23
#> 
#>  Factor 
#>          
#> Depth     20_30 0_10 10_20
#>   Count      23    0     0
#>   Percent   100    0     0
#> Mode 20_30 
#>          
#> Strates      FM   CAC    FS FMa  JA
#>   Count    9.00  5.00  5.00 2.0 2.0
#>   Percent 39.13 21.74 21.74 8.7 8.7
#> Mode FM
\end{verbatim}

Les valeurs moyennes par catégories les plus importantes se localisent dans les forêts marécageuses (21,24 ± 15,25 tC/ha) suivies des forêts secondaires (16 ± 1,83 tC/ha) et des forêts matures (15,93 ± 5,45). Les jachères demeurent la catégorie à faible taux de carbone (11,4 ± 1,3 tC/ha) Tableau:\ref{tab:Stockscarbonesolde20-30cmNdikinimeki3}. Ce sont pour la plupart des terres remaniées ayant subies de fortes activités agricoles qui ont pour corollaire un appauvrissement en matières organiques. Les tests statistiques montrent dans l'ensemble qu'à la profondeur de 20-30cm, il n'existe pas de différences significatives dans la distribution des stocks de carbone. Seules les forêts secondaires et jachères présentent une différence significative (Figure:\ref{fig:ttestsol20-30cmFCNdikinimeki}).

\begin{table}[H]
\centering
\caption{\label{tab:Stockscarbonesolde20-30cmNdikinimeki3}Stocks de carbone du sol de 20 à 30 cm dans les différentes catégories d'utilisation de terres dans la FC de Ndikiniméki.}
\centering
\begin{threeparttable}
\begin{tabular}[t]{lr}
\toprule
Strates & total.Carbone\\
\midrule
FMa & 21.24\\
FS & 16.00\\
FM & 15.93\\
CAC & 14.44\\
JA & 11.41\\
\bottomrule
\end{tabular}
\begin{tablenotes}
\item \textit{Note: } 
\item FM: Forêt Mature; FS: Forêt Secondaire; FMa: Forêt Marécageuse; JA: Jachère; CAC: Cacaoyère; Carbone.total:stock de carbone total en tC/ha.
\end{tablenotes}
\end{threeparttable}
\end{table}

\begin{table}[H]
\centering
\caption{\label{tab:Stocksmoyenssol20-30cmFCNdikinimeki}Calcul des stocks moyens de carbone du sol de 20 à 30 cm dans les TUT de la FC de Ndikiniméki.}
\centering
\begin{threeparttable}
\begin{tabular}[t]{lrrrr}
\toprule
Strates & length & mean & sd & range\\
\midrule
FM & 9 & 15.93 & 5.46 & 18.09\\
CAC & 5 & 14.44 & 1.74 & 4.79\\
FS & 5 & 16.00 & 1.83 & 4.54\\
FMa & 2 & 21.24 & 15.25 & 21.57\\
JA & 2 & 11.41 & 1.30 & 1.84\\
\bottomrule
\end{tabular}
\begin{tablenotes}
\item \textit{Note: } 
\item length: taille de l'échantillon; mean: moyenne; sd: écart type; range: amplitude; FM: Forêt Mature; FS: Forêt Secondaire; FMa: Forêt Marécageuse; JA: Jachère; CAC: Cacaoyère.
\end{tablenotes}
\end{threeparttable}
\end{table}

\begin{figure}

{\centering \includegraphics[width=0.8\linewidth]{03-resultats_et_discussions_files/figure-latex/ttestsol20-30cmFCNdikinimeki-1} 

}

\caption{t.test des stocks de carbone du sol de 20 à 30 cm dans la FC de Ndikiniméki.}\label{fig:ttestsol20-30cmFCNdikinimeki}
\end{figure}

\hypertarget{estimation-des-stocks-de-carbone-du-sol-dans-la-fc-de-nanga-eboko}{%
\subsection{Estimation des stocks de carbone du sol dans la FC de Nanga-Eboko}\label{estimation-des-stocks-de-carbone-du-sol-dans-la-fc-de-nanga-eboko}}

L'analyse de la distribution des stocks de carbone du sol dans la FC de Nanga-Eboko s'est faite dans les mêmes conditions que celle de la FC de Ndikiniméki à savoir 69 échantillons de terres prélevés repartis entre les 06 catégories d'occupation des terres. Parmi ces échantillons, 21 (environ 30\%) appartiennent aux Forêts secondaires adultes, 24 (environ 34\%) aux Forêts secondaires jeunes et aux zones marécageuses, 18 (environ 26\%) aux Forêts galeries et Savanes, et 06 (environ 9\%) aux zones de Cultures. Les résultats d'analyses de l'ensemble des échantillons montrent que les stocks de carbone du sol vont de 4,37 à 35,07 tC/ha pour une moyenne de 15,08 ± 5,92 tC/ha:

\begin{verbatim}
#> Description of Nanga_StC.Sol
#> 
#>  Numeric 
#>          mean   sd median  min   max valid.n
#> MCsolHa 15.08 5.92  14.12 4.37 35.07      69
#> 
#>  Factor 
#>          
#> Depth      0_10 10_20 20_30
#>   Count   23.00 23.00 23.00
#>   Percent 33.33 33.33 33.33
#> Mode >1 mode 
#>          
#> Strates     FSA   FSJ   MAR   FGA   SAV CUL
#>   Count   21.00 12.00 12.00  9.00  9.00 6.0
#>   Percent 30.43 17.39 17.39 13.04 13.04 8.7
#> Mode FSA
\end{verbatim}

En se rapportant aux catégories d'utilisation des terres, la moyenne de stock de carbone la plus élevée a été enregistrée dans les zones de cultures (16,38 ± 9,04 tC/ha). Elles sont suivies par les Forêts secondaires adultes (15,93 ± 5,9 tC/ha), les Forêts galeries (15,64 ± 5,43 tC/ha), les Forêts secondaires jeunes (15,37 ± 2,88 tC/ha), les Savanes (15,33 ± 4,77 tC/ha) et les Marécages (12,01 ± 7,52 tC/ha) Figure:\ref{fig:GraphiquedensitesolNanga-Eboko3}). En terme de stocks cumulés, les Cultures arrivent en tête avec près de 49,14 tC/ha tandis que les zones de Marécages enregistrent le plus faible taux (36,02 tC/ha) Figures:\ref{fig:BarplotcarbonesolNanga-Eboko3} et \ref{fig:BarplotcarboneprofondeursolNanga-Eboko3}. Les analyses statistiques montrent qu'il n'existe pas de différences significatives entre les différentes catégories d'utilisation des terres dans la distribution des stocks de carbone du sol (Figure:\ref{fig:ttestsolNanga-Eboko3}).

\begin{figure}

{\centering \includegraphics[width=0.8\linewidth]{03-resultats_et_discussions_files/figure-latex/GraphiquedensitesolNanga-Eboko3-1} 

}

\caption{Graphiques de densité des stocks de carbone de la sol dans la FC de Nanga-Eboko.}\label{fig:GraphiquedensitesolNanga-Eboko3}
\end{figure}

\begin{table}[H]
\centering
\caption{\label{tab:StocksmoyenssolNanga-Eboko3}Calcul des stocks moyens de carbone du sol dans les TUT de la FC de Nanga-Eboko.}
\centering
\begin{threeparttable}
\begin{tabular}[t]{lrrrr}
\toprule
Strates & length & mean & sd & range\\
\midrule
CUL & 6 & 16.38 & 9.05 & 23.93\\
FSA & 21 & 15.93 & 5.91 & 25.83\\
FGA & 9 & 15.65 & 5.44 & 15.33\\
FSJ & 12 & 15.38 & 2.89 & 7.67\\
SAV & 9 & 15.33 & 4.78 & 14.12\\
\addlinespace
MAR & 12 & 12.01 & 7.53 & 24.56\\
\bottomrule
\end{tabular}
\begin{tablenotes}
\item \textit{Note: } 
\item length: taille de l'échantillon; mean: moyenne; sd: écart type; range: amplitude; FSA: Forêts Secondaires Adultes; FGA: Forêts Galeries; FSJ: Forêts Secondaires Jeunes; MAR: Marécages; SAV: Savanes; CUL: Cultures.
\end{tablenotes}
\end{threeparttable}
\end{table}

\begin{table}[H]
\centering
\caption{\label{tab:DescriptionsommairesolNanga-Eboko3} Répartition des stocks de carbone selon les profondeurs du sol dans les TUT de la FC de Nanga-Eboko.}
\centering
\begin{threeparttable}
\begin{tabular}[t]{lrrr}
\toprule
  & 0\_10 & 10\_20 & 20\_30\\
\midrule
CUL & 22.49 & 15.39 & 11.27\\
FGA & 20.97 & 14.57 & 11.41\\
FSA & 20.20 & 14.40 & 13.19\\
MAR & 18.36 & 10.52 & 7.16\\
SAV & 17.29 & 15.61 & 13.10\\
\addlinespace
FSJ & 16.60 & 16.47 & 13.07\\
\bottomrule
\end{tabular}
\begin{tablenotes}
\item \textit{Note: } 
\item 0\_10: Profodeur de 0 à 10 cm; 10\_20: Profondeur de 10 à 20 cm; 20\_30 cm: Profondeur de 20 à 30 cm; FSA: Forêts Secondaires Adultes; FGA: Forêts Galeries; FSJ: Forêts Secondaires Jeunes; MAR: Marécages; SAV: Savanes; CUL: Cultures.
\end{tablenotes}
\end{threeparttable}
\end{table}

\begin{verbatim}
#> Description of Nanga_Depth.Sol
#> 
#>  Numeric 
#>        mean median  var   sd valid.n
#> 0_10  19.32  19.28 5.20 2.28       6
#> 10_20 14.49  14.98 4.35 2.09       6
#> 20_30 11.53  12.24 5.36 2.31       6
\end{verbatim}

\begin{table}[H]
\centering
\caption{\label{tab:StockscarbonesolNanga-Eboko3}Stocks de carbone du sol dans les différentes catégories d'utilisation de terres dans la FC de Nanga-Eboko.}
\centering
\begin{threeparttable}
\begin{tabular}[t]{lr}
\toprule
Strates & Carbone.total\\
\midrule
CUL & 16.38\\
FSA & 15.93\\
FGA & 15.65\\
FSJ & 15.38\\
SAV & 15.33\\
\addlinespace
MAR & 12.01\\
\bottomrule
\end{tabular}
\begin{tablenotes}
\item \textit{Note: } 
\item length: taille de l'échantillon; mean: moyenne; sd: écart type; range: amplitude; FSA: Forêts Secondaires Adultes; FGA: Forêts Galeries; FSJ: Forêts Secondaires Jeunes; MAR: Marecages; SAV: Savanes; CUL: Cultures.
\end{tablenotes}
\end{threeparttable}
\end{table}

\begin{figure}

{\centering \includegraphics[width=0.8\linewidth]{03-resultats_et_discussions_files/figure-latex/ttestsolNanga-Eboko3-1} 

}

\caption{t.test des stocks de carbone du sol dans la FC de Nanga-ebiko.}\label{fig:ttestsolNanga-Eboko3}
\end{figure}

\begin{figure}

{\centering \includegraphics[width=0.8\linewidth]{03-resultats_et_discussions_files/figure-latex/BarplotcarbonesolNanga-Eboko3-1} 

}

\caption{Répartition des stocks de carbone du sol dans les TUT de la FC de Nanga-Eboko.}\label{fig:BarplotcarbonesolNanga-Eboko3}
\end{figure}

\begin{figure}

{\centering \includegraphics[width=0.8\linewidth]{03-resultats_et_discussions_files/figure-latex/BarplotcarboneprofondeursolNanga-Eboko3-1} 

}

\caption{Répartition des stocks de carbone du sol suivant la profondeur dans les TUT de la FC de Nanga-Eboko.}\label{fig:BarplotcarboneprofondeursolNanga-Eboko3}
\end{figure}

\begin{figure}

{\centering \includegraphics[width=0.8\linewidth]{03-resultats_et_discussions_files/figure-latex/GradientcarbonesolNanga-Eboko3-1} 

}

\caption{Gradient de distribution des stocks de carbone  du sol dans la FC de Nanga-Eboko.}\label{fig:GradientcarbonesolNanga-Eboko3}
\end{figure}

\hypertarget{stocks-de-carbone-dans-les-profondeurs-de-0-uxe0-10-cm}{%
\subsubsection{Stocks de carbone dans les profondeurs de 0 à 10 cm}\label{stocks-de-carbone-dans-les-profondeurs-de-0-uxe0-10-cm}}

Les résultats montrent que les stocks de carbone dans la première couche varit de 8,15 à 35,07 tC/ha pour une moyenne de 19,17 ± 6,76 tC/ha respectivement dans les Marécages et les Forêts secondaires adultes:

\begin{verbatim}
#> Description of Nanga_StC_Sol.0_10
#> 
#>  Numeric 
#>          mean   sd median  min   max valid.n
#> MCsolHa 19.17 6.76  17.51 8.15 35.07      23
#> 
#>  Factor 
#>          
#> Depth     0_10 10_20 20_30
#>   Count     23     0     0
#>   Percent  100     0     0
#> Mode 0_10 
#>          
#> Strates     FSA   FSJ   MAR   FGA   SAV CUL
#>   Count    7.00  4.00  4.00  3.00  3.00 2.0
#>   Percent 30.43 17.39 17.39 13.04 13.04 8.7
#> Mode FSA
\end{verbatim}

La moyenne de stocks la plus élevée se situe dans les Cultures (22,48 ± 13,18 tC/ha) tandis que la plus faible est dans les Forêts secondaires jeunes (16,6 ± 2,37 tC/ha) Tableau:\ref{tab:Stockscarbonesolde0-10cmNanga-Eboko3}. Les tests statistiques ne montrent aucune différence significative dans la distribution des stocks de carbone pour cette couche (Figure:\ref{fig:ttestsol0-10cmNanga-Eboko3}).

\begin{table}[H]
\centering
\caption{\label{tab:Stockscarbonesolde0-10cmNanga-Eboko3}Stocks de carbone du sol de 0 à 10 cm dans les différentes catégories d'utilisation de terres dans la FC de Nanga-Eboko.}
\centering
\begin{threeparttable}
\begin{tabular}[t]{lr}
\toprule
Strates & total.Carbone\\
\midrule
CUL & 22.49\\
FGA & 20.97\\
FSA & 20.20\\
MAR & 18.36\\
SAV & 17.29\\
\addlinespace
FSJ & 16.60\\
\bottomrule
\end{tabular}
\begin{tablenotes}
\item \textit{Note: } 
\item length: taille de l'échantillon; mean: moyenne; sd: écart type; range: amplitude; FSA: Forêts Secondaires Adultes; FGA: Forêts Galeries; FSJ: Forêts Secondaires Jeunes; MAR: Marecages; SAV: Savanes; CUL: Cultures.
\end{tablenotes}
\end{threeparttable}
\end{table}

\begin{table}[H]
\centering
\caption{\label{tab:Stocksmoyenssol0-10cmFCNanga-Eboko}Calcul des stocks moyens de carbone du sol de 0 à 10 cm dans les TUT de la FC de Nanga-Eboko.}
\centering
\begin{threeparttable}
\begin{tabular}[t]{lrrrr}
\toprule
Strates & length & mean & sd & range\\
\midrule
FSA & 7 & 20.20 & 7.66 & 22.61\\
FSJ & 4 & 16.60 & 2.37 & 5.50\\
MAR & 4 & 18.36 & 9.45 & 20.77\\
FGA & 3 & 20.97 & 6.00 & 10.76\\
SAV & 3 & 17.29 & 4.65 & 8.46\\
\addlinespace
CUL & 2 & 22.49 & 13.18 & 18.64\\
\bottomrule
\end{tabular}
\begin{tablenotes}
\item \textit{Note: } 
\item length: taille de l'échantillon; mean: moyenne; sd: écart type; range: amplitude; FSA: Forêts Secondaires Adultes; FGA: Forêts Galeries; FSJ: Forêts Secondaires Jeunes; MAR: Marécages; SAV: Savanes; CUL: Cultures.
\end{tablenotes}
\end{threeparttable}
\end{table}

\begin{figure}

{\centering \includegraphics[width=0.8\linewidth]{03-resultats_et_discussions_files/figure-latex/ttestsol0-10cmNanga-Eboko3-1} 

}

\caption{t.test des stocks de carbone du sol de 0 à 10 cm dans la FC de Nanga-Eboko.}\label{fig:ttestsol0-10cmNanga-Eboko3}
\end{figure}

\hypertarget{stocks-de-carbone-dans-les-profondeurs-de-10-uxe0-20-cm}{%
\subsubsection{Stocks de carbone dans les profondeurs de 10 à 20 cm}\label{stocks-de-carbone-dans-les-profondeurs-de-10-uxe0-20-cm}}

Dans la couche intermédiaire les stocks de carbone varient de 4,37 tC/ha dans les Marécages à 21,85 tC/ha dans les Cultures pour une moyenne générale de 14,35 ± 4,41 tC/ha:

\begin{verbatim}
#> Description of Nanga_StC_Sol.10_20
#> 
#>  Numeric 
#>          mean   sd median  min   max valid.n
#> MCsolHa 14.35 4.41  14.75 4.37 21.85      23
#> 
#>  Factor 
#>          
#> Depth     10_20 0_10 20_30
#>   Count      23    0     0
#>   Percent   100    0     0
#> Mode 10_20 
#>          
#> Strates     FSA   FSJ   MAR   FGA   SAV CUL
#>   Count    7.00  4.00  4.00  3.00  3.00 2.0
#>   Percent 30.43 17.39 17.39 13.04 13.04 8.7
#> Mode FSA
\end{verbatim}

Les Forêts secondaires jeunes détiennent la moyenne de carbone du sol la plus élevée dans cette couche (16,47 ± 3,37 tC/ha) tandis que les Marécages en renferment beaucoup moins (10,52 ± 5,01 tC/ha) Tableau:\ref{tab:Stockscarbonesolde10-20cmNanga-Eboko3}. Du point de vue statistique il n'existe pas de différence significative à cette profondeur dans la distribution des stocks de carbone entre les différentes strates qui constituent la FC de Nanga-Eboko (Figure:\ref{fig:ttestsol10-20cmNanga-Eboko3}).

\begin{table}[H]
\centering
\caption{\label{tab:Stockscarbonesolde10-20cmNanga-Eboko3}Stocks de carbone du sol de 10 à 20 cm dans les différentes catégories d'utilisation de terres dans la FC de Nanga-Eboko.}
\centering
\begin{threeparttable}
\begin{tabular}[t]{lr}
\toprule
Strates & total.Carbone\\
\midrule
FSJ & 16.47\\
SAV & 15.61\\
CUL & 15.39\\
FGA & 14.57\\
FSA & 14.40\\
\addlinespace
MAR & 10.52\\
\bottomrule
\end{tabular}
\begin{tablenotes}
\item \textit{Note: } 
\item length: taille de l'échantillon; mean: moyenne; sd: écart type; range: amplitude; FSA: Forêts Secondaires Adultes; FGA: Forêts Galeries; FSJ: Forêts Secondaires Jeunes; MAR: Marecages; SAV: Savanes; CUL: Cultures.
\end{tablenotes}
\end{threeparttable}
\end{table}

\begin{table}[H]
\centering
\caption{\label{tab:Stocksmoyenssol10-20cmFCNanga-Eboko}Calcul des stocks moyens de carbone du sol de 10 à 20 cm dans les TUT de la FC de Nanga-Eboko.}
\centering
\begin{threeparttable}
\begin{tabular}[t]{lrrrr}
\toprule
Strates & length & mean & sd & range\\
\midrule
FSA & 7 & 14.40 & 3.42 & 9.92\\
FSJ & 4 & 16.47 & 3.38 & 7.58\\
MAR & 4 & 10.52 & 5.01 & 12.19\\
FGA & 3 & 14.57 & 2.27 & 4.43\\
SAV & 3 & 15.61 & 5.95 & 11.89\\
\addlinespace
CUL & 2 & 15.39 & 9.14 & 12.92\\
\bottomrule
\end{tabular}
\begin{tablenotes}
\item \textit{Note: } 
\item length: taille de l'échantillon; mean: moyenne; sd: écart type; range: amplitude; FSA: Forêts Secondaires Adultes; FGA: Forêts Galeries; FSJ: Forêts Secondaires Jeunes; MAR: Marécages; SAV: Savanes; CUL: Cultures.
\end{tablenotes}
\end{threeparttable}
\end{table}

\begin{figure}

{\centering \includegraphics[width=0.8\linewidth]{03-resultats_et_discussions_files/figure-latex/ttestsol10-20cmNanga-Eboko3-1} 

}

\caption{t.test des stocks de carbone du sol de 10 à 20 cm dans la FC de Nanga-Eboko.}\label{fig:ttestsol10-20cmNanga-Eboko3}
\end{figure}

\hypertarget{stocks-de-carbone-dans-les-profondeurs-de-20-uxe0-30-cm}{%
\subsubsection{Stocks de carbone dans les profondeurs de 20 à 30 cm}\label{stocks-de-carbone-dans-les-profondeurs-de-20-uxe0-30-cm}}

Dans la dernière couche (20 à 30 cm) les stocks de carbone varient de 8,15 tC/ha dans les Marécages à 35,07 tC/ha dans les forêts secondaires adultes pour une moyenne de 19,17 ± 6,76 tC/ha:

\begin{verbatim}
#> Description of Nanga_StC_Sol.0_10
#> 
#>  Numeric 
#>          mean   sd median  min   max valid.n
#> MCsolHa 19.17 6.76  17.51 8.15 35.07      23
#> 
#>  Factor 
#>          
#> Depth     0_10 10_20 20_30
#>   Count     23     0     0
#>   Percent  100     0     0
#> Mode 0_10 
#>          
#> Strates     FSA   FSJ   MAR   FGA   SAV CUL
#>   Count    7.00  4.00  4.00  3.00  3.00 2.0
#>   Percent 30.43 17.39 17.39 13.04 13.04 8.7
#> Mode FSA
\end{verbatim}

La moyenne la plus élévée parmi les strates se situe dans les forêts secondaires adultes (13,19 tC/ha) tandis que la plus faible est de 7,16 tC/ha et se localise dans les Marécages (Tableau:\ref{tab:Stockscarbonesolde20-30cmNanga-Eboko}). Contrairement aux deux couches précédentes, les tests statistiques montrent des différences significatives dans la distribution des stocks de carbone entre les Marécages-Forêts secondaires adultes d'une part et Marécages-Forêts secondaires jeunes d'autre part (Figure:\ref{fig:ttestsol20-30cmNanga-Eboko}). Cette difference pourrait traduire l'influence de la nature du sol dans le processus de stockage du carbone organique (Figure:\ref{fig:GradientcarbonesolNanga-Eboko3}). En effet les sols des zones marécageuses étaient de type sableux tandis ceux des strates forestières étaient de nature argilo-sableux. Les sols forestiers connaissent un apport constant de matières oraganiques provenant soit du bois mort soit de la chute des feuilles des arbres. En se décomposant lentement, ces débris enrichissent les couches superficielles et profondes du sol par un processus d'enfouissement. Le fait que les sols sableux des zones marécageuses de la FC de Nanga-Eboko sont moins riches en carbone peut laisser penser qu'il s'agirait d'un dépôt récent de sédiment qui recouvrirait des couches plus anciennes et probablement riches en matières organiques comme cela a pu être observé sur les terrains marécageux de la FC de Ndikiniméki.

\begin{table}[H]
\centering
\caption{\label{tab:Stockscarbonesolde20-30cmNanga-Eboko}Stocks de carbone du sol de 20 à 30 cm dans les différentes catégories d'utilisation de terres dans la FC de Nanga-Eboko.}
\centering
\begin{threeparttable}
\begin{tabular}[t]{lr}
\toprule
Strates & total.Carbone\\
\midrule
FSA & 13.19\\
SAV & 13.10\\
FSJ & 13.07\\
FGA & 11.41\\
CUL & 11.27\\
\addlinespace
MAR & 7.16\\
\bottomrule
\end{tabular}
\begin{tablenotes}
\item \textit{Note: } 
\item length: taille de l'échantillon; mean: moyenne; sd: écart type; range: amplitude; FSA: Forêts Secondaires Adultes; FGA: Forêts Galeries; FSJ: Forêts Secondaires Jeunes; MAR: Marecages; SAV: Savanes; CUL: Cultures.
\end{tablenotes}
\end{threeparttable}
\end{table}

\begin{table}[H]
\centering
\caption{\label{tab:Stocksmoyenssol20-30cmFCNanga-Eboko}Calcul des stocks moyens de carbone du sol de 20 à 30 cm dans les TUT de la FC de Nanga-Eboko.}
\centering
\begin{threeparttable}
\begin{tabular}[t]{lrrrr}
\toprule
Strates & length & mean & sd & range\\
\midrule
FSA & 7 & 13.19 & 3.63 & 10.71\\
FSJ & 4 & 13.07 & 1.70 & 3.20\\
MAR & 4 & 7.16 & 2.31 & 5.51\\
FGA & 3 & 11.41 & 2.41 & 4.63\\
SAV & 3 & 13.10 & 4.58 & 9.17\\
\addlinespace
CUL & 2 & 11.27 & 4.80 & 6.79\\
\bottomrule
\end{tabular}
\begin{tablenotes}
\item \textit{Note: } 
\item length: taille de l'échantillon; mean: moyenne; sd: écart type; range: amplitude; FSA: Forêts Secondaires Adultes; FGA: Forêts Galeries; FSJ: Forêts Secondaires Jeunes; MAR: Marécages; SAV: Savanes; CUL: Cultures.
\end{tablenotes}
\end{threeparttable}
\end{table}

\begin{figure}

{\centering \includegraphics[width=0.8\linewidth]{03-resultats_et_discussions_files/figure-latex/ttestsol20-30cmNanga-Eboko-1} 

}

\caption{t.test des stocks de carbone du sol de 20 à 30 cm dans la FC de Nanga-Eboko.}\label{fig:ttestsol20-30cmNanga-Eboko}
\end{figure}

\hypertarget{footnotes-and-citations}{%
\chapter{Footnotes and citations}\label{footnotes-and-citations}}

\hypertarget{footnotes}{%
\section{Footnotes}\label{footnotes}}

Footnotes are put inside the square brackets after a caret \texttt{\^{}{[}{]}}. Like this one \footnote{This is a footnote.}.

\hypertarget{citations}{%
\section{Citations}\label{citations}}

Reference items in your bibliography file(s) using \texttt{@key}.

For example, we are using the \textbf{bookdown} package \citep{R-bookdown} (check out the last code chunk in index.Rmd to see how this citation key was added) in this sample book, which was built on top of R Markdown and \textbf{knitr} \citep{xie2015} (this citation was added manually in an external file book.bib).
Note that the \texttt{.bib} files need to be listed in the index.Rmd with the YAML \texttt{bibliography} key.

The \texttt{bs4\_book} theme makes footnotes appear inline when you click on them. In this example book, we added \texttt{csl:\ chicago-fullnote-bibliography.csl} to the \texttt{index.Rmd} YAML, and include the \texttt{.csl} file. To download a new style, we recommend: \url{https://www.zotero.org/styles/}

The RStudio Visual Markdown Editor can also make it easier to insert citations: \url{https://rstudio.github.io/visual-markdown-editing/\#/citations}

\hypertarget{blocks}{%
\chapter{Blocks}\label{blocks}}

\hypertarget{equations}{%
\section{Equations}\label{equations}}

Here is an equation.

\begin{equation} 
  f\left(k\right) = \binom{n}{k} p^k\left(1-p\right)^{n-k}
  \label{eq:binom}
\end{equation}

You may refer to using \texttt{\textbackslash{}@ref(eq:binom)}, like see Equation \eqref{eq:binom}.

\hypertarget{theorems-and-proofs}{%
\section{Theorems and proofs}\label{theorems-and-proofs}}

Labeled theorems can be referenced in text using \texttt{\textbackslash{}@ref(thm:tri)}, for example, check out this smart theorem \ref{thm:tri}.

\begin{theorem}
\protect\hypertarget{thm:tri}{}\label{thm:tri}For a right triangle, if \(c\) denotes the \emph{length} of the hypotenuse
and \(a\) and \(b\) denote the lengths of the \textbf{other} two sides, we have
\[a^2 + b^2 = c^2\]
\end{theorem}

Read more here \url{https://bookdown.org/yihui/bookdown/markdown-extensions-by-bookdown.html}.

\hypertarget{callout-blocks}{%
\section{Callout blocks}\label{callout-blocks}}

The \texttt{bs4\_book} theme also includes special callout blocks, like this \texttt{.rmdnote}.

You can use \textbf{markdown} inside a block.

\begin{verbatim}
#>   day time  temp activ
#> 1 346  840 36.33     0
#> 2 346  850 36.34     0
#> 3 346  900 36.35     0
#> 4 346  910 36.42     0
#> 5 346  920 36.55     0
\end{verbatim}

It is up to the user to define the appearance of these blocks for LaTeX output.

You may also use: \texttt{.rmdcaution}, \texttt{.rmdimportant}, \texttt{.rmdtip}, or \texttt{.rmdwarning} as the block name.

The R Markdown Cookbook provides more help on how to use custom blocks to design your own callouts: \url{https://bookdown.org/yihui/rmarkdown-cookbook/custom-blocks.html}

\hypertarget{sharing-your-book}{%
\chapter{Sharing your book}\label{sharing-your-book}}

\hypertarget{publishing}{%
\section{Publishing}\label{publishing}}

HTML books can be published online, see: \url{https://bookdown.org/yihui/bookdown/publishing.html}

\hypertarget{pages}{%
\section{404 pages}\label{pages}}

By default, users will be directed to a 404 page if they try to access a webpage that cannot be found. If you'd like to customize your 404 page instead of using the default, you may add either a \texttt{\_404.Rmd} or \texttt{\_404.md} file to your project root and use code and/or Markdown syntax.

\hypertarget{metadata-for-sharing}{%
\section{Metadata for sharing}\label{metadata-for-sharing}}

Bookdown HTML books will provide HTML metadata for social sharing on platforms like Twitter, Facebook, and LinkedIn, using information you provide in the \texttt{index.Rmd} YAML. To setup, set the \texttt{url} for your book and the path to your \texttt{cover-image} file. Your book's \texttt{title} and \texttt{description} are also used.

This \texttt{bs4\_book} provides enhanced metadata for social sharing, so that each chapter shared will have a unique description, auto-generated based on the content.

Specify your book's source repository on GitHub as the \texttt{repo} in the \texttt{\_output.yml} file, which allows users to view each chapter's source file or suggest an edit. Read more about the features of this output format here:

\url{https://pkgs.rstudio.com/bookdown/reference/bs4_book.html}

Or use:

  \bibliography{referencess.bib}

\end{document}
